\documentclass[11pt]{ctexart}
\ifdefined\pdfinfoomitdate\pdfinfoomitdate=1\fi
\ifdefined\pdftrailerid\pdftrailerid{}\fi
\ifdefined\pdfsuppressptexinfo\pdfsuppressptexinfo=15\fi
\usepackage[margin=1in]{geometry}
\usepackage{amsmath,amssymb,amsthm,mathtools}
\usepackage{microtype}
\usepackage{enumitem,booktabs,array}
\usepackage{tikz}
\usepackage{placeins,flafter}
\usepackage[titles]{tocloft}
\usepackage[hyphens]{url}
\usetikzlibrary{arrows.meta,positioning,calc}
\usepackage{xcolor}
\definecolor{linkblue}{RGB}{25,63,105}
\usepackage[colorlinks=true,allcolors=linkblue,pdfencoding=auto,psdextra]{hyperref}
\usepackage[nameinlink,capitalise,noabbrev]{cleveref}
\hypersetup{
  pdftitle={唯一游戏定理（中文对照版）},
  pdfsubject={从 3SAT 到唯一游戏的确定性归约},
  pdfauthor={OpenAI（原文）；openai-math-zh 项目（翻译）},
  pdfkeywords={Unique Games, approximation hardness, shortcode, projection games}
}

% —— 中文环境名 ——
\newtheorem{theorem}{定理}[section]
\newtheorem{lemma}[theorem]{引理}
\newtheorem{proposition}[theorem]{命题}
\newtheorem{corollary}[theorem]{推论}
\theoremstyle{definition}
\newtheorem{definition}[theorem]{定义}
\theoremstyle{remark}
\newtheorem{remark}[theorem]{注记}
\crefname{theorem}{定理}{定理}
\crefname{lemma}{引理}{引理}
\crefname{proposition}{命题}{命题}
\crefname{corollary}{推论}{推论}
\crefname{definition}{定义}{定义}
\crefname{remark}{注记}{注记}
\crefname{equation}{式}{式}
\crefname{section}{节}{节}
\crefname{appendix}{附录}{附录}
\crefname{figure}{图}{图}
\crefname{table}{表}{表}
\numberwithin{equation}{section}
\newcommand{\F}{\mathbb F}
\newcommand{\E}{\mathbb E}
\newcommand{\eps}{\varepsilon}
\DeclareMathOperator{\OPT}{OPT}
\DeclareMathOperator{\val}{val}
\DeclareMathOperator{\Hom}{Hom}
\DeclareMathOperator{\im}{im}
\DeclareMathOperator{\rank}{rank}
\DeclareMathOperator{\Span}{span}
\DeclareMathOperator{\supp}{supp}
\DeclareMathOperator{\TV}{TV}
\newcommand{\abs}[1]{\lvert #1\rvert}
\newcommand{\norm}[1]{\lVert #1\rVert}
\newcommand{\ip}[2]{\langle #1,#2\rangle}
\setlist[itemize]{topsep=4pt,itemsep=2pt,parsep=0pt}
\setlist[enumerate]{topsep=4pt,itemsep=2pt,parsep=0pt}
\setlength{\emergencystretch}{2em}
\allowdisplaybreaks[1]

\title{唯一游戏定理\\{\large The Unique Games Theorem（中文对照版）}}
\author{OpenAI（原文）\\中文翻译：openai-math-zh 项目}
\date{2026 年 9 月 23 日（原文日期）}

\begin{document}
\maketitle
\begin{abstract}
我们证明唯一游戏猜想。对每个固定的
\(\varepsilon,\delta\in(0,1/2)\)，我们给出从 \(\mathrm{3SAT}\) 到固定
有限字母表上唯一游戏的确定性多项式时间归约，完备性至少
\(1-\varepsilon\)，可靠性至多 \(\delta\)。

\smallskip
\emph{（本文件是 openai/math 手稿的中文对照翻译；原文结果未经同行评审。
数学公式、定理编号与结构与英文原文保持一致，术语以仓库 glossary.md
为准；译文与原文冲突时以英文原文为准。）}
\end{abstract}

\begingroup
\small
\setlength{\baselineskip}{10.5pt}
\setlength{\cftbeforesecskip}{3pt}
\tableofcontents
\endgroup
\clearpage
\section{引言}\label{sec:introduction}

一个置换约束规定：对另一端点的每个标记，本端点恰有一个对应标记。
唯一游戏猜想断言：即使几乎所有约束都能被满足，这类约束系统也可能
难以近似。我们正面解决该猜想，给出一个显式的无权归约，从
\(\mathrm{3SAT}\) 出发。

\subsection{定理}

\emph{唯一游戏实例}由有限顶点集、有限字母表 \(K\)、以及非空的定向
约束列表 \(e=(u_e,v_e,\pi_e)\) 组成，其中 \(\pi_e\) 是 \(K\) 的一个
置换。染色 \(a\) 满足 \(e\) 当
\(a(v_e)=\pi_e(a(u_e))\)。其值是被满足约束的比例，
\(\val(G)\) 是染色的最大值。当顶点与约束列表被完整写出（包括每个
置换的表）时，称实例\emph{显式且无权}。简单二部实例有两个不相交的
顶点集，且每对相对顶点上至多一条约束。

\begin{theorem}[唯一游戏定理]\label{thm:ugc}
对每个固定的 \(\varepsilon,\delta\in(0,1/2)\)，存在整数 \(s\ge1\)
与一个确定性多项式时间归约，从 \(\mathrm{3SAT}\) 到 \(K=\F_2^s\) 上
显式无权唯一游戏实例 \(G_\varphi\)，使得
\[
\begin{aligned}
 \varphi\text{ 可满足}
 &\quad\Longrightarrow\quad \val(G_\varphi)\ge1-\varepsilon,\\
 \varphi\text{ 不可满足}
 &\quad\Longrightarrow\quad \val(G_\varphi)\le\delta.
\end{aligned}
\]
图为简单二部图，且每条约束都是 \(K\) 的一个平移
\(a(v_e)=a(u_e)+c_e\)。字母表、运行时间多项式的次数及其常数只依赖
\(\varepsilon,\delta\)。
\end{theorem}

因此，该结果以其通常的边值表述正面解决唯一游戏
猜想~\cite[Conjecture~2.5]{Khot10}。两个误差被独立指定，且先于
字母表选取。二元平移约束是该猜想的一个已知的等价表述
~\cite[Corollary~13 and Section~11.1]{KKMO07}；本归约直接产生该
形式。

Theorem~\ref{thm:ugc} 为已有的针对割、覆盖、约束满足、排序、删除与
聚类问题的难度归约提供了唯一游戏假设
\cite{KKMO07,KhotRegev2008,Raghavendra08,GuruswamiEtAl2011OrderingCSP,CKKRS2006,DemaineEtAl2006Clustering}。
那些归约及由此得到的近似阈值均出自上述作者。
Section~\ref{sec:consequences} 给出精确的问题表述、严格阈值与归约
限定条件，并附各推论的证明。

\subsection{证明、编码与量词}

核心难点在于：当每个标记只容许唯一应答时，如何保持接近 1 的完备
性。证明验证与近似难度之间的早期联系由 Feige、Goldwasser、
Lov\'asz、Safra 与 Szegedy~\cite{FGLSS96}，由 Arora 与 Safra 经由
证明复合~\cite{AroraSafra98}，以及由 Arora、Lund、Motwani、Sudan
与 Szegedy 经由常数查询 PCP 定理~\cite{ALMSS98} 发展而来。
Bellare、Goldreich 与 Sudan 为验证者复合引入了长码（long code）与
折叠~\cite[Section~3.3]{BGS98}。答案 \(a\) 的长码是表
\(f\mapsto f(a)\)，以答案集上的布尔函数为指标。在 H\aa stad 的奇偶
难度证明中，傅里叶分析从使验证者足够频繁接受的任意证明表恢复出
外层博弈策略，而无需先恢复邻近码字。他的近可满足奇偶间隙是我们的
起始难度输入~\cite{Hastad01}。

Khot 于 2002 年提出唯一游戏猜想，并给出基于半定舍入的算法
~\cite{Khot02}。若干算法保证可以解释其量词顺序为何重要。对固定的
目标误差，Trevisan 把半定舍入与图分解结合，在输入误差按实例规模的
倒数对数速率缩小时得到多项式时间保证~\cite{Trevisan08}。
Charikar、Makarychev 与 Makarychev 用随机化半定舍入，从值至少
\(1-\epsilon\) 的 \(q\) 标记实例得到期望值
\(1-O(\sqrt{\epsilon\log q})\)~\cite{CMM06}；该保证随字母表增大而
失准，而猜想允许在误差选定后使用大的固定字母表。Kolla 的谱枚举
给出以标记扩展图（顶点为「变量–标记」对）适当谱子空间维数计的运行
时间界~\cite{Kolla11}。Arora、Barak 与 Steurer 用谱分解与枚举，对
足够小的固定完备性误差得到近似满足解，其保证与固定字母表规模无关
~\cite{ABS10}；对固定的正误差，其运行时间界是次指数的而非多项式
的。这些是已知保证的各自局限。

整性间隙关心的是松弛的强度。Khot 与 Vishnoi 构造了半定值任意接近
1、整数值任意接近 0 的唯一游戏实例，即使加上三角不等式也如此
~\cite{KV15}。Barak、Gopalan、H\aa stad、Meka、Raghavendra 与
Steurer 基于 Reed--Muller 码引入短码（short code），作为长码在字母
表归约与定量整性间隙等应用中的更短替代~\cite{BGHMRS15}。这些构造
约束的是特定松弛；难度归约需要单独论证。Bafna、Barak、Kothari、
Schramm 与 Steurer 对具有适当低度平方和超压缩性证书的图（含噪声
超立方与噪声 Reed--Muller 短码图）的近可满足平移博弈，得到多项式
时间常值保证~\cite{BBKSS21}。下文使用的秩一矩阵短码是另一种图，且
我们的论证不假设这些证书。

\subsection{Grassmann 扩张与完备性障碍}

最接近的前置难度突破涉及完备性不完美的二对一与二对二博弈。二对一
约束是每个右标记恰有两个原像的投影。二对二约束把两侧标记各分成
对，并在这些对之间建立双射匹配，允许匹配对内的四种组合。Khot、
Minzer 与 Safra 的 STOC 2017 工作提出 Grassmann 图归约，其非标准的
可靠性保证以一个结构性假设为条件~\cite{KMS17}。Dinur、Khot、
Kindler、Minzer 与 Safra 的 STOC 2018 工作提出具有普通边满足可靠性
的归约，以一致性假设为条件~\cite{DKKMS25}；其扩张工作提出了相关的
结构问题并建立了部分情形~\cite{DKKMS21Expansion}。Barak、Kothari
与 Steurer 通过二元矩阵表示把扩张与所需的一致性测试联系起来
~\cite{BKS18}。Khot、Minzer 与 Safra 的扩张定理于 2018 年完成了
这条线~\cite{KMS23}。

这些结果对几乎可满足的二对一与二对二博弈建立了难度，并作为推论给
出完备性接近二分之一、可靠性任意小的唯一游戏间隙。使唯一游戏的
完备性任意接近 1 是一个单独的问题。\footnote{一篇独立的姊妹工作
证明：在固定字母表上，完备性完美、固定小可靠性下精确二对一博弈的
难度~\cite[Theorem~1.1]{OpenAIPerfectTwoToOne2026}。它不是本证明的
输入。}

Barak、Kothari 与 Steurer 的二次矩阵短码~\cite{BKS18} 是本文使用
的表述：其顶点是二元矩阵，其步骤添加秩一矩阵。他们的双应答规则
容许旧答案或其被更新左因子平移后的答案；只要求相等只能给出二分之
一的完备性保证。他们讨论过更高次的测试，作为改进完备性的可能路线
~\cite[Section~1.6]{BKS18}。我们使用的逆命题（由 Grassmann 扩张
导出）见 Theorem~\ref{thm:shortcode}。Ellis、Kindler 与 Lifshitz
随后给出更短的扩张证明，使用线性映射空间上的超压缩性~\cite{EKL26}；
我们的外部输入仍是 Khot--Minzer--Safra 定理。我们的证明保留已有的
二次逆定理，并用一个非线性映射改变噪声。下面的比较解释这一选择。

\subsection{构造与证明}

均匀秩一测试展示了完备性问题。设 \(M\) 是二元 \(\ell\times m\)
矩阵，固定非零答案 \(z\in\F_2^m\)。求值函数 \(f_z(M)=Mz\) 满足
\[
 f_z(M+a l^\top)-f_z(M)=a(l^\top z).
\]
对独立的均匀 \(a\in\F_2^\ell\) 与 \(l\in\F_2^m\)，该求值以概率
\((1+2^{-\ell})/2\) 保持其值，接近二分之一。我们的测试改变噪声，
同时保留普通秩一测试作为可靠性分析的结构工具。

对任意指定的 \(p>0\)，Lemma~\ref{lem:latent} 固定一个秩阈值，并对
每个足够大的 \(\ell\) 构造包含 \(K=\F_2^\ell\) 的二元空间
\(\mathcal V\)、映射 \(C:\mathcal V\to K\) 与 \(\mathcal V\) 上的
噪声律 \(\mu\)，使得对 \(a\sim\mu\)，
\[
 C(x+h)=C(x)+h\quad(h\in K),\qquad
 \Pr[C(x+a)\ne C(x)]\le p
\]
其中 \(x\) 均匀且与 \(a\) 独立。同时，任何限制到 \(K\) 后秩足够大
的线性观察族以至少 \(1/8\) 的概率检测出 \(a\)。秩阈值只依赖
\(p\)，且在 \(\ell\) 之前选定。该构造使用有限域上的二次块：其非
线性输出在递归下越来越稳定，而一个秩势控制线性观察的损失。
Section~\ref{sec:latent} 证明这一分离，并把输出平移实现为子空间
\(K\)。

对同一个非零 \(z\)，设 \(P\in\Hom(\F_2^m,\mathcal V)\) 均匀。非线性
求值 \(P\mapsto C(Pz)\) 在更新 \(P\mapsto P+a l^\top\) 下比较
\[
 C(Pz)\quad\text{与}\quad C\bigl(Pz+a(l^\top z)\bigr)
\]
两个值。此处 \(Pz\) 在 \(\mathcal V\) 中均匀，与潜在噪声 \(a\)
独立，而 \(l^\top z\) 是独立的公平比特。因此其相等测试以至多
\(p/2\) 的概率失败。

归约从三个不同变量上的二元奇偶方程出发。每条方程有四个满足的
三元组。方程元组 \(U=(e_1,\ldots,e_k)\) 上的表是从其解平面的乘积
到 \(\mathcal V\) 的仿射映射，由一个 \(m=1+2k\) 列的矩阵表示：一个
截距列加每条方程两个斜率列。当某个方程块只通过一个实际比特影响
一张表时，其约化键（reduced key）记录该变量与整个诱导仿射映射。
投影表与其拉回具有相同的键。把键的常数 \(K\)-平移归组为一个顶点，
就折叠了答案：代表元处的标记决定每个平移后的答案，而共享的键在
两个域上给出相同的恢复答案。所得的唯一游戏测试利用潜在噪声，检验
一张表与其秩一扰动上恢复答案的相等性。一个近满足的全局奇偶赋值在
其满足的每个元组上给出上述理想求值，从而提供接近 1 的完备性。
键共享改编了 Khot 与 Safra 的虚拟表~\cite[Section~1.1]{KS13}，而
平移轨道改编了 Bellare、Goldreich 与 Sudan 的折叠原理
~\cite[Section~3.3]{BGS98}。

对可靠性，设某染色以至少 \(0.99\) 的概率通过该测试。矩阵频率是
元组 \(g=(g_i)_{i=0}^{m-1}\in(\mathcal V^*)^m\)（每列一个泛函），
其特征为 \((-1)^{\sum_i g_i(P\mathbf e_i)}\)，其中 \(\mathbf e_i\)
是 \(\F_2^m\) 的标准基向量。把秩一更新的相位增量对均匀选择子
\(l\) 取平均：
\[
\begin{aligned}
 \E_{a\sim\mu,l}(-1)^{\sum_i l_i g_i(a)}
   &=\Pr_{a\sim\mu}[g_i(a)=0\text{ 对每个 }i],\\
 \E_{a\text{ 在 }K\text{ 中均匀},l}(-1)^{\sum_i l_i g_i(a)}
   &=2^{-\dim\Span\{g_i|_K\}}.
\end{aligned}
\]
第一个乘子非负，且由潜在检测，当限制秩大时至多 \(7/8\)。因此高
接受率迫使固定数量的低秩傅里叶质量；第二个乘子在 \(K\) 中的普通
均匀秩一扰动下把这一质量转化为正的接受率。使用与潜在数据一起固定
的线性补 \(\mathcal V=K\oplus W\)，把表写成 \(P=(M,T)\)，其中
\(M\) 与 \(T\) 是其取值于 \(K\) 与 \(W\) 的仿射分量，用含常数列的
矩阵表示。普通扰动改变 \(M\) 而固定 \(T\)。于是，对正比例的对
\((U,T)\)，作为 \(M\) 的函数的恢复答案经常在普通测试下保持取值。
对其余的可靠性论证，取 \(k\) 为足够大的整数立方，潜在与短码参数
固定。在每个这样的对上，逆短码定理给出与某个目标
\(Mz+u\)（\(z\in\F_2^m\)、\(u\in K\)）的一致性，该一致性在由至多
固定数目 \(r_{\rm s}\) 个行方程与至多 \(r_{\rm s}\) 个列方程定义的
切片上成立。

对独立的均匀线性映射 \(A:K\to\F_2^{r_{\rm s}}\)，关于 \((M,T)\)
的可见建议是 \((A,T,AM)\)。连同 \(U\)，它确定了 \(T\) 以及 \(U\)
的答案空间上 \(K\) 值仿射映射的行纤维 \(\{M':AM'=AM\}\)（含其常数
项）。一个擦除论证证明：在 \(U\) 上的无限制表采样下，固定正比例的
此类数据容许额外的列方程，以及一个在非空切片上见证一致性的目标
\(Mz+u\)。对每个这样的数据 \((U,A,T,AM)\)，按只依赖该可见数据与
固定染色的固定规则选取一个证书。折叠随后使该证书确定 \(U\) 上的
一个短的、非空的合法答案集。

建议实验把方程元组 \(U\) 与问题 \(O\) 比较：\(O\) 以概率 \(\beta\)
独立地把每条方程替换为一个均匀选取的变量。设 \(\pi\) 为实际答案的
所得投影。裁判在 \(O\) 上采样均匀表 \((M_O,T_O)\) 及具有上述律的
独立映射 \(A\)。在实际答案空间上，它把可见数据
\((A,T_O\circ\pi,AM_O\circ\pi)\) 连同 \(U\)、把 \((A,T_O,AM_O)\)
连同 \(O\) 发给两个玩家。取下文的稀疏 \(\beta\)，对投影选择的
取平均后，完整的拉回表律接近抽取所用的无限制表律。在正比例的联合
样本上，所选的 \(U\)-证书转移到 \(O\) 上的非空切片并保持一致性。

两个策略由各自的可见数据定义。当其建议容许证书时，\(U\)-策略从
其短答案集中均匀采样；在其他输入上返回任意合法答案。\(O\)-策略
在自己的行纤维 \(\{M':AM'=AM_O\}\) 上使用傅里叶解码器。在分析中，
转移来的证书提供傅里叶质量，证明该解码器以正概率到达投影后的
\(U\)-答案集；解码器本身不接收该证书。独立的 \(U\)-样本随后给出
完整答案一致性，概率至少为固定的 \(\gamma>0\)，与 \(k\) 无关。
稀疏问题比较与傅里叶解码改编了 Khot 与 Safra 的方法~\cite{KS13}；
相关建议与消失坐标的使用遵循 Khot、Minzer 与 Safra
~\cite[Sections~3.3--3.4]{KMS17}，以及 Minzer 与 Zheng 对该框架的
后续使用~\cite[Section~3]{MZ26}。

在 NO 奇偶实例上，每对建议策略的一致性趋于零。取 \(\beta=k^{-2/3}\)
使 \(k\beta^2\to0\)（表律比较所需），同时 \(k\beta\to\infty\) 留下
大量单例坐标。当每条被观察的行在其保留比特上的系数为零时，称该
单例是干净的。在分析中，以 \(A\)、干净位置、\(O\) 上的原建议截距、
以及其他位置上采样的问题与建议斜率为条件。干净问题对于是保持独立
的方程–变量律。对这些固定数据的每个取值，由每个原局部输入自身的
干净问题重构它，即定义了重复外层博弈的策略。并行重复界定它们在
干净坐标上的成功率；对那些坐标的个数取平均，把原建议一致性界定在
\(\exp(-c k^{1/3})\) 以内。对大 \(k\)，这与解码后的下界矛盾。

并行重复我们使用 Dinur 与 Steurer 的显式投影博弈界
~\cite[Corollary~1.2]{DS14}。Raz 对一般双证明者博弈建立了指数衰减
~\cite{Raz98}，Holenstein 简化了证明~\cite{Holenstein2009}；那些界
依赖答案字母表。Rao 对投影博弈去除了这一依赖~\cite[Theorem~4]{Rao11}。
Dinur--Steurer 界先用于上述干净坐标，最后用于把矩阵博弈权重取整并
把每条边细分为四条边之后。其字母表无关性允许最终重复指数在潜在
字母表之前固定，从而完成显式无权归约。

Figure~\ref{fig:overview} 展示该归约以及确立可靠性的两个不相容
界。
\begin{figure}[htbp]
\centering
\begin{tikzpicture}[
  x=1cm,y=1cm,
  every node/.style={font=\small},
  box/.style={draw=linkblue!80!black,rounded corners=2pt,
    text width=5.9cm,align=center,inner sep=8pt,minimum height=11mm},
  branch/.style={box,text width=2.70cm,inner sep=6pt,minimum height=15mm},
  heading/.style={font=\small\bfseries,align=center},
  reduction/.style={-{Latex[length=2.1mm]},thick,draw=linkblue!80!black},
  implication/.style={-{Latex[length=2.1mm]},thick,dashed,draw=black!70}
]
\node[heading] at (0,0) {归约};
\node[heading] at (7.4,0) {NO 输入上的可靠性};

\node[box] (parity) at (0,-1.0)
  {$\mathrm{3SAT}$\\$\longrightarrow$ 带权奇偶间隙实例};
\node[box] (matrix) at (0,-2.8)
  {潜在矩阵测试\\
   稳定非线性映射：Lemma~\ref{lem:latent}};
\node[box] (ugc) at (0,-6.2)
  {无权简单二部\\唯一游戏\\
   Theorem~\ref{thm:ugc}};
\draw[reduction] (parity) -- (matrix);
\draw[reduction] (matrix) -- node[midway,fill=white,inner sep=4pt,
    text width=5.0cm,align=center]
  {取整与细分；\\并行重复} (ugc);

\node[branch] (assumption) at (5.65,-1.0)
  {假设测试接受率\\至少 $0.99$};
\node[branch] (no) at (9.15,-1.0)
  {NO 奇偶\\实例};
\node[branch] (decoding) at (5.65,-3.3)
  {解码后的建议策略\\一致概率至少 $\gamma>0$\\
   Proposition~\ref{prop:matrix-decoding}};
\node[branch] (clean) at (9.15,-3.3)
  {每一对策略：\\一致概率至多\\$\exp(-c k^{1/3})$\\
   Lemma~\ref{lem:clean}};
\node[box] (contradiction) at (7.4,-6.2)
  {对足够大的 $k$ 得到矛盾：\\
   $\exp(-c k^{1/3})<\gamma$};
\draw[implication] (assumption) -- (decoding);
\draw[implication] (no) -- (clean);
\draw[implication] (decoding) -- (contradiction);
\draw[implication] (clean) -- (contradiction);
\end{tikzpicture}
\caption{实线箭头描述归约；虚线箭头描述其可靠性论证。解码下界来自
高测试接受率，而干净坐标上界来自 NO 奇偶间隙，并适用于每个建议
策略。常数 $c,\gamma>0$ 在元组长度 $k$ 增长之前固定。}
\label{fig:overview}
\end{figure}


Section~\ref{sec:preliminaries} 固定博弈与概率约定。
Section~\ref{sec:latent} 构造 gadget，Section~\ref{sec:game} 定义
置换测试。Sections~\ref{sec:matrix} 与~\ref{sec:clean} 给出建议
实验的两个不相容界。Section~\ref{sec:finish} 完成参数选取与显式
唯一游戏归约。Section~\ref{sec:consequences} 从完成的定理与所引
归约导出近似推论。Appendix~\ref{app:finite-consequences} 给出那里
使用的简单强左表述与有限排序归约。

\FloatBarrier
\section{约定}\label{sec:preliminaries}

\paragraph{带权博弈。}
中间实例可以有自环、重复的约束出现、以及和为 1 的正有理权重。其值
是某个染色所满足的最大总权重。即使约束重合，各出现仍保留独立的
ID。所有有理数都用二进制分子与分母编码。

二部\emph{投影博弈}在两侧每个顶点上有一个问题、有限的答案集，以及
每条约束上从左答案到右答案的映射。策略为每个问题选一个答案。依赖
问题的答案集可以与同规模的固定字母表等同；对奇偶方程的四个满足
三元组我们就这么用。每个证明者只用自己的问题。与问题独立的私有
或共享随机性给出确定性策略的混合，不会增大值。

并行幂 \(G^{\otimes t}\) 采样 \(t\) 个独立的约束出现，发送有序的
问题元组，当每个坐标都接受时接受。答案可以依赖证明者的整个元组。
重复的问题 ID 保留独立的坐标位置与独立的局部答案副本。这就是重复
定理所用的普通经典乘积博弈。

\paragraph{二元线性代数与傅里叶分析。}
除特别指明更大域外，所有向量空间与映射都是 \(\F_2\) 上的有限维。
记 \(\Hom(E,F)\) 为线性映射空间，\(E^*=\Hom(E,\F_2)\) 为对偶。
子空间 \(S\le E^*\) 在子空间 \(K\le E\) 上的限制为其像
\(S|_K\le K^*\)；其限制秩为 \(\dim(S|_K)\)。

对二元空间 \(E\) 上的实函数 \(f\)，使用归一化傅里叶系数
\[
 \widehat f(g)=\E_{x\in E}f(x)(-1)^{g(x)},\qquad
 f(x)=\sum_{g\in E^*}\widehat f(g)(-1)^{g(x)},\qquad
 \sum_g|\widehat f(g)|^2=\E_x|f(x)|^2.
\]
这些恒等式源于：非零线性泛函取零与取一的次数相同，故其特征均值
为零。因此这些特征构成正交规范基。在矩阵空间上，这是与其对偶矩阵
空间的逐项配对；定义傅里叶解码器时会指明张量记号。

\paragraph{概率与有限构造。}
从向量空间均匀采样包括零。除非声明依赖，样本相互独立。对有限概率
律，
\[
 \TV(P,Q)=\frac12\sum_x|P(x)-Q(x)|.
\]
它控制任何 \([0,1]\) 值函数的期望之差。我们只在正概率事件上作
条件化。

随机实验描述有限的带权约束列表。归约确定性地枚举这些列表。每个
辅助维数与元组长度都在输入公式变化之前固定；最终证明同时检查该
参数顺序与二进制编码代价。傅里叶见证与解码器只用于界定博弈值，
不必高效。

\section{潜在字母表 gadget}\label{sec:latent}

本节的构造把一个非线性可观测与同加性噪声下的线性可观测族分离。
本节的所有向量空间与对偶都在 $\F_2$ 上，除非指明别的域。对偶的
子空间的「限制」指其在限制映射下的像。

\begin{lemma}[潜在字母表 gadget]\label{lem:latent}
对每个 $p_*>0$ 存在整数 $r_*$，使得对每个足够大的整数 $\ell$，
以下对象存在：包含指定子空间 $K=\F_2^\ell$ 的有限维二元空间
$\mathcal V$、映射 $C:\mathcal V\to K$、以及 $\mathcal V$ 上的有理
概率分布 $\mu$。它们满足
\begin{equation}\label{eq:latent-equivariance}
 C(x+k)=C(x)+k
 \qquad(x\in\mathcal V,\ k\in K),
\end{equation}
并且，对独立的均匀 $x\in\mathcal V$ 与 $a\sim\mu$，
\begin{equation}\label{eq:latent-small-change}
 \Pr[C(x+a)\ne C(x)]\le p_*.
\end{equation}
此外，每个满足 $\dim(G|_K)\ge r_*$ 的线性子空间
$G\le\mathcal V^*$ 都满足
\begin{equation}\label{eq:latent-detection}
 \Pr_{a\sim\mu}[\text{存在 }g\in G\text{ 使 }g(a)\ne0]
 \ge\frac18.
\end{equation}
特别地，$r_*$ 在选取 $\ell$ 之前固定。对有理 $p_*$，该阈值与所有
有限数据都可由一个会终止的确定性算法找到。
\end{lemma}

我们先构造一个输出空间 $B$ 固定且更小的 gadget：其中的平移诱导
$B$ 的平移。在把诱导同一平移的平移等同之后，该 gadget 将检测每个
限制填满 $B^*$ 的线性族。随后我们组合由单射 $B\to K$ 索引的副本：
$K$ 上限制秩足够高的族在大多数副本上限制到整个 $B^*$。这第二步
容许 $K$ 的每个足够大的维数，同时保持秩阈值固定。

固定 $B$ 的构造从二次块开始。在独立均匀输入上，其某个更新空间中
的均匀非零扰动以概率 $\theta$ 保持非线性输出不变。一个适当一般的
秩 $r$ 线性观察族在同一更新空间上至多损失一维，而均匀选取该空间
导致损失的概率至多为 $3r\theta$。一个调和势把这一点转化为每层
至多 $3\theta$ 的期望损失。坐标的随机变化使非一般的族在下降过程
中变稀少，即使后面的线性映射依赖前面的选择。因此级联使非线性变化
概率变小，同时留下常数概率的正秩。在最后的叶子处，存活下来的族
以常数概率检测出一个独立的均匀符号。

\subsection{二次块}

固定整数 $r_0\ge1$，其规模最终只依赖 $p_*$。记 $H_0=0$ 与
$H_r=\sum_{j=1}^r 1/j$。设 $F=\F_{2^d}$、$q=2^d$，并把
$X=B=F^3$ 视为二元空间，各维数 $b=3d$。整数 $d$ 将按 $r_0$ 选得
足够大。

用迹配对把 $X^*$ 与 $B^*$ 都等同于 $F^3$：向量 $z\in F^3$ 表示
特征 $y\mapsto\operatorname{Tr}_{F/\F_2}(z\cdot y)$，其中点号表示
域点积。我们简写迹为 $\operatorname{Tr}$。迹配对非退化，且
$\operatorname{Tr}(t^2)=\operatorname{Tr}(t)$。$F$ 的每个元素都有
唯一的平方根。

定义
\begin{equation}\label{eq:latent-quadratic}
 Q(x_1,x_2,x_3)=(x_2x_3,x_1x_3,x_1x_2),
 \qquad
 D(x,a)=Q(x+a)+Q(x)+Q(a).
\end{equation}
块可观测为
\[
 C_{\mathrm{blk}}:X\oplus B\longrightarrow B,
 \qquad C_{\mathrm{blk}}(x,y)=y+Q(x).
\]
它在竖直的 $B$ 副本下等变：$C_{\mathrm{blk}}(x,y+k)=C_{\mathrm{blk}}(x,y)+k$。
对 $z\in B^*$，标量观察 $z\circ C_{\mathrm{blk}}$ 在该平移下改变
$z(k)$。我们将把它与限制到竖直 $B$ 上同为 $z$ 的线性泛函比较。

对 $F^3$ 的一维 $F$-子空间 $A=Fv$，令
\begin{equation}\label{eq:latent-block-subspaces}
 D_A=\{y\in B:v\cdot y=0\},
 \qquad
 U_A=\{(a,Q(a)+w):a\in A,\ w\in D_A\}\subseteq X\oplus B.
\end{equation}
这些定义与所选的非零生成元 $v$ 无关。这样的线有 $q^2+q+1$ 条；令
\begin{equation}\label{eq:latent-theta}
 \theta=\frac{1}{q^2+q+1}.
\end{equation}

\begin{lemma}\label{lem:latent-block}
对每条线 $A$，空间 $U_A$ 是维数 $b$ 的二元线性空间。对
$a\in A$ 与 $x\in X$ 有 $D(x,a)\in D_A$；对 $a,a'\in A$ 有
$D(a,a')=0$。$D_A$ 在 $B^*=F^3$ 中的迹零化子是 $A$，且诸 $U_A$
共同张成 $X\oplus B$。
\end{lemma}

\begin{proof}
$D(x,a)$ 的三个坐标是
\[
 x_2a_3+a_2x_3,\qquad
 x_1a_3+a_1x_3,\qquad
 x_1a_2+a_1x_2.
\]
当 $a=tv$ 时，$v\cdot D(x,a)$ 中每一项出现两次，故其和为零。当两个
自变量都是 $v$ 的倍数时，极化式的每个坐标为零。于是 $Q|_A$ 是二元
可加的，且 $U_A$ 由 $A\oplus D_A$ 的参数化是二元线性双射。其维数
为 $d+2d=b$。

每个 $z\in A$ 在迹配对下零化 $D_A$。该零化子的二元维数为
$3d-2d=d$，故恰为 $A$。最后，对每条坐标线 $A=Fe_i$，空间 $U_A$
包含 $(Fe_i,0)$（因为 $Q(te_i)=0$）以及 $\{0\}\oplus D_A$。这三条
坐标线与其竖直超平面共同张成 $X\oplus B$。
\end{proof}

\paragraph{非线性响应。}
$U_A$ 的形状使更新的二次项消去。对 $(a,Q(a)+w)\in U_A$，
\[
 C_{\mathrm{blk}}\bigl((x,y)+(a,Q(a)+w)\bigr)
       +C_{\mathrm{blk}}(x,y)
 =w+D(x,a).
\]
由 Lemma~\ref{lem:latent-block}，右端落在 $D_A$ 中。下面的计算给出
递归将要使用的精确律。

\begin{lemma}[对一次块更新的响应]\label{lem:latent-block-noise}
固定线 $A$。对独立的均匀 $(x,y)\in X\oplus B$ 与
$\zeta\in U_A\setminus\{0\}$，
\[
 \Pr\bigl[C_{\mathrm{blk}}((x,y)+\zeta)
             \ne C_{\mathrm{blk}}(x,y)\bigr]=1-\theta.
\]
\end{lemma}

\begin{proof}
先取 $\zeta$ 在整个 $U_A$ 中均匀，仍与均匀聚合独立。其参数化
$(a,Q(a)+w)$ 中 $a\in A$ 与 $w\in D_A$ 独立均匀。对每个固定的
$x,a$，输出差 $w+D(x,a)$ 在 $D_A$ 中均匀，故以概率 $1-q^{-2}$ 非零。
零扰动占质量 $q^{-3}$ 且从不改变输出。对非零扰动条件化给出
\[
 \frac{1-q^{-2}}{1-q^{-3}}
 =\frac{q(q+1)}{q^2+q+1}=1-\theta.
\]
\end{proof}

该引理处理独立的均匀聚合与均匀非零更新。当更新来自递归的更低层时，
我们会在应用引理之前验证这一条件律。

\paragraph{线性限制。}
固定二元子空间 $S\le B^*$。限制到竖直 $B$ 上为 $z\in S$ 的
$X\oplus B$ 上的线性泛函有唯一形式 $(\gamma_z,z)\in X^*\oplus B^*$。
于是「为每个 $z\in S$ 二元线性地选取一个这样的延拓」由二元线性映射
$\gamma:S\to X^*$ 指定。我们问这些所选延拓中哪些限制到 $U_A$ 上
为零。该限制恰在其核包含非零元素时失去秩。

\begin{lemma}[块限制的核]\label{lem:latent-block-kernel}
设 $S\le B^*$ 且 $\gamma:S\to X^*$ 二元线性。对每条线 $A$，定义
\[
 \kappa_{A,\gamma}:S\longrightarrow U_A^*,
 \qquad \kappa_{A,\gamma}(z)=(\gamma_z,z)|_{U_A}.
\]
在迹等同下，把 $\gamma_z(x)=\operatorname{Tr}(g(z)\cdot x)$ 写成
相应的二元线性映射 $g:S\to F^3$。对每个非零 $z\in S$，
\[
 \kappa_{A,\gamma}(z)=0
 \quad\Longleftrightarrow\quad
 A=Fz
 \quad\text{且}\quad
 z\cdot g(z)=\sqrt{z_1z_2z_3}.
\]
\end{lemma}

\begin{proof}
若 $z$ 在核中，则它零化每个竖直向量 $(0,w)$（$w\in D_A$）。
Lemma~\ref{lem:latent-block} 给出 $z\in A$，而 $z\ne0$ 意味着
$A=Fz$。在其余向量 $(tz,Q(tz))$（$t\in F$）上，该泛函取值
\begin{equation}\label{eq:latent-kernel-trace}
 \operatorname{Tr}\bigl(t\,g(z)\cdot z+t^2z_1z_2z_3\bigr).
\end{equation}
$z\cdot Q(tz)$ 的三个被加项相等，而特征二中三等于一。另外
\[
 \operatorname{Tr}(t^2c)=\operatorname{Tr}(t\sqrt c).
\]
迹配对的非退化性表明：Equation~\eqref{eq:latent-kernel-trace} 对
每个 $t$ 为零，当且仅当 $z\cdot g(z)=\sqrt{z_1z_2z_3}$。结合
$z\in A$，该条件也是充分的：竖直部分已被零化。
\end{proof}

\subsection{一般的特征空间}

核判据提示两种充分控制。我们将使用这样的一族：没有域线包含两个
非零特征，且对任一延拓选择，只有少数特征满足对齐方程。对该选择
的一致性随后允许提升依赖较早的随机坐标变化。

\begin{definition}
维数 $r\ge1$ 的二元子空间 $S\le F^3$ 称为\emph{一般的}，若 $S$ 中
没有两个不同的非零元素在 $F$ 上成比例，且对每个二元线性映射
$g:S\to F^3$，至多 $3r$ 个非零 $z\in S$ 满足
\begin{equation}\label{eq:latent-alignment}
 z\cdot g(z)=\sqrt{z_1z_2z_3}.
\end{equation}
\end{definition}

在选取均匀线 $A$ 之前先固定一个映射 $\gamma:S\to X^*$，并令
$r_A=\dim\im\kappa_{A,\gamma}$，其中 $r=\dim S>0$。
Lemma~\ref{lem:latent-block-kernel} 现在给出损失界。若 $S$ 是一般
的，每条域线至多含 $S$ 的一个非零元素，故核的维数至多为一。对与
$\gamma$ 关联的固定映射 $g$，至多 $3r$ 个特征满足
Equation~\eqref{eq:latent-alignment}；故至多 $3r$ 条线可能造成任何
秩损失。因此
\begin{equation}\label{eq:latent-generic-rank-loss}
 r-r_A\le1,\qquad
 \Pr_A[r_A<r]\le3r\theta,\qquad
 \E_A[H_r-H_{r_A}]\le3\theta.
\end{equation}
最后的界用了 $H_r-H_{r-1}=1/r$。
对任意的秩 $r$ 空间，损失仍要求对某非零 $z\in S$ 有 $A=Fz$。因此
\begin{equation}\label{eq:latent-arbitrary-rank-loss}
 \Pr_A[r_A<r]\le(2^r-1)\theta\le2^r\theta,
 \qquad H_r-H_{r_A}\le H_r.
\end{equation}
这些估计关心限制下的秩。只有当正秩的族到达叶子后，最终噪声的检测
才会随之而来。

还需证明具有所述一般性的空间大量存在。

\begin{lemma}\label{lem:latent-generic}
对每个固定的 $r\ge1$，$F^3$ 中均匀随机选取的二元 $r$ 维子空间在
$d\to\infty$ 时为一般的概率趋于 1。
\end{lemma}

\begin{proof}
选独立的均匀 $x_{ij}\in F$（$1\le i\le3$，$1\le j\le r$），并用
\[
 z_i(t)=\sum_{j=1}^r x_{ij}t_j,
 \qquad t\in\F_2^r
\]
参数化一个二元线性映射。每个非零输入的像在 $F^3$ 中均匀，故该映射
非单射的概率至多 $(2^r-1)q^{-3}$。在单射的条件下，其像是均匀的
二元 $r$-子空间：每个这样的子空间有相同数目的有序二元基。不同的
非零二元输入在 $F$ 上也线性无关，故其像是 $F^3$ 中独立均匀元素。
因此除非单射事件外，并界表明非比例条件以 $O_r(q^{-2})$ 的概率
失败。

剩下要对所有 $g$ 同时证明对齐条件。固定 $\F_2^r$ 中超过 $3r$ 个
互异非零输入的集合 $E$。在单射参数化下，二元线性 $g$ 在一组基上的
取值可以是 $F^3$ 中的任意向量。于是 $E$ 上的对齐蕴含向量
\[
 c=\bigl(\sqrt{z_1(t)z_2(t)z_3(t)}\bigr)_{t\in E}
\]
落在具有 $3r$ 列的矩阵
\[
 M_{t,(i,j)}=t_jz_i(t),
 \qquad t\in E,\quad 1\le i\le3,\quad 1\le j\le r
\]
的 $F$-列张成中。反之，属于该列张成即给出这样一个 $g$ 的基取值。
特别地，这一列隶属条件已经涵盖了所有可能的映射 $g$。

先把 $x_{ij}$ 视为独立不定元并令 $L=\F_2(x_{ij})$。我们断言：即使
把系数扩到代数闭包 $\overline L$，$c$ 也不在 $M$ 的列张成中。以下
$r^3$ 个元素在 $L$ 上线性无关：
\[
 \sqrt{x_{1j}x_{2k}x_{3l}},
 \qquad 1\le j,k,l\le r.
\]
事实上，通分后，一个依赖关系将是这些不定元的平方根之间的多项式
恒等式。所示单项式有不同的奇偶向量，而关于 $x_{ij}$ 的多项式系数
在这些平方根中只有偶次幂。不同奇偶向量的项不可能相消。

$c$ 在这些无关元素中的展开的系数向量为
\[
 v_{jkl}=(t_jt_kt_l)_{t\in E}.
\]
若 $c$ 属于 $\overline L$ 上的列张成，则每个在 $L$ 上零化 $M$ 的
行向量都会零化 $c$。刚证的无关性蕴含它会零化每个 $v_{jkl}$。由
有限维空间 $L^E$ 中的双重零化，每个 $v_{jkl}$ 都会属于 $L$ 本身上
的列张成。

对 $h\ge1$，设 $V_h\le L^E$ 为 $t_1,\ldots,t_r$ 的二元坐标中次数
至多 $h$ 的所有正次数单项式的赋值张成。$v_{jkl}$ 中的重复指标给出
一次、二次与三次的所有单项式，因为在赋值集上 $t_i^2=t_i$。前述
结论与列的形状因此蕴含
\[
 V_3\subseteq\operatorname{col}_L(M)\subseteq V_2\subseteq V_3.
\]
三个空间相等。乘以任意 $t_i$ 把 $V_2$ 送入 $V_3$，故保持其公共
值。由于该空间含一次单项式，它含每个正次数单项式的赋值。这些赋值
张成整个 $L^E$：对任意非零二元点，其通常的示性多项式是因子 $t_i$
与 $1+t_i$ 的乘积且至少含一个因子 $t_i$，从而无常数单项式；其在
$E$ 上的限制给出相应的点示性。于是列张成的维数将为 $|E|>3r$，
矛盾。

为从这一代数命题过渡到有限域，把 $M$ 与 $c$ 的元素都平方。所得
矩阵有多项式元素，附加列的元素为 $z_1(t)z_2(t)z_3(t)$。令
$h=\rank_L M\le3r$。平方在 $\overline L$ 上保持秩，故增广平方
矩阵有 $h+1$ 阶非零子式。该子式是次数至多 $2h+3\le6r+3$ 的非零
多项式。所有 $h+1$ 阶的未增广子式恒为零，故未增广秩在特化时不会
增大。因此特化后的列隶属迫使这个特定的非零多项式为零。对独立的
均匀 $x_{ij}\in F$，由多项式零点界~\cite[Corollary~1]{Schwartz80}
（另见~\cite[Section~3.1]{Zippel79} 较早的稀疏多项式表述），该事件
的概率至多 $(6r+3)/q$。

若对齐条件失败，则它在某个大小为 $3r+1$ 的集合 $E$ 上失败。这样的
集合只有有限多个（只依赖 $r$）；若 $2^r-1\le3r$ 则没有，条件自动
成立。并界加上对单射的条件化即证引理。
\end{proof}

一般性因此控制所有可能的映射 $g$。现在固定一个域，使截至 $r_0$ 的
每个秩上非一般空间都稀少；此阶段不需要选取任何线性映射。取正的
$\eps_0$ 满足
\begin{equation}\label{eq:latent-generic-error}
 H_{r_0}2^{r_0}\eps_0\le1,
 \qquad \eps_0<1.
\end{equation}
由 Lemma~\ref{lem:latent-generic}，可以固定足够大的 $d$ 使
$b=3d\ge r_0$，且对每个 $1\le r\le r_0$，$B^*$ 的均匀秩 $r$ 子空间
非一般的概率至多 $\eps_0$。特别地，存在一般的秩 $r_0$ 子空间，
我们一劳永逸地固定它。

\subsection{级联及其噪声}

在级联二次块时，我们保持输出空间 $B$ 固定。每层选取一条噪声分支。
下一个引理计算非线性变化概率衰减的速度；随后我们将沿同一分支追踪
线性限制的秩。

我们构造二元输入空间 $P_t$、平移子空间 $R_t\le P_t$、线性满射
$\lambda_t:R_t\to B$ 与映射 $C_t:P_t\to B$，使
\begin{equation}\label{eq:latent-recursive-equivariance}
 C_t(u+h)=C_t(u)+\lambda_t(h)
 \qquad(h\in R_t).
\end{equation}
于是平移 $h\in R_t$ 表示输出平移 $\lambda_t(h)$；不同的平移可以
表示同一输出平移。这种平移作用的形式在聚合下保持。最后，对值为
零的平移作商，将把输出平移变成输入的实际子空间。

在高度零处取 $P_0=R_0=B$ 与 $C_0=\lambda_0=\operatorname{id}_B$。
设 $\mathcal I$ 由所有对 $(A,J)$ 组成，其中 $A$ 是 $F^3$ 中的一条
线，$J:B\to U_A$ 是二元线性同构。每条线有相同数目的这样的同构。

我们在单个固定的父映射中为每个 $(A,J)$ 安排一个子节点。选取均匀
的噪声子节点随即采样一条均匀的线与一个新鲜的均匀同构。
对 $t\ge1$，取 $P_t=\prod_{(A,J)\in\mathcal I}P_{t-1}$ 并设
\begin{equation}\label{eq:latent-recursive-output}
 (x,y)=\sum_{A,J}J C_{t-1}(u_{AJ}),
 \qquad C_t(u)=C_{\mathrm{blk}}(x,y)=y+Q(x).
\end{equation}
子节点平移的元组通过其嵌入的输出平移之和改变该聚合。为使父节点
的变化是纯输出平移，我们要求该和的 $X$-分量为零。据此定义
\begin{equation}\label{eq:latent-recursive-shifts}
 R_t=\left\{(h_{AJ})\in\prod_{A,J}R_{t-1}:
       \sum_{A,J}J\lambda_{t-1}(h_{AJ})\in\{0\}\oplus B\right\},
\end{equation}
并当该聚合为 $(0,k)$ 时令 $\lambda_t(h)=k$。这些是线性定义，且
Equation~\eqref{eq:latent-recursive-equivariance} 立即由
Equation~\eqref{eq:latent-recursive-output} 得到。满射性由
Lemma~\ref{lem:latent-block} 归纳可得：子节点的输出平移可以任意
选取，而其聚合映射是到 $X\oplus B$ 的满射。
Equation~\eqref{eq:latent-recursive-equivariance} 与满射性还表明
$C_t$ 把均匀输入送到 $B$ 的均匀元素。

设 $\xi_t$ 为 $P_t$ 中如下的加性噪声。高度零处它在 $B$ 中均匀。
在正高度处，选均匀的子节点 $(A,J)$，把 $\xi_{t-1}$ 的独立副本嵌入
该子节点，其余子节点置零。等价地，在乘积树中选均匀的叶子，并在该
叶子处加 $B$ 的均匀符号。每条路径选择可以采样为一次均匀的 $A$ 加
一次均匀的 $J$，逐层独立。所有这些分布都是有理的。
Figure~\ref{fig:latent-recursion} 展示递归的一层与所选的噪声分支。
\begin{figure}[!htbp]
\centering
\begin{tikzpicture}[
 every node/.style={font=\small,align=center},
 block/.style={draw=linkblue!80!black,rounded corners=2pt,
 text width=3.3cm,inner sep=7pt,minimum height=14mm},
 flow/.style={-{Latex[length=2mm]},thick}]
\node[block] (left) at (0,0)
 {一个子节点输入\\$u_{AJ}\in P_{t-1}$\\$J C_{t-1}(u_{AJ})$};
\node[block,draw=red!70!black] (selected) at (4.8,0)
 {被选中的子节点输入\\$u_{A_*J_*}+\xi_{t-1}$\\$J_*C_{t-1}(u_{A_*J_*}+\xi_{t-1})$};
\node[block] (right) at (9.6,0)
 {其余子节点输入\\$u_{AJ}\in P_{t-1}$\\$J C_{t-1}(u_{AJ})$};
\node at (2.4,0) {$\cdots$};
\node at (7.2,0) {$\cdots$};
\node[block,text width=7cm] (sum) at (4.8,-2)
 {对所有子节点 $(A,J)\in\mathcal I$ 求和\\
 噪声聚合 $(x',y')=(x,y)+J_*\Delta$};
\draw[flow] (left.south) -- (sum.north west);
\draw[flow,draw=red!70!black] (selected) -- (sum);
\draw[flow] (right.south) -- (sum.north east);
\node[block,text width=7cm] (out) at (4.8,-3.8)
 {噪声父输出 $y'+Q(x')\in B$};
\draw[flow] (sum) -- (out);
\end{tikzpicture}
\caption{潜在递归的一层（在其加性噪声下）。输入空间是全部子节点
指标 $(A,J)$ 上的乘积；只有一个均匀选中的子节点被扰动。每个子
节点输出被嵌入 $U_A\le X\oplus B$，全部嵌入输出相加，二次映射在
父节点处作用一次。此处 $(x,y)$ 是初始聚合，$\Delta$ 是被选子节点
的输出差，故噪声聚合为 $(x',y')=(x,y)+J_*\Delta$。被省略的子节点
遵循同一规则。高亮的扰动可能在父节点处消失。}
\label{fig:latent-recursion}
\end{figure}

\FloatBarrier

\begin{lemma}\label{lem:latent-noise}
对独立的均匀 $u\in P_t$ 与噪声 $\xi_t$，
\begin{equation}\label{eq:latent-change-probability}
 p_t:=\Pr[C_t(u+\xi_t)\ne C_t(u)]
     =(1-q^{-3})(1-\theta)^t.
\end{equation}
\end{lemma}

\begin{proof}
高度零处的恒等式显然。考虑更高的一层，并在所选子节点 $(A,J)$ 上
条件化。记其输出差为 $\Delta$。所有子节点副本有相同的函数与相同
的内部噪声律，故 $\Delta$ 的律与所选指标无关。

我们还需要子节点初始输出的条件律。对任意固定的子节点扰动 $a$，
Equation~\eqref{eq:latent-recursive-equivariance} 给出
\[
 C_{t-1}(u+a+h)+C_{t-1}(u+h)
 =C_{t-1}(u+a)+C_{t-1}(u)
 \qquad(h\in R_{t-1}).
\]
于是该差在这些平移下不变，而初始输出被 $B$ 的每个元素平移。在
均匀输入下，即使内部扰动也被指定，其初始输出在给定差时仍均匀。
其余子节点的初始输出保持独立均匀。由于聚合映射是到 $X\oplus B$ 的
满射，初始聚合 $(x,y)$ 在给定 $(A,J,\Delta)$ 时均匀。特别地，它与
聚合扰动 $J\Delta$ 独立。

在 $\Delta\ne0$ 与 $A$ 的条件下，同构 $J$ 仍均匀，因为 $\Delta$ 的
律与所选子节点指标无关。于是 $J\Delta$ 在 $U_A\setminus\{0\}$ 中
均匀。刚证明的条件均匀性使初始聚合与该更新独立。因此
Lemma~\ref{lem:latent-block-noise} 给出：在子节点差非零的条件下，
父节点变化概率为 $1-\theta$。当 $\Delta=0$ 时聚合不变。于是
$p_t=p_{t-1}(1-\theta)$，恒等式得证。
\end{proof}

\subsection{线性提升与秩损失}

非线性变化概率现在随高度衰减。对固定 $B$ 的 gadget，剩余目标是：
只要族的限制填满 $B^*$ 就能检测。我们将通过平移作用追踪这些输出
特征的每一种二元线性延拓选择来证明它。

对 $S\le B^*$，$P_t$ 上的\emph{提升}指二元线性映射
$z\mapsto\psi_z$（从 $S$ 到 $P_t^*$），满足
\begin{equation}\label{eq:latent-lift-definition}
 \psi_z|_{R_t}=z\circ\lambda_t.
\end{equation}
定义域为 $B^*$ 的提升称为\emph{满提升}。接下来我们描述父节点处的
每一种可能的提升，不对其选取方式作任何假设。

\begin{lemma}\label{lem:latent-lift-factorization}
设 $P_t$（$t\ge1$）上给定了定义域为 $S$ 的提升。对每个子节点
$(A,J)$，限制到该子节点的平移空间会经由 $\lambda_{t-1}$ 分解。所得
子节点输出空间 $B$ 上的特征形如
\begin{equation}\label{eq:latent-child-characters}
 \phi_{AJ}(z)=J^*\bigl((\gamma_z,z)|_{U_A}\bigr)\in B^*,
\end{equation}
其中 $\gamma:S\to X^*$ 是对所有子节点公共的单一二元线性映射。
\end{lemma}

\begin{proof}
支撑在单个子节点上且属于 $\ker\lambda_{t-1}$ 的平移，是输出平移为
零的父平移。每个 $\psi_z$ 在其上为零，证得第一个断言。

设 $T=\prod_{A,J}B$ 为子节点输出平移的空间，聚合映射为
$\mathcal A:T\to X\oplus B$。该映射是满的。子节点特征之和减去 $z$
作用在 $\mathcal A$ 的第二分量上，在 $\mathcal A$ 的第一分量为零
时为零：每个这样的平移元组都可由子平移提供，给出一个父平移。因此
该差唯一地经由 $\mathcal A$ 的第一分量（满射到 $X$）分解。称此
特征为 $\gamma_z$。唯一性表明 $\gamma_z$ 线性依赖 $z$。在支撑于
单个子节点的元组上求值即得 Equation~\eqref{eq:latent-child-characters}。
\end{proof}

选定子节点后，令 $S'=\im\phi_{AJ}$ 并取 $\phi_{AJ}$ 的线性截断
$\iota:S'\to S$。把 $\psi_{\iota(z')}$ 限制到该子节点的输入上，即
得 $P_{t-1}$ 上定义域为 $S'$ 的提升。我们可以为所有这样的截断固定
确定性规则。特别地，它们只依赖已选的路径，从不依赖未来的路径指标
或最终噪声符号。下降树时，每个存活的泛函都是原线性族某个成员的
限制。

分解由此把递归与局部限制问题连接起来。对固定的父提升，其唯一的
映射 $\gamma$ 对每条出线 $A$ 公共，且
$\phi_{AJ}=J^*\kappa_{A,\gamma}$。由于 $J^*$ 是同构，不改变核与秩。
因此当出线均匀选取时，可以应用
Equations~\eqref{eq:latent-generic-rank-loss}
与~\eqref{eq:latent-arbitrary-rank-loss}。下一个证明表明：入射
同构使一般性在后续每个正秩处都典型，即使提升本身依赖较早的路径。

\begin{proposition}\label{prop:latent-full-lift-detection}
令
\begin{equation}\label{eq:latent-depth}
 s=\left\lfloor\frac{H_{r_0}}{8\theta}\right\rfloor.
\end{equation}
$P_s$ 上的每个满提升以至少 $1/4$ 的概率检测 $\xi_s$：
\begin{equation}\label{eq:latent-full-lift-detection}
 \Pr[\text{存在 }z\in B^*\text{ 使 }\psi_z(\xi_s)\ne0]
 \ge\frac14.
\end{equation}
\end{proposition}

\begin{proof}
把顶端的满提升限制到维数 $r_0$ 的固定一般子空间。沿随机叶子路径
行进，在每个所选子节点处把提升换成刚描述的像提升。记其经 $j$ 步
后的定义域为 $S_j$，$r_j=\dim S_j$，初始时 $j=0$ 处 $r_j=r_0$。
灭绝后继续用零空间。

我们详细说明允许在每层使用一般性的条件化。设 $\mathcal F_j$ 为
第 $j$ 步的出线选择 $(A_j,J_j)$ 之前所有路径选择的信息。当前的
空间、提升与映射 $\gamma_j$ 都是 $\mathcal F_j$-可测的。在
$\mathcal F_j$ 与 $A_j$ 的条件下，子空间
\[
 E_{j,A_j}
 =\im\bigl[z\mapsto(\gamma_{j,z},z)|_{U_{A_j}}\bigr]
 \le U_{A_j}^*
\]
固定。下一个定义域与其秩为
\[
 S_{j+1}=J_j^*E_{j,A_j},
 \qquad r_{j+1}=\dim E_{j,A_j}.
\]
秩在 $J_j$ 揭晓之前固定。均匀二元同构 $J_j:B\to U_{A_j}$ 使
$S_{j+1}$ 成为 $B^*$ 中该秩的均匀随机子空间。因此，若 $B_j$ 为
「$r_j>0$ 且 $S_j$ 非一般」的示性变量，则
\begin{equation}\label{eq:latent-fresh-orientation}
 \E[B_{j+1}\mid\mathcal F_j,A_j]\le\eps_0.
\end{equation}
下一个提升可以依赖 $J_j$。这对论证不构成限制：一般性界定所选
子空间上的每个二元线性映射 $g$，包括由该提升给出的映射。

令 $\Delta_j=H_{r_j}-H_{r_{j+1}}\ge0$。在 $\mathcal F_j$ 的条件
下，出线 $A_j$ 独立均匀，而上述秩损失估计适用于固定的当前提升。
因此
\begin{equation}\label{eq:latent-conditional-harmonic}
 \E[\Delta_j\mid\mathcal F_j]
 \le3\theta+H_{r_0}2^{r_0}\theta B_j.
\end{equation}
对 $j=0$，空间是一般的，故其期望下降至多 $3\theta$。对 $j\ge1$，
用 Equation~\eqref{eq:latent-fresh-orientation} 把
Equation~\eqref{eq:latent-conditional-harmonic} 对前面的同构取平均。
这给出无条件的界
\[
 \E\Delta_j
 \le\bigl(3+H_{r_0}2^{r_0}\eps_0\bigr)\theta
 \le4\theta.
\]
没有假设入射取向与其后续提升之间的独立性。

对期望下降求和并用 Equation~\eqref{eq:latent-depth}：
\[
 \E[H_{r_0}-H_{r_s}]\le4s\theta\le\frac{H_{r_0}}2.
\]
灭绝需要总下降 $H_{r_0}$，故 $\Pr[r_s=0]\le1/2$。在叶子处，
$P_0=R_0=B$，其提升由特征本身组成。在任意存活的秩 $r_s>0$ 的条件
下，$B$ 中的独立均匀符号以概率 $1-2^{-r_s}\ge1/2$ 被检测。每个
存活的泛函都是原族沿所选叶子的限制，故这样的检测蕴含原族对嵌入
噪声 $\xi_s$ 的检测。两个界相乘即证
Equation~\eqref{eq:latent-full-lift-detection}。
\end{proof}

\subsection{商与放大}

一旦平移作用由实际子空间实现，固定 $B$ 的构造就完成了。令
\[
 V_0=P_s/\ker\lambda_s,
\]
其中核取在 $R_s\le P_s$ 内，并设 $q_0:P_s\to V_0$ 为商映射。
Equation~\eqref{eq:latent-recursive-equivariance} 使 $C_s$ 在该核下
不变，故它下降为映射 $\overline C:V_0\to B$。商
$R_s/\ker\lambda_s$ 由 $\lambda_s$ 等同于 $B$，给出包含
$i_B:B\to V_0$ 使
\[
 q_0(h)=i_B(\lambda_s h)\qquad(h\in R_s).
\]
下文把该子空间等同于 $B$。其平移作用满足
\[
 \overline C(x+b')=\overline C(x)+b'
 \qquad(b'\in B).
\]
设 $\nu$ 为 $\xi_s$ 的噪声律的前推。$P_s$ 上的均匀输入前推为
$V_0$ 上的均匀输入，故输出变化概率仍为 $p_s$。

限制为整个 $B^*$ 的每个线性子空间 $G_0\le V_0^*$ 都满足
\begin{equation}\label{eq:latent-base-detection}
 \Pr_{a\sim\nu}[\text{存在 }g\in G_0\text{ 使 }g(a)\ne0]
 \ge\frac14.
\end{equation}
事实上，取限制的线性截断 $\sigma:B^*\to G_0$ 并定义
$\psi_z=q_0^*\sigma(z)=\sigma(z)\circ q_0$。对 $h\in R_s$，
\[
 \psi_z(h)=\sigma(z)(q_0h)=z(\lambda_s h).
\]
于是该拉回是满提升。而且 $\psi_z(\xi_s)=\sigma(z)(q_0\xi_s)$，
故提升的检测就是所选截断对前推噪声的检测。
Proposition~\ref{prop:latent-full-lift-detection} 给出该界，而截断
的检测蕴含 $G_0$ 的检测。

\begin{proof}[Lemma~\ref{lem:latent} 的证明]
选足够大的 $r_0$ 使
\[
 \exp(1-H_{r_0}/8)\le p_*.
\]
再按 Equation~\eqref{eq:latent-generic-error} 选 $\eps_0$，并用
Lemma~\ref{lem:latent-generic} 固定 $d$（从而 $b=3d$ 与 $\theta$）。
构造 Equation~\eqref{eq:latent-depth} 的高度 $s$ 与上述商
$(V_0,B,\overline C,\nu)$。Lemma~\ref{lem:latent-noise} 给出
\begin{equation}\label{eq:latent-small-base-change}
 p_s\le(1-\theta)^s
 \le\exp(-\theta s)
 \le\exp(\theta-H_{r_0}/8)
 \le p_*.
\end{equation}
令 $r_*=b+2$。至此所有选取都只依赖 $p_*$。

现在放大输出空间。使用随机单射 $B\to K$ 的要点是：$K$ 上限制秩
至少 $b+2$ 的族沿大多数单射限制到整个 $B^*$。于是为每个单射取
基 gadget 的一个副本，即可把其满限制检测界转化为所需的高秩界。

固定 $\ell\ge b+2$ 并记 $K=\F_2^\ell$。设 $\mathcal L$ 为所有二元
线性单射 $L:B\to K$ 的有限集。取乘积 $P=\prod_{L\in\mathcal L}V_0$
并定义
\[
 \widetilde C((x_L)_L)=\sum_{L\in\mathcal L}L\overline C(x_L).
\]
平移子空间的乘积 $T=\prod_L B\le P$ 通过线性映射
\[
 \Lambda:T\longrightarrow K,
 \qquad \Lambda((b_L)_L)=\sum_L Lb_L
\]
作用于该输出。该映射是满的，因为 $K$ 的每个非零向量都在某个单射
$B\to K$ 的像中。对其在 $T$ 内的核作商：
\[
 \mathcal V=P/\ker\Lambda.
\]
映射 $\widetilde C$ 下降为 $C:\mathcal V\to K$，而
$T/\ker\Lambda$ 等同于所需的平移子空间 $K$。这证明
Equation~\eqref{eq:latent-equivariance}。

在 $P$ 上选均匀的 $L\in\mathcal L$，仅对该副本用 $\nu$ 的独立样本
更新。设 $\mu$ 为取商后的律。在均匀乘积输入上，输出和变化当且仅当
所选副本的输出变化（因为 $L$ 是单射）。均匀输入前推为商上的均匀
输入，故变化概率仍为 $p_s$。Equation~\eqref{eq:latent-small-base-change}
证明 Equation~\eqref{eq:latent-small-change}。这里用到的每个噪声
分布都由有限均匀选取经前推得到，因而是有理的。

还需对任意 $G\le\mathcal V^*$ 建立检测。设 $S=G|_K\le K^*$，维数
$r\ge r_*$。把 $G$ 拉回到 $P$，再限制到由 $L$ 索引的 $V_0$ 副本；
把所得 $V_0^*$ 的子空间记为 $G_L$。它限制到该副本的平移子空间
$B$ 上恰为
\[
 G_L|_B=L^*S.
\]
缘由：支撑在副本 $L$ 上的平移 $b'\in B$ 在商中表示平移 $Lb'\in K$，
故限制到 $z\in S$ 的特征在该处取值 $z(Lb')$。

映射 $L^*:S\to B^*$ 不满，当且仅当存在非零 $b'\in B$ 使
$Lb'\in S^\perp$。对每个这样的固定 $b'$，均匀单射 $L$ 使 $Lb'$
在 $K$ 的非零向量中均匀。故
\[
 \Pr_L[Lb'\in S^\perp]
 =\frac{2^{\ell-r}-1}{2^\ell-1}
 \le2^{-r}.
\]
并界给出
\[
 \Pr_L[L^*S\ne B^*]
 \le(2^b-1)2^{-r}
 \le2^{b-r}\le\frac14.
\]
在每个其余的副本上，Equation~\eqref{eq:latent-base-detection} 给出
$G_L$ 从而 $G$ 对嵌入噪声的检测概率至少 $1/4$。对所选副本取平均，
检测概率至少 $\frac34\cdot\frac14=\frac{3}{16}\ge\frac18$，证明了
Equation~\eqref{eq:latent-detection}。构造对每个 $\ell\ge b+2$
有效，且 $r_*=b+2$ 已固定，如所需。

为完备起见：当 $p_*$ 有理时，所有选取都是有效的。选整数 $m\ge0$
使 $2^{-m}\le p_*$，再选 $r_0$ 使 $H_{r_0}\ge8(m+1)$；由
$e^{-m}\le2^{-m}$，这保证 $r_0$ 的初始界。选满足
Equation~\eqref{eq:latent-generic-error} 的有理 $\eps_0$。对相继的
整数 $d$（$3d\ge r_0$），枚举 $\F_{2^d}^3$ 的二元子空间与一般性
定义中的线性映射。这精确检验一般性，并计算每个秩 $1,\ldots,r_0$
处非一般空间的比例。Lemma~\ref{lem:latent-generic} 保证最终所有
这些比例都至多 $\eps_0$。固定第一个这样的域与一个一般的秩 $r_0$
空间。其余的乘积、同构、单射、商、函数表与有理噪声律都由有限枚举
与二元线性代数获得。这给出所述的会终止的构造。
\end{proof}

\section{外层问题与置换测试}
\label{sec:game}

现在用潜在 gadget 构造置换实例。其顶点编码奇偶方程元组上的仿射表。
当两张表在相同的保留坐标上诱导同一仿射映射（含同一截距）时，我们
把二者精确等同。这一公共表示将使我们在 Section~\ref{sec:matrix}
中比较方程元组与稀疏投影元组。

\subsection{带权奇偶方程与不同位置}

\emph{带权奇偶实例}由有限变量集 \(\mathcal X\)、有限的非空方程出现
集 \(\mathcal E\)、以及和为 1 的正有理权重 \(\omega_e\) 组成。一个
出现是一条有序方程
\[
 x_{v_{e,1}}+x_{v_{e,2}}+x_{v_{e,3}}=b_e,
 \qquad b_e\in\F_2.
\]
其值是被满足方程的加权比例，\(\OPT(\mathcal E)\) 记最大值。即使
有序方程重合，不同的出现也有不同的 ID。

\begin{theorem}[带权奇偶间隙]
\label{thm:parity}
对每个固定的有理 \(\xi>0\)，存在从布尔可满足性到带权奇偶实例的
确定性多项式时间归约，权重为二进制编码的有理数，使得
\[
 \begin{array}{ll}
 \textup{YES：}&\OPT(\mathcal E)\ge 1-\xi,\\
 \textup{NO：}&\OPT(\mathcal E)\le \frac12+\xi.
 \end{array}
\]
运行时间多项式可以依赖 \(\xi\)。
\end{theorem}

这是 \cite[Theorem~5.4 and its proof]{Hastad01} 中 H\aa stad 构造
的间隙推论。更具体地，该证明使用固定的二进验证者误差 \(\delta>0\)，
给出完备性 \(1-\delta\) 与可靠性 \((1+\delta)/2\)，并把验证者的
选择作为方程枚举出来、以其有理概率为权重。取 \(\delta\le\xi\) 即
得所示陈述。该枚举是确定性归约；验证者的随机选择描述了所得的带权
实例。

\begin{lemma}[使位置互异]
\label{lem:distinct-parity}
由带权奇偶实例 \(\mathcal E\) 可以在多项式时间内确定性构造实例
\(\mathcal E'\)：其中每个出现的三个变量 ID 互不相同。其权重是
多项式位长的正有理数，且
\[
 \OPT(\mathcal E')\ge\OPT(\mathcal E),
 \qquad
 \OPT(\mathcal E')
 \le \frac{1+\OPT(\mathcal E)}2+\frac3{100}.
\]
特别地，先取 \(\xi<1/100\) 的 Theorem~\ref{thm:parity}，再做该
预处理，给出完备性至少 \(1-\xi\)、NO 最优值至多 \(4/5\)。
\end{lemma}

\begin{proof}
把每个变量 \(v\) 换成一百个克隆 \((v,c)\)（\(c\in[100]\)）。对每
个原出现，均匀选取其三个克隆，只要各位置 underlying 变量相同就
强制选互异的克隆。列出所得的全部出现，并赋予相应的条件权重。每个
原出现至多有 \(100^3\) 种选择。给所有克隆其原变量的标记即证完备性
断言。

对上界，固定任意克隆标记，令
\[
 t_v=\frac1{100}\sum_{c=1}^{100}(-1)^{x_{(v,c)}}.
\]
先独立选克隆，不以互异性为条件。出现 \(e\) 的满足概率是
\[
 \frac{1+(-1)^{b_e}t_{v_{e,1}}t_{v_{e,2}}t_{v_{e,3}}}{2}.
\]
把每个偏差四舍五入为多数标记，零偏差任意处理。对被该取整赋值违背
的每条方程，分子中的带符号乘积非正，故其满足概率至多 \(1/2\)；
其余方程的概率至多 1。加权平均因此至多
\((1+\OPT(\mathcal E))/2\)。独立选取的克隆之间出现禁用碰撞的概率
由对三对位置的并界至多为 \(3/100\)。以其补为条件会把任何
\([0,1]\) 值函数的期望改变至多 \(3/100\)。界得证。除原出现权重外，
所有条件概率的分母都是常数规模。
\end{proof}

在构造的其余部分，\(\mathcal E\) 表示预处理后的实例，故每条方程
都有三个互异的变量 ID。其\emph{普通「方程对变量」博弈}按 \(\omega\)
采样 \(e\)、从 \([3]\) 均匀采样 \(j\)。第一个证明者收到出现 ID
\(e\)，以该方程的一个满足三元组作答；第二个收到变量 ID
\(v_{e,j}\)，以一个比特作答。接受表示位置 \(j\) 处相等。由于方程
内的 ID 互异，问题对决定了投影约束。

对任意第二证明者策略，其比特构成一个全局赋值。在该赋值违背的方程
上，每个合法的第一证明者答案至少在一个位置不一致。因此这个记作
\(G_{\mathcal E}\) 的普通投影博弈满足
\begin{equation}
 \val(G_{\mathcal E})
 \le 1-\frac{1-\OPT(\mathcal E)}3
 \le \frac{14}{15}
 \quad\text{若 }\OPT(\mathcal E)\le\frac45.
 \label{eq:base-game-gap}
\end{equation}
答案字母表规模分别为 4 与 2。

\subsection{元组、局部答案副本与齐次坐标}

固定正整数 \(k\)。方程问题 \(U=(e_1,\ldots,e_k)\) 从
\(\omega^{\otimes k}\) 有放回采样。其答案空间是 \(k\) 个仿射解平面
的乘积。块总是使用局部副本：即使同一方程或变量 ID 出现在多个块中，
它们的局部答案也是分离的坐标。

对参数 \(0<\beta<1\)，在每个块中独立地：以概率 \(1-\beta\) 保留该
方程，或以概率 \(\beta\) 只保留均匀所选位置上的变量 ID。所得问题
记为 \(O\)，\(\pi\) 是从 \(U\) 上的答案到 \(O\) 上的答案的坐标
投影。问题 \(O\) 在每个有序块位置记录一个方程 ID 或变量 ID。
单例块接受任一比特。

把方程答案用其前两个比特 \((x_i,y_i)\) 表示；第三个是
\(b_{e_i}+x_i+y_i\)。单例用其实际比特表示。加一个公共齐次坐标
\(t_0\)。于是
\[
 L_U=\F_2^{1+2k},\qquad m=1+2k,
\]
而 \(L_O\) 有一个齐次坐标、每个保留方程两个坐标、每个单例一个
坐标。在任一空间中，第一个坐标为 1 的向量恰是合法答案。把读取该
坐标的泛函记为 \(e_0^\top\)。

仿射投影 \(\pi\) 延拓为保持 \(t_0\) 的线性满射
\(\bar\pi:L_U\to L_O\)。在保留的方程块上它是恒等；在投影为单例的
块上，其坐标是
\begin{equation}
 \begin{cases}
 x_i,&\text{位置 }1,\\
 y_i,&\text{位置 }2,\\
 b_{e_i}t_0+x_i+y_i,&\text{位置 }3.
 \end{cases}
 \label{eq:outer-projection}
\end{equation}
把由这些逐块选择得到的每个映射称为\emph{容许投影}；这一类映射与
\(\beta\) 无关。满射性由固定 \(t_0\) 并在每个块中独立选取坐标
得到。对某些块被省略的问题，我们沿用同样的约定。

\subsection{精确表键与折叠}

跨投影问题等同表的做法遵循 Khot--Safra 的虚拟表方法
~\cite[Section~1.1]{KS13}。我们先把每张表用它所经过的坐标表示。
这将给投影表与其拉回相同的表示。

固定 Lemma~\ref{lem:latent} 的潜在数据 \((\mathcal V,K,C,\mu)\)，
其中 \(K=\F_2^\ell\)，并固定线性补 \(\mathcal V=K\oplus W\)。问题
\(Q\) 上的\emph{表}是从其答案空间到 \(\mathcal V\) 的仿射映射，
等价地其唯一线性延拓 \(P\in\Hom(L_Q,\mathcal V)\)。此处 \(Q\)
可以是方程问题或其某个容许投影，同一定义也适用于有省略块的问题。
把 \(P\) 中齐次坐标 \(t_0\) 的系数记为 \(c\)。

在方程块中，检查其斜率映射 \(\F_2^2\to\mathcal V\) 的秩，其贡献
记为 \(a_i x_i+d_i y_i\)。若秩为 2，保留该方程出现与该贡献。若秩
为 0，省略该块。剩余情形秩为 1，系数恰有
\[
 (a_i,d_i)=(v,0),\quad (0,v),\quad (v,v),
 \qquad v\ne0
\]
三种形式之一。这些形式分别只保留位置一、二、三处的变量，贡献为
实际比特 \(z_i\) 上的 \(v z_i\)。对第三种形式，
\[
 v(x_i+y_i)=v z_i+b_{e_i}v t_0,
 \qquad z_i=b_{e_i}t_0+x_i+y_i,
\]
故该替换同时给齐次截距加上 \(b_{e_i}v\)。在单例块中，恰在其斜率
非零时保留其变量 ID 与原贡献；否则省略该块。

保留的块构成 \((Q,P)\) 的\emph{支撑}。它们保持自己在 \([k]\) 中的
指标、方程/变量标签与 ID。保留的方程保持其有序的前两个坐标，保留
的单例用其实际比特坐标。于是不同块中的重复 ID 仍有分离的局部坐标。
支撑可以为空；无论何种情形齐次坐标都在。

设 \(J_3\) 是被约化到第三比特的方程块的集合，把这种块中的非零
斜率记为 \(v_i\)。支撑上的\emph{约化表} \(\widetilde P\) 具有上述
保留块贡献，及齐次截距
\begin{equation}
 \widetilde c=c+\sum_{i\in J_3} b_{e_i}v_i.
 \label{eq:reduced-table-intercept}
\end{equation}
为核对完整表，设 \(L_{\rm supp}\) 为支撑的齐次答案空间，
\(\rho_{Q,P}:L_Q\to L_{\rm supp}\) 保留 \(t_0\) 与指定的块坐标。
它在保留的方程上、以及在 \(Q\) 中已有的保留单例块上是恒等；在
被约化的方程上使用 Equation~\eqref{eq:outer-projection} 的相应
实际比特坐标；并丢弃被省略的块。在每个块中独立选坐标可知
\(\rho_{Q,P}\) 满射。各块恒等式与
Equation~\eqref{eq:reduced-table-intercept} 给出
\[
 P=\widetilde P\rho_{Q,P}.
\]
满射性使 \(\widetilde P\) 成为具有此分解的唯一线性映射。我们把
\((Q,P)\) 的\emph{键}定义为其支撑连同这整张约化表。于是两张表有
相同键，当且仅当它们的保留坐标一致（相同的位置、标签与 ID），且
其诱导仿射映射在这些坐标上一致（包括截距）。

接下来核对与每个容许投影 \(\bar\pi:L_U\to L_O\) 的相容性。设
\(P'\in\Hom(L_O,\mathcal V)\)。\(O\) 中保留的方程在 \(P'\) 及其
拉回 \(P'\bar\pi\) 中有相同的贡献，故其约化相同。斜率为零的单例在
两侧都消失。斜率非零 \(v\) 的单例拉回为上述三种秩一模式之一，约化
到同一块位置上的同一变量 ID。对第三位置的单例，拉回给截距加上
\(b_{e_i}v\)，而其约化再加一次 \(b_{e_i}v\)。两项在 \(\F_2\) 上
相消。其余截距修正都来自保留的方程块，且在两侧一致。因此支撑与
完整约化表一致：\((U,P'\bar\pi)\) 与 \((O,P')\) 有相同的键。

把键的全域定义为方程问题 \(U\in\mathcal E^k\) 上所有表的键。刚证
的相容性表明：每个容许投影问题上的表的键都属于这个全域——从任何
相容的 \(U\) 拉回即可。

现在应用 Bellare、Goldreich 与 Sudan 的折叠原理
~\cite[Section~3.3]{BGS98}：一个代表元上的取值决定其平移上的取值。
给键加常数 \(h\in K\) 指给其约化截距加 \(h\)，支撑与斜率不动。
这由把呈现表 \(P\) 换成 \(P+h e_0^\top\) 实现，故保持键全域。
由于键记录了截距，该作用是自由的。我们的置换实例的顶点就是该作用
的一条轨道，由其唯一的、约化截距 \(K\) 分量为零的键表示。

对方程问题或容许投影问题 \(Q\) 上的表 \(P\)，令 \(c(Q,P)\in K\)
为其键的约化截距的 \(K\) 分量，\(v(Q,P)\) 为去掉该分量后的键所
代表的轨道顶点。顶点的染色 \(\lambda\) 定义恢复表答案
\[
 F_Q(P)=\lambda(v(Q,P))+c(Q,P).
\]
加常数只改变被去掉的分量，故
\begin{equation}
 F_Q(P+h e_0^\top)=F_Q(P)+h
 \qquad(h\in K).
 \label{eq:table-folding}
\end{equation}
同键的表有相同的代表与被去分量。因此拉回相容性给出：对每个容许的
\(U,O,\bar\pi\) 与 \(O\) 上的每张表 \(P'\)，
\begin{equation}
 F_U(P'\bar\pi)=F_O(P').
 \label{eq:key-sharing}
\end{equation}
于是投影问题使用由同一顶点染色决定的答案，无需额外的一致性测试。

\subsection{潜在矩阵测试与完备性}

置换实例记作 \(\mathcal G\)，独立采样
\[
 U\sim\omega^{\otimes k},\qquad
 P\text{ 在 }\Hom(L_U,\mathcal V)\text{ 中均匀},\qquad
 a\sim\mu,\qquad l\text{ 在 }\F_2^m\text{ 中均匀},
\]
并检验
\begin{equation}
 F_U(P)=F_U(P+a l^\top).
 \label{eq:matrix-test}
\end{equation}
此处 \(a l^\top\) 表示映射 \(z\mapsto a(l^\top z)\)。每个样本指定
两个键轨道顶点之间的一条边。若其恢复偏移为 \(c_0,c_1\)，则其约束
是固定字母表 \(K\) 上的置换
\[
 \lambda(v_1)=\lambda(v_0)+c_0+c_1.
\]
边出现以采样概率为权重。此阶段允许自环与重复边。涉及 \(O\) 的
稀疏投影实验用于分析这些同样的顶点标记；它不是 \(\mathcal G\) 中
的额外测试。

\begin{lemma}[外层矩阵测试的完备性]
\label{lem:outer-completeness}
若 \(\mathcal E\) 有值至少 \(1-\xi\) 的赋值，且潜在映射满足
\[
 \Pr_{x\text{ 在 }\mathcal V\text{ 中均匀},\ a\sim\mu}
       [C(x+a)\ne C(x)]\le p_*,
\]
则
\[
 \val(\mathcal G)\ge 1-k\xi-\frac{p_*}{2}.
\]
\end{lemma}

\begin{proof}
固定近满足的全局比特赋值。给每个约化支撑如下合法答案：单例用其
全局比特；保留的方程用其前两个全局比特，并置第三个为满足该方程的
唯一值。这些选择只依赖支撑（包括其出现 ID），因此对每个精确键
独立于其来源地定义了一个答案。

在轨道代表元处，用该答案对其约化表求值并应用 \(C\)。这定义一个
顶点染色。\(C\) 在 \(K\) 下的等变性蕴含：其恢复答案是对实际约化表
求值再作用 \(C\)，与 Equation~\eqref{eq:table-folding} 一致。

以至少 \(1-k\xi\) 的概率，\(U\) 中的每条方程都被全局赋值满足。
在这样的元组上，\emph{每张}表使用的支撑答案都是同一个完整答案
\(z_U\in L_U\)（首坐标为 1）的投影。于是
Equation~\eqref{eq:matrix-test} 的两侧是
\[
 C(Pz_U)\quad\text{与}\quad
 C(Pz_U+a(l^\top z_U)).
\]
由于 \(z_U\ne0\)，向量 \(Pz_U\) 在 \(\mathcal V\) 中均匀，与
\(a,l\) 独立，而 \(l^\top z_U\) 是独立的公平比特。因此测试在这些
元组上以至多 \(p_*/2\) 的概率失败。加上元组含未满足方程的概率即
得该界。
\end{proof}

\section{矩阵抽取与傅里叶解码}
\label{sec:matrix}

我们把矩阵测试的高接受率转化为两个策略在稀疏投影问题 $U,O$ 上的
一致性。每个策略只用其自身问题与有界数目的矩阵建议行。所得一致性
概率为正，且与元组长度 $k$ 及奇偶实例无关。表函数只通过其折叠与
共享恒等式以及接受界进入证明，故该证明适用于任何具备这些性质的
族；每个顶点染色的恢复函数都是例子。精确假设在固定参数并定义建议
实验之后陈述。

\subsection{逆短码定理}

结构性输入描述的是在均匀秩一扰动下经常保持取值的函数。这样的函数
在「固定有界数目的行、列线性组合」所得的集合上与一个仿射求值
一致。

记 $\operatorname{Mat}_{\ell,m}=\Hom(\F_2^m,\F_2^\ell)$。
\emph{行与列切片}是由方程
\begin{equation}
 \mathcal S=
 \{M:d_i^\top M=s_i^\top\ (1\le i\le n_r),\quad
          Mq_j=t_j\ (1\le j\le n_c)\}
 \label{eq:shortcode-slice}
\end{equation}
规定的非空集，其中 $d_i\in\F_2^\ell$、$s_i\in\F_2^m$、
$q_j\in\F_2^m$、$t_j\in\F_2^\ell$。计数 $n_r,n_c$ 不必是其方程组
的秩。

\begin{theorem}[二次短码的相等可靠性]
\label{thm:shortcode}
对每个 $\eta\in(0,1)$，存在只依赖 $\eta$ 的 $\alpha\in(0,1]$、
整数 $r_{\rm s}\ge1$ 与整数 $\ell_0$，具有以下性质。对每个
$\ell\ge\ell_0$ 有 $m_0=m_0(\eta,\ell)$，使得对每个 $m\ge m_0$ 与
每个函数 $f:\operatorname{Mat}_{\ell,m}\to\F_2^\ell$，若
\begin{equation}
 \Pr_{M,a,l}\bigl[f(M)=f(M+a l^\top)\bigr]\ge\eta,
 \label{eq:shortcode-acceptance}
\end{equation}
则存在如 Equation~\eqref{eq:shortcode-slice} 的切片 $\mathcal S$
（$n_r,n_c\le r_{\rm s}$）与某个 $z\in\F_2^m$、$u\in\F_2^\ell$，使
\begin{equation}
 \Pr_{M\mid\mathcal S}[f(M)=Mz+u]\ge\alpha.
 \label{eq:shortcode-conclusion}
\end{equation}
此处 $M$、$a\in\F_2^\ell$、$l\in\F_2^m$ 独立均匀，而
$M\mid\mathcal S$ 在切片上均匀。对 $f$ 不施加折叠条件。
\end{theorem}

\begin{proof}
我们从 Khot--Minzer--Safra 的 Grassmann 扩张定理
\cite[Theorem~1.12 in the 2018 revision]{KMS23} 导出此推论，使用
Barak--Kothari--Steurer 的矩阵图卡
\cite[Definition~3.2 and Lemmas~3.3--3.7]{BKS18}。

环境二元空间 $E$ 上的 Grassmann 图以 $\ell$ 维子空间为顶点；两个
顶点相邻当其交的维数为 $\ell-1$。非空顶点集的\emph{保留率}是：
集合中的均匀顶点有一个均匀邻居仍在集合中的概率。所引定理说：保留
率至少固定 $\zeta>0$ 的集合在某个区间
\[
 \{L:A_0\subseteq L\subseteq B_0\},
 \qquad \dim A_0+\operatorname{codim}_E B_0\le r
\]
中有正的相对密度。密度界与 $r$ 只依赖 $\zeta$。定理对所有足够大的
$\ell$、进而所有足够大的 $\dim E$ 成立。

相等接受是各非空纤维 $f^{-1}(y)$ 在矩阵步下保留率的加权平均（权重
为均匀质量）。故某纤维 $S$ 的保留率至少 $\eta$。两个扰动因子之一
为零的概率是
\[
 p_0=2^{-\ell}+2^{-m}-2^{-\ell-m}.
\]
把维数阈值选得使 $p_0\le\eta/2$。在两个因子都非零的条件下，$S$ 的
保留率至少 $(\eta-p_0)/(1-p_0)\ge\eta/2$。

对每个矩阵令
\[
 L_M=\{(x,M^\top x):x\in\F_2^\ell\}
 \subseteq E:=\F_2^\ell\oplus\F_2^m.
\]
映射 $M\mapsto L_M$ 是到图卡 $\mathcal C$（到第一加项投影可逆的
子空间）的双射。由于
\[
 \dim(L_M\cap L_N)=\ell-\rank(M-N),
\]
图卡内的邻居恰对应非零的秩一矩阵差。每个图卡顶点有
$(2^\ell-1)(2^m-1)$ 个这样的邻居，而其完整 Grassmann 度为
$(2^\ell-1)(2^{m+1}-2)$：先在顶点中选超平面，再在该超平面的商中选
一维延拓（非顶点本身）。故恰有一半邻居在图卡内。在 $\F_2$ 上，
每个非零秩一矩阵都有唯一分解 $a l^\top$（因子非零）。因此提升集
$\widetilde S=\{L_M:M\in S\}$ 的 Grassmann 保留率至少 $\eta/4$。

以 $\zeta=\eta/4$ 应用扩张定理，得某个区间上的密度 $\alpha$，其中
$\dim A_0+\operatorname{codim}_E B_0\le r_{\rm s}$。该区间与
$\mathcal C$ 的交只会提高 $\widetilde S$ 的相对密度（因
$\widetilde S\subseteq\mathcal C$），且交非空。$A_0$ 的基
$(d_i,s_i)$ 用方程 $d_i^\top M=s_i^\top$ 刻画
$A_0\subseteq L_M$。对 $E$ 的标准点积，$B_0^\perp$ 的基 $(t_j,q_j)$
用 $Mq_j=t_j$ 刻画 $L_M\subseteq B_0$。二者的联立解集即为所需
形式的切片，两类方程各至多 $r_{\rm s}$ 个。在该切片的至少 $\alpha$
比例上，$f$ 等于定义 $S$ 的那个值。取 $z=0$、$u$ 等于该值即得
结论。

常数 $\alpha,r_{\rm s}$ 的选取只用 $\eta$。由于
$\dim E=\ell+m$，固定 $\ell$ 后剩余的维数阈值被吸收进
$m_0(\eta,\ell)$。未用任何折叠条件。纤维论证是
\cite[Lemma~2.3]{BKS18} 的相等推论；该文的编号引用均指其 arXiv
版本 1。
\end{proof}

Theorem~\ref{thm:shortcode} 中的常数与维数阈值可以有效地选取，
当 $\eta$ 有理时 $\alpha$ 亦有理。这由 \cite[Theorem~2.13 and
Section~2.4]{KMS23} 的定量界与 \cite[Lemmas~2.7--2.8]{KMS23} 的
比较估计得到；定理编号同样指 2018 修订版。

现在固定下文所需参数。从 Lemma~\ref{lem:latent} 选 $p_*>0$ 与相应
的 $r_*$。令
\begin{equation}
 \eta=2^{-r_*}/16,
 \qquad
 2^{-(\ell-r_{\rm s})}<\alpha/8,
 \qquad
 2^{-\ell}<\eta/4,
 \label{eq:matrix-parameters}
\end{equation}
其中 $\alpha,r_{\rm s}$ 先由 Theorem~\ref{thm:shortcode} 给出，
然后选足够大的 $\ell$ 满足该定理、潜在引理以及所示不等式。在此
$\ell$ 值处取潜在数据 $(\mathcal V,K,C,\mu)$，并固定线性分裂
$\mathcal V=K\oplus W$。所有这些对象都在 $k$ 增长之前固定。
特别地，$\ell>r_{\rm s}$。我们把表写成 $P=(M,T)$，其中
$M:L_U\to K$、$T:L_U\to W$。

\subsection{建议实验与解码陈述}

我们先指定两个策略的局部输入，其答案一致性将与外层博弈间隙矛盾。
该实验对每个正整数 $k$ 与每个 $0<\beta<1$ 定义。设 $\mathcal Q_k$
由所有方程问题 $U\in\mathcal E^k$ 以及由它们经
Section~\ref{sec:game} 的坐标投影得到的容许投影问题 $O$ 组成。
问题 $U,O$ 与投影 $\bar\pi:L_U\to L_O$ 的律见 Section~\ref{sec:game}：
在每个块中独立地，方程以概率 $1-\beta$ 被保留，均匀选取的变量位置
以概率 $\beta$ 被保留。

\begin{definition}[投影建议实验]
\label{def:projected-advice}
采样 $U,O,\pi$，再采样独立的均匀映射 $A:K\to\F_2^{r_{\rm s}}$ 与
独立均匀映射
\[
 T_O:L_O\to W,\qquad M_O:L_O\to K.
\]
令 $S_O=AM_O$，并把以下输入发给两个证明者：
\begin{equation}
 \begin{aligned}
  &\text{$U$ 证明者：}&& (U,A,T_O\bar\pi,S_O\bar\pi),\\
  &\text{$O$ 证明者：}&& (O,A,T_O,S_O).
 \end{aligned}
 \label{eq:projected-advice}
\end{equation}
所给建议之外的映射 $M_O$ 被隐藏。每个证明者返回一个实际答案，即
其齐次答案空间中首比特为 1 的向量。当其答案 $a_U,a_O$ 满足
$\bar\pi a_U=a_O$ 时成功。
\end{definition}

$U$ 消息包含映射 $(T_U,S_U)=(T_O\bar\pi,S_O\bar\pi)$，由在 $O$ 上
采样的映射补零得到。我们将先在另一种律下找到有用的建议：$U$ 上的
无限制均匀表 $(M,T)$，取 $S_U=AM$。随后在把投影选择平均之后，把
完整的补零表 $(M_O\bar\pi,T_O\bar\pi)$ 与该无限制表比较。这一比较
是从抽取论证通向实验的桥梁。

\begin{proposition}[矩阵解码]
\label{prop:matrix-decoding}
按 Equation~\eqref{eq:matrix-parameters} 指定的顺序选取潜在与短码
参数，并在 $k$ 之前固定 $\mathcal V=K\oplus W$。存在只依赖这些
固定参数的常数 $\gamma>0$ 与整数 $k_0\ge8$，具有以下性质。

对每个立方数 $k\ge k_0$ 与每个奇偶实例，设对每个 $Q\in\mathcal Q_k$
给定了函数
\[
 F_Q:\Hom(L_Q,\mathcal V)\longrightarrow K.
\]
假设对每个这样的 $Q$、每个 $P\in\Hom(L_Q,\mathcal V)$ 与每个
$c\in K$ 满足折叠，并对每个容许投影 $\bar\pi:L_U\to L_O$ 与每个
$P'\in\Hom(L_O,\mathcal V)$ 满足共享：
\begin{equation}
 \begin{aligned}
 F_Q(P+c e_0^\top)&=F_Q(P)+c,\\
 F_U(P'\bar\pi)&=F_O(P').
 \end{aligned}
 \label{eq:matrix-family-identities}
\end{equation}
若该族以至少 $0.99$ 的概率通过 Equation~\eqref{eq:matrix-test} 的
表测试，则取 $\beta=k^{-2/3}$ 的投影建议实验具有局部策略，其完整
答案一致的概率至少 $\gamma$。
\end{proposition}

$\mathcal G$ 顶点的每个染色都给出这样的族：其恢复函数由
Equations~\eqref{eq:table-folding} 与~\eqref{eq:key-sharing} 满足
Equation~\eqref{eq:matrix-family-identities}。抽象表述记录的是：
解码论证不使用顶点构造的任何其他性质。在定义两个应答规则时该族
已固定，并可以内置于规则之中。

接下来两个小节对任意正整数 $k$ 与满足
Equation~\eqref{eq:matrix-family-identities} 的固定族成立。它们
使用 $U$ 上的无限制均匀表，尚未限制 $\beta$。在把该律与投影实验
比较时才特化 $k$ 与 $\beta$。

\subsection{比较两种矩阵扰动}

均匀矩阵扰动具有 \cite[Lemma~2.10 in the 2018 revision]{KMS23}
计算的按秩索引的傅里叶谱。我们把它与潜在 gadget 提供的噪声比较。

\begin{lemma}[谱比较]
\label{lem:matrix-spectral}
若表族以至少 $0.99$ 的概率通过 Equation~\eqref{eq:matrix-test}，则
\begin{equation}
 \E_{U,T}\Pr_{M,a,l}
   [F_U(M,T)=F_U(M+a l^\top,T)]
 \ge 0.9\,2^{-r_*},
 \label{eq:surrogate-acceptance}
\end{equation}
其中 $a$ 在 $K$ 中均匀，其余矩阵与向量样本在各自空间中均匀。因此
至少 $10\eta$ 的对 $(U,T)$ 的内层概率至少 $4\eta$。
\end{lemma}

\begin{proof}
对 $\sigma\in K^*$ 定义符号函数
$f_{U,\sigma}(P)=(-1)^{\sigma(F_U(P))}$。矩阵域上的频率是泛函元组
$g=(g_i)_{i=0}^{m-1}\in\mathcal V^*$；其在 $P$ 处的特征为
$(-1)^{\sum_i g_i(Pe_i)}$。所有傅里叶系数用均匀概率测度。平移
$a l^\top$ 把该特征乘以 $(-1)^{\sum_i l_i g_i(a)}$。对均匀 $l$ 取
平均得特征值
\[
 \Lambda_g=\Pr_{a\sim\mu}[g_i(a)=0\text{ 对每个 }i].
\]
它在 $[0,1]$ 中。由 Lemma~\ref{lem:latent}，只要
\[
 \dim\Span\{g_i|_K:0\le i<m\}\ge r_*
\]
它就至多为 $7/8$。事实上 $g_i$ 的张成是 $\mathcal V^*$ 的子空间，
引理适用于它。

用一个均匀输出特征展开相等，再用傅里叶正交性，把测试接受率表示为
\[
 \E_{U,\sigma}\sum_g
      \abs{\widehat f_{U,\sigma}(g)}^2\Lambda_g.
\]
Parseval 给出每个符号函数的总质量平方为 1。若 $B$ 是限制秩至少
$r_*$ 上的平均质量平方，则 $0.99\le1-B/8$，故 $B\le0.08$。特别
地，至少 $0.9$ 的平均质量落在低于 $r_*$ 的秩上。

当 $a$ 改为在 $K$ 中均匀时，$g$ 处的特征值为
$2^{-\dim\Span\{g_i|_K\}}$。上述质量因此给出
Equation~\eqref{eq:surrogate-acceptance} 的下界。该扰动固定
$T$，故其中的表达式恰是固定 $(U,T)$ 处普通短码相等测试的平均。
其下界为 $14.4\eta$。设 $q$ 为内层接受至少 $4\eta$ 的概率，则该
平均至多 $4\eta+q$。于是 $q\ge10.4\eta\ge10\eta$。
\end{proof}

\subsection{获得许多有用的行纤维}

逆定理在每个合格 $(U,T)$ 纤维上给出一个结构化切片。$U$ 证明者
关于 $K$ 值矩阵的信息是对 $(A,AM)$。我们需要：这些被观察的行以正
概率决定一个有用的切片。我们还要识别该切片提供的合法答案短集。
下面的擦除论证证明概率界，折叠随后使答案集非空。

对这一抽取步骤，照旧采样 $U$，并取独立均匀映射 $T:L_U\to W$、
$A:K\to\F_2^{r_{\rm s}}$、$M:L_U\to K$。对固定的 $U,T,A$ 与映射
$S_0:L_U\to\F_2^{r_{\rm s}}$，集合 $\{M:AM=S_0\}$ 是其\emph{行
纤维}。称行建议 $(U,T,A,S_0)$ 为\emph{好的}，若存在至多 $r_{\rm s}$
个列规定与某个目标 $Mz+u$，使列规定与 $AM=S_0$ 的联立解集非空，
且在该集合的均匀测度下 $F_U(M,T)=Mz+u$ 以至少 $\alpha/2$ 的概率
成立。

这是行纤维上的解析条件。目标向量 $z$ 起初不必是实际答案。固定的
列值允许 $z$ 在其列方向张成的模下变化，$u$ 相应变化。因此切片
自然提供可能目标向量的一个陪集。其合法部分将是 $U$ 证明者采样的
答案集。

\begin{lemma}[好行建议及其短答案集]
\label{lem:good-row-advice}
对所有足够大的 $k$，矩阵测试接受率至少 $0.99$ 的表族满足
\begin{equation}
 \Pr_{U,T,A,M}[(U,T,A,AM)\text{ 是好的}]
 \ge g_0,
 \qquad
 g_0=10\eta^2\,2^{-\ell r_{\rm s}}>0.
 \label{eq:good-advice-mass}
\end{equation}
该下界与所需的 $k$ 之大与奇偶实例及表族无关。

对每个正整数 $k$ 与每个满足
Equation~\eqref{eq:matrix-family-identities} 的族，可以确定性地为
每个好行建议选取至多 $r_{\rm s}$ 个列规定 $Mq_j=t_j$
（$q_j\in L_U$、$t_j\in K$）与目标 $Mz+u$（$z\in L_U$、$u\in K$、
$e_0^\top z=1$），满足好的定义。记 $Z=\Span\{q_j\}$，所选见证
决定答案集
\begin{equation}
 \mathcal A=(z+Z)\cap\{v\in L_U:e_0^\top v=1\},
 \qquad 1\le |\mathcal A|\le 2^{r_{\rm s}}.
 \label{eq:row-answer-set}
\end{equation}
见证与 $\mathcal A$ 只依赖 $(U,T,A,S_0)$ 与固定的表族。切片及其
目标保持好的定义中的一致性保证。
\end{lemma}

\begin{proof}
固定替代接受率至少 $4\eta$ 的 $(U,T)$，记 $f(M)=F_U(M,T)$。令
$B$ 为矩阵域中所有好行纤维（对一切可能的 $A,S_0$）的并。设其均匀
测度小于 $\eta$。在 $B$ 的表项上修改 $f$，把每一项换成 $K$ 的独立
均匀元素。每个矩阵只被替换一次，即使多条好行纤维在该处重叠。短码
步的两个端点都是均匀矩阵。因此每次这样的修改使接受率至多改变
$2\Pr[B]$，修改后的函数接受率至少 $2\eta$。

我们证明：可以选取随机替换，使其作为与每个目标 $Mz+u$ 匹配的
贡献，在每个行与列的联立切片中少于 $\alpha/4$ 的规模。令
$r=r_{\rm s}$。至多 $r$ 个行规定与 $r$ 个列规定的非空切片，其
齐次空间为
\[
 \Hom\bigl(\F_2^m/\Span\{q_j\},\ \bigcap_i\ker d_i^\top\bigr).
\]
其规模 $n_{\mathcal S}$ 至少为 $2^{(\ell-r)(m-r)}$。用零方程补齐
规定列表后，「切片–目标」描述的个数至多 $2^{(2r+1)(\ell+m)}$。

对固定的描述，只统计随机化表项处的匹配，对其他表项赋以确定的零
贡献。它们是 $[0,1]$ 中的独立变量，总均值至多
$2^{-\ell}n_{\mathcal S}<\alpha n_{\mathcal S}/8$。Hoeffding 不等
式 \cite[Theorem~2, equation~(2.6)]{Hoeffding1963} 把出现至少
$\alpha n_{\mathcal S}/4$ 个这种匹配的概率界定在
$\exp(-\alpha^2 n_{\mathcal S}/32)$ 以内。对所有描述的并界因此
至多
\[
 \exp\left((2r+1)(\ell+m)\log2
       -\frac{\alpha^2}{32}2^{(\ell-r)(m-r)}\right),
\]
当 $m=1+2k$ 增大时趋于零。固定具有所述性质的替换。

考虑任何容许类型的切片。若其行纤维是好的，则该纤维的每个表项都
属于 $B$，其修改后的一致性小于 $\alpha/4$。若行纤维不好，则原
切片上与每个目标的一致性小于 $\alpha/2$。未修改的匹配是原匹配的
子集，而随机化匹配再加少于 $\alpha/4$。修改后的一致性于是小于
$3\alpha/4$。用零补齐较短的行系统表明：这两种情形覆盖
Theorem~\ref{thm:shortcode} 中的每个切片。该定理随即与修改后
函数至少 $2\eta$ 的接受率矛盾。其适用性不要求修改保持折叠。

因此 $\Pr[B]\ge\eta$。对每个 $M\in B$，至少有一个矩阵 $A$ 使某个
好行纤维含 $M$。均匀选取的 $A$ 命中该特定矩阵的概率为
$2^{-\ell r_{\rm s}}$。故在这个固定的 $(U,T)$ 处，好数据的概率
至少 $\eta2^{-\ell r_{\rm s}}$。Lemma~\ref{lem:matrix-spectral}
提供至少 $10\eta$ 个这样的对 $(U,T)$，证明了
Equation~\eqref{eq:good-advice-mass}。所有计数与集中界只依赖固定
参数与 $m$，如所需。

其余断言不使用高测试接受率或 $k$ 之大。固定任意正整数 $k$、满足
Equation~\eqref{eq:matrix-family-identities} 的族与任意好行建议。
固定一个见证，把其列规定写成 $Mq_j=t_j$，令 $Z=\Span\{q_j\}$。
切片非空使规定的值相容：它们确定线性映射 $t:Z\to K$，在整个切片
上 $t(w)=Mw$。把 $z$ 换成 $z+w$、$u$ 换成 $u+t(w)$（$w\in Z$）
因此保持目标。只需证明陪集 $z+Z$ 含首比特为 1 的向量。

否则 $z$ 与 $Z$ 中每个向量首比特都为零。令 $H=\ker A$。对每个
$h\in H$，平移 $M\mapsto M+h e_0^\top$ 保持行纤维、列规定与目标。
由折叠它把 $F_U(M,T)$ 平移 $h$。该作用自由，且其每条 $H$-轨道中
至多一个矩阵满足目标等式。一致性至多
\[
 |H|^{-1}=2^{-\ell+\rank A}
 \le2^{-(\ell-r_{\rm s})}<\alpha/8,
\]
与好性矛盾。选 $w\in Z$ 使 $z+w$ 首比特为 1，并对 $z,u$ 作相应
改变。它们保持原一致性比例。用有限可能性上的固定顺序选一个见证
与这样的代表。该选择只用行建议与固定表族。其向量 $z$ 落在
$\mathcal A$ 中，故集合非空；由 $\dim Z\le r_{\rm s}$，其规模
至多 $2^{r_{\rm s}}$。
\end{proof}

\subsection{转移所选切片}

与稀疏投影问题的比较遵循 Khot--Safra 的光滑外层博弈方法
~\cite[Section~3 and Lemma~3.1]{KS13}，共享线性建议如
Khot--Minzer--Safra~\cite[Sections~3.3--3.4]{KMS17}。
Lemma~\ref{lem:good-row-advice} 给出正概率的行建议：在 $U$ 上的
无限制采样下决定一个切片与一个短的合法答案集。在实际实验中，映射
在 $O$ 上采样并补零到 $U$。我们把完整的补零表（含仿射截距）与
无限制表比较（投影选择被平均掉），随后把所选切片转移到 $O$。

从现在起取大于 1 的立方数 $k$，并在 Definition~\ref{def:projected-advice}
中令 $\beta=k^{-2/3}$。

取 $\beta=k^{-2/3}$ 有两个目的：$k\beta^2\to0$ 使补零表律在
Lemma~\ref{lem:projected-tv} 中逼近无限制律，而 $k\beta\to\infty$
使单例坐标的期望数目发散。一旦建议空间固定，每个单例都有一个一致
为正的概率：每条可见建议行在其变量比特上的系数为零。
Lemma~\ref{lem:clean} 用这些坐标界定一致性。

\begin{lemma}[补零矩阵渐近均匀]
\label{lem:projected-tv}
在每个 $U$ 的条件下，分布
\[
 (M,T)=(M_O\bar\pi,T_O\bar\pi)
\]
与 $\Hom(L_U,K)\times\Hom(L_U,W)$ 上的无限制均匀分布的全变差距离
为 $o(1)$（$k$ 沿立方数趋于无穷）。该界对 $U$ 与奇偶实例一致。
投影选择在此比较中被平均掉。
\end{lemma}

\begin{proof}
把两个分量合成一张 $\mathcal V$ 值表，令 $n=|\mathcal V|$。在未
投影的方程块上，两个斜率是 $\mathcal V$ 中独立的均匀元素。在单例
块上，补零后的斜率对是
\[
 (a,0),\qquad (0,a),\qquad (a,a),
 \qquad a\text{ 在 }\mathcal V\text{ 中均匀}
\]
的等权混合。称这两个对律为 $P_0$ 与 $Q_0$。似然比 $dQ_0/dP_0$ 在
$(0,0)$ 处为 $n$，在三条所示直线上的其余 $3(n-1)$ 个点处为
$n/3$，其他处为零。于是
\[
 \chi^2(Q_0\Vert P_0)
 :=\E_{P_0}\left[\left(\frac{dQ_0}{dP_0}-1\right)^2\right]
 =\frac{n-1}{3}.
\]
因此单块混合 $(1-\beta)P_0+\beta Q_0$ 相对 $P_0$ 的卡方散度为
$\beta^2(n-1)/3$。各块独立采样（即使其问题 ID 重合）。把似然比
相乘给出散度
\begin{equation}
 \left(1+\frac{n-1}{3}\beta^2\right)^k-1=o(1),
 \label{eq:projected-chi-square}
\end{equation}
因为 $n$ 固定而 $k\beta^2=k^{-1/3}$ 趋于零。全变差至多该表达式
平方根的一半。

对方程 $e$ 投影到第三比特的情形，补零截距获得项 $b_e a$（$a$ 为
该单例的斜率）。原截距独立且在 $\mathcal V$ 中均匀。加上这些项
之和后，即使在所有投影选择与斜率都已给出之后，它仍均匀独立。前面
的斜率估计因此证明了整张表的断言。斜率律与
Equation~\eqref{eq:projected-chi-square} 的界都不依赖方程 ID 或其
右端。
\end{proof}

暂时固定好的 $U$-建议及其来自 Lemma~\ref{lem:good-row-advice} 的
所选见证与答案集 $\mathcal A$。对容许投影令 $z'=\bar\pi z$ 与
$Z'=\bar\pi Z$。由于 $\bar\pi$ 保持首坐标，
\begin{equation}
 \bar\pi\mathcal A
 =(z'+Z')\cap\{v\in L_O:e_0^\top v=1\}.
 \label{eq:projected-answer-set}
\end{equation}
事实上，右端陪集中的每个向量在 $z+Z$ 中有原像，且原像有相同的
首坐标。我们将证明 $O$ 证明者的私有傅里叶样本到达这个投影后的
答案集。

见证切片并不自动在投影下存活。例如列条件 $Mq=t$ 变成
$M_O\bar\pi q=t$；若 $\bar\pi q=0$ 而 $t\ne0$ 则不相容。更一般
地，投影后的列方向可能出现新的依赖。下一个论证使用完整表律的
比较，找出正概率的样本：其中所选切片保持非空并保住足够的一致性。

为这一分析，记 $D=(U,O,\pi,A,T_O,S_O)$ 为同时决定两个证明者消息
与投影的联合数据。这些数据在约束 $AM_O=S_O$ 之外不指定 $M_O$。

\begin{lemma}[所选见证的转移]
\label{lem:projected-witness}
设矩阵测试接受率至少 $0.99$。对所有足够大的立方数 $k$，
Definition~\ref{def:projected-advice} 中概率至少为
\begin{equation}
 \gamma_1=\frac{\alpha g_0}{8}\,2^{-\ell r_{\rm s}}>0
 \label{eq:projected-witness-mass}
\end{equation}
的联合数据 $D$ 具有如下性质：补零后的 $U$-建议是好的，带所选
答案集 $\mathcal A$ 及其见证。记 $z'=\bar\pi z$，该见证转移到
非空切片
\[
 AM_O=S_O,\qquad M_O\bar\pi q_j=t_j,
\]
且在其上
\begin{equation}
 \Pr\bigl[F_O(M_O,T_O)=M_Oz'+u\mid
           AM_O=S_O,\ M_O\bar\pi q_j=t_j\text{ 对所有 }j\bigr]
 \ge\alpha/4.
 \label{eq:transferred-witness}
\end{equation}
此式中的概率取遍指定切片上的均匀矩阵。
\end{lemma}

\begin{proof}
先用 $U,A,M,T$ 的无限制采样，令 $S_0=AM$。设 $\mathcal J$ 为事件
「$(U,T,A,S_0)$ 是好的且 $M$ 满足其所选列规定」，$\mathcal M$
为与所选目标相等（在坏数据之外声明该事件为假）。在好行纤维内，
非空列切片施加至多 $\ell r_{\rm s}$ 个额外的标量线性方程。故
Lemma~\ref{lem:good-row-advice} 给出
\[
 \Pr[\mathcal J]\ge g_0\,2^{-\ell r_{\rm s}},
 \qquad
 \Pr[\mathcal J\cap\mathcal M]\ge(\alpha/2)\Pr[\mathcal J].
\]
从而
\begin{equation}
 \Pr[\mathcal J\cap\mathcal M]
       -(\alpha/4)\Pr[\mathcal J]
 \ge\frac{\alpha g_0}{4}\,2^{-\ell r_{\rm s}}.
 \label{eq:witness-positive-difference}
\end{equation}

该表达式中的两个事件只依赖 $(U,A,M,T)$。把独立的 $A$ 加入
Lemma~\ref{lem:projected-tv} 不改变其变差界。因此可以用补零实验
中的值比较两个事件概率。若全变差距离为 $v_k=o(1)$，
Equation~\eqref{eq:witness-positive-difference} 中的差至多改变
$(1+\alpha/4)v_k$。对所有足够大的 $k$，补零后的差至少 $\gamma_1$。

现在对 $D$ 条件化。剩余的 $M_O$ 在行纤维 $AM_O=S_O$ 上均匀。
补零后的 $U$-建议（从而其所选见证与答案集）固定。在好数据上，令
$p_D$ 为投影后列切片的条件概率；当 $p_D>0$ 时，令 $q_D$ 为其
条件相等比例。在坏数据上置 $p_D=0$，当 $p_D=0$ 时置 $q_D=0$。
共享给出 $F_U(M_O\bar\pi,T_O\bar\pi)=F_O(M_O,T_O)$，故补零后的差
恰为
\[
 \E_D\bigl[p_D(q_D-\alpha/4)\bigr]\ge\gamma_1.
\]
被积项在「好性、$p_D>0$、$q_D\ge\alpha/4$」组成的事件之外非正，
且处处至多 1。因此该事件概率至少 $\gamma_1$，且恰具有所述性质。
律比较只用于仅依赖 $(U,A,M,T)$ 的事件；此处对 $D$ 的条件化只是
指出它们的正差在投影实验中的位置。
\end{proof}

\subsection{解码行建议}

从傅里叶质量抽取外层策略的方法遵循 H\aa stad
~\cite[Section~1.1]{Hastad01} 与 Khot--Safra 的傅里叶列表方法
~\cite[Section~1.1]{KS13}。此处投影后的证明者在自己的行纤维上采样
傅里叶质量平方。转移来的切片将表明：该私有样本以正概率到达投影后
的答案集 $\bar\pi\mathcal A$。

现在由各自的消息定义两个策略。联合数据与转移来的见证只用于分析
其一致性；$O$ 证明者不使用 $U$ 证明者的见证或答案集。

$U$ 证明者用其可见的行建议执行 Lemma~\ref{lem:good-row-advice}
中的确定性选择。当建议是好的时，它从 $\mathcal A$ 中均匀采样
$a_U$。这是至多 $2^{r_{\rm s}}$ 个实际答案的非空集。在其他输入上
返回任意实际答案。

$O$ 证明者从自己的行纤维出发。给定 $(O,A,T_O,S_O)$，令 $H=\ker A$，
并按固定规则选满足 $AM_0=S_O$ 的代表 $M_0$。这样的代表在实验的
每个输入上都存在。把行纤维写成 $M_0+X$，其中 $X=\Hom(L_O,H)$。
采样均匀的 $\sigma\in K^*$，并作符号函数
\[
 f_\sigma(N)=(-1)^{\sigma(F_O(M_0+N,T_O))},\qquad N\in X.
\]
以概率 $|\widehat f_\sigma(\Phi)|^2$ 采样傅里叶频率
$\Phi\in X^*$。把 $X^*=H^*\otimes L_O$ 等同，纯张量上的配对定义为
\[
 \langle\varphi\otimes v,N\rangle=\varphi(Nv).
\]
对 $h\in H$，收缩把 $\Phi$ 送到向量 $\Phi(h)\in L_O$。
若 $\sigma|_H\ne0$，按固定规则选使 $\sigma(h)=1$ 的 $h\in H$ 并
返回 $\Phi(h)$。若限制平凡，返回任意实际答案。

该规则在正质量频率处总返回实际答案。事实上
Equation~\eqref{eq:matrix-family-identities} 的折叠恒等式给出
\[
 f_\sigma(N+h e_0^\top)=-f_\sigma(N)
 \qquad\text{当 }\sigma(h)=1.
\]
在傅里叶系数中换变量可知其非零要求
$(-1)^{\langle\Phi,h e_0^\top\rangle}=-1$，等价于
$e_0^\top\Phi(h)=1$。

Figure~\ref{fig:advice} 概括两个解码器的信息与私有选择。转移来的
切片只用于分析其一致概率。
\begin{figure}[tbp]
\centering
\begin{tikzpicture}[
  x=1cm,y=1cm,
  every node/.style={font=\small},
  box/.style={draw=linkblue!80!black,rounded corners=2pt,
    text width=6.0cm,align=center,inner sep=8pt},
  input/.style={box,minimum height=16mm},
  rule/.style={box,minimum height=35mm},
  output/.style={box,minimum height=10mm},
  flow/.style={-{Latex[length=2.1mm]},thick,draw=black!70}
]
\node[box,text width=13.3cm,minimum height=12mm] (sample) at (3.7,0)
  {\textbf{裁判}\\[2pt]
   采样 $U,O,\bar\pi,A,T_O,M_O$；令 $S_O=AM_O$。
   只发送下述输入。};

\node[input] (uinput) at (0,-1.9)
  {\textbf{方程解码器}\\
   $(U,A,T_U,S_U)$\\
   $T_U=T_O\bar\pi,\quad S_U=S_O\bar\pi$};
\node[input] (oinput) at (7.4,-1.9)
  {\textbf{投影解码器}\\
   $(O,A,T_O,S_O)$};
\draw[flow] (sample.south -| uinput.north) -- (uinput.north);
\draw[flow] (sample.south -| oinput.north) -- (oinput.north);

\node[rule] (urule) at (0,-4.8)
  {从本解码器的可见建议中选取固定的切片见证\\及其合法答案集\\[2pt]
   $\mathcal A=(z+Z)\cap\{v:e_0^\top v=1\}$。\\[6pt]
   \textit{私有随机性：}\\
   从 $\mathcal A$ 中均匀采样 $a_U$。};
\node[rule] (orule) at (7.4,-4.8)
  {令 $H=\ker A$。选 $M_0$\\使 $AM_0=S_O$。\\[5pt]
   \textit{私有随机性：}\\
   采样均匀的 $\sigma\in K^*$，再以概率\\
   $|\widehat f_\sigma(\Phi)|^2$ 采样 $\Phi$。\\[5pt]
   若 $\sigma|_H\ne0$，取固定的 $h\in H$\\
   使 $\sigma(h)=1$；令 $a_O=\Phi(h)$。};
\draw[flow] (uinput) -- (urule);
\draw[flow] (oinput) -- (orule);

\node[output] (uoutput) at (0,-7.9)
  {$a_U\in L_U,\qquad e_0^\top a_U=1$};
\node[output] (ooutput) at (7.4,-7.9)
  {$a_O\in L_O,\qquad e_0^\top a_O=1$};
\draw[flow] (urule) -- (uoutput);
\draw[flow] (orule) -- (ooutput);

\node[box,text width=8.2cm,minimum height=11mm] (verify) at (3.7,-9.6)
  {\textbf{裁判检验}\quad $\bar\pi a_U=a_O$};
\draw[flow] (uoutput.south) -- (verify.north -| uoutput.south);
\draw[flow] (ooutput.south) -- (verify.north -| ooutput.south);
\end{tikzpicture}
\caption{两个解码器使用独立的随机带。二者都不接收采样到的 $M_O$
或投影 $\bar\pi$；右侧解码器从可见建议中自行选取代表 $M_0$。每个
所示的建议映射都经由齐次坐标包含其仿射截距。转移来的切片只用于
分析这些私有规则：右侧解码器不使用左侧的见证或答案集。若左侧建议
不好，或右侧 $\sigma|_H=0$，该解码器返回任意合法答案。}
\label{fig:advice}
\end{figure}

\FloatBarrier

\begin{lemma}[转移切片给出一致]
\label{lem:fourier-coset-decoding}
固定满足 Lemma~\ref{lem:projected-witness} 的联合数据。在这些数据
的条件下，刚定义的两个规则使用独立随机性，在 $\bar\pi$ 之下一
致的概率至少
\begin{equation}
 2^{-\ell r_{\rm s}}(\alpha/8)^2\,2^{-r_{\rm s}}.
 \label{eq:conditional-decoding}
\end{equation}
\end{lemma}

\begin{proof}
令 $Z'=\bar\pi Z$，故 $\dim Z'\le r_{\rm s}$。行纤维中投影后的列
规定把限制 $N|_{Z'}$ 固定为线性映射 $b:Z'\to H$。所得仿射切片
非空；取其中一个 $N_*$。对每个 $\sigma\in K^*$ 定义
\[
 C_\sigma=\E\bigl[f_\sigma(N)(-1)^{\sigma(Nz')}
                       \mid N|_{Z'}=b\bigr].
\]
Equation~\eqref{eq:transferred-witness} 的特征展开给出
\[
 \E_\sigma\bigl[(-1)^{\sigma(M_0z'+u)}C_\sigma\bigr]
 \ge\alpha/4,
 \qquad\text{故}\qquad
 \E_\sigma|C_\sigma|\ge\alpha/4.
\]
恰有 $2^{-\dim H}$ 比例的特征在 $H$ 上限制平凡。由
$\dim H\ge\ell-r_{\rm s}$ 与 $|C_\sigma|\le1$，
Equation~\eqref{eq:matrix-parameters} 蕴含
\begin{equation}
 \E_\sigma\bigl[\mathbf 1_{\{\sigma|_H\ne0\}}|C_\sigma|\bigr]
 \ge\alpha/8.
 \label{eq:nontrivial-correlation}
\end{equation}

在 $Z'$ 上为零的矩阵在 $X^*$ 中的零化子是 $H^*\otimes Z'$。令
$\theta_\sigma=(\sigma|_H)\otimes z'$，在仿射切片上的傅里叶展开
给出
\[
 C_\sigma
 =\sum_{\Phi\in\theta_\sigma+H^*\otimes Z'}
       \widehat f_\sigma(\Phi)
       (-1)^{\langle\Phi+\theta_\sigma,N_*\rangle}.
\]
该和有 $2^{\dim H\dim Z'}\le2^{\ell r_{\rm s}}$ 项。
Cauchy--Schwarz（先对每个和、再对 $\sigma$ 的平均）由
Equation~\eqref{eq:nontrivial-correlation} 给出
\begin{align*}
 &\E_\sigma\left[
   \mathbf 1_{\{\sigma|_H\ne0\}}
   \sum_{\Phi\in\theta_\sigma+H^*\otimes Z'}
             |\widehat f_\sigma(\Phi)|^2\right]\\
 &\hspace{15mm}\ge
 2^{-\ell r_{\rm s}}
 \E_\sigma\bigl[\mathbf 1_{\{\sigma|_H\ne0\}}C_\sigma^2\bigr]
 \ge 2^{-\ell r_{\rm s}}(\alpha/8)^2.
\end{align*}
这就是 $O$ 解码器采到非平凡限制与所示陪集中频率的概率下界。

每个这样的频率形如 $\Phi=(\sigma|_H)\otimes z'+\Psi$
（$\Psi\in H^*\otimes Z'$）。对解码器选取的 $h$，
\[
 \Phi(h)=z'+\Psi(h)\in z'+Z'.
\]
如已证，正质量频率给出首比特 1。Equation~\eqref{eq:projected-answer-set}
因此把返回向量置于 $\bar\pi\mathcal A$。独立地，$U$ 解码器从
$\mathcal A$ 均匀采样，以至少 $1/|\mathcal A|\ge2^{-r_{\rm s}}$
的概率选中原像。各界相乘即证
Equation~\eqref{eq:conditional-decoding}。
\end{proof}

\begin{proof}[Proposition~\ref{prop:matrix-decoding} 的证明]
使用上述两个解码器。Lemma~\ref{lem:projected-witness} 提供概率
至少 $\gamma_1$ 的联合数据集，Lemma~\ref{lem:fourier-coset-decoding}
在该集上给出所述条件一致性。于是可取
\[
 \gamma=\gamma_1\,2^{-\ell r_{\rm s}}(\alpha/8)^2\,2^{-r_{\rm s}}>0.
\]
对 $k$ 之大的需要只来自 Theorem~\ref{thm:shortcode}、
Lemma~\ref{lem:good-row-advice} 的集中界与 Lemma~\ref{lem:projected-tv}
的全变差估计。它们对奇偶实例与表族一致。见证选取与傅里叶采样不必
高效：它们在可靠性论证的有限博弈中定义策略。最后，对证明者独立
随机带的取平均也给出一致概率不低于此的确定性策略。
\end{proof}

\section{干净坐标与可靠性间隙}
\label{sec:clean}

Proposition~\ref{prop:matrix-decoding} 把矩阵测试的高接受率转化为
两个携带相关线性建议的局部策略之间的一致。现在我们在 NO 奇偶实例
上界定该建议实验的值。消失建议的使用是 \cite[Sections~3.3--3.4]{KMS17}
光滑外层博弈框架的一部分。我们针对本文的以出现为指标的问题给出
完整的局部重构论证。关键在于找出两类建议斜率都为零的单例坐标。
在这些坐标之外的数据固定后，其问题对具有普通的独立「方程对变量」
律。于是每一方都可以从自己剩余的问题与固定数据重构其原始建议输入。
这就定义了普通并行重复的策略，其值界定原一致性。

\subsection{经典投影博弈重复}

我们使用 Dinur 与 Steurer 并行重复定理的如下形式。此处的值是普通
经典值：每个证明者只用自身问题作答。投影约束把一侧的每个答案映到
其在另一侧的所需答案。

\begin{theorem}[投影博弈重复]
\label{thm:repetition}
设 $G$ 为有限二部投影博弈，其约束出现上有任意概率分布。问题集与
答案字母表可以是任意有限集，并允许并行的约束出现。博弈
$G^{\otimes n}$ 采样 $n$ 个独立出现，给每个证明者发送其问题元组，
当所有 $n$ 个投影约束都成立时接受。对每个整数 $n\ge1$ 与每个
$0<g<1$，
\begin{equation}
 \val(G)\le 1-g
 \quad\Longrightarrow\quad
 \val(G^{\otimes n})\le (1-g^2/16)^n.
 \label{eq:projection-repetition}
\end{equation}
该界与问题集及字母表规模无关。
\end{theorem}

这是 \cite{DS14} 的 Corollary~1.2。其 Section~2.1 容许非负边权与
平行边，Section~2.2 定义了此处使用的独立乘积。因此该定理适用于
我们构造中任何有限的带权问题分布。必要时交换两个证明者以匹配投影
方向。私有随机化与独立于问题的共享随机性不会增大经典值：在问题
到来之前用该随机性采样完整的确定性答案表，再对其成功概率取平均。

\subsection{干净坐标上消失的建议}

设 $\mathcal E$ 为 Section~\ref{sec:game} 的奇偶实例（每条方程出现
有三个互异变量 ID）。回忆其普通「方程对变量」博弈
$G_{\mathcal E}$：第一个证明者收到方程并以满足三元组作答；第二个
收到从该方程均匀选取的变量并以比特作答。在 NO 实例上，
Equation~\eqref{eq:base-game-gap} 给出
$\val(G_{\mathcal E})\le14/15$。

回忆 Section~\ref{sec:matrix} 的建议实验：取大于 1 的立方数 $k$，
并令 $\beta=k^{-2/3}$。有序元组 $U$ 由 $k$ 个独立的方程样本组成。
在每个块中独立地：以概率 $\beta$ 保留均匀选取的变量位置，否则保留
整条方程。这产生 $O$ 与 Equation~\eqref{eq:outer-projection} 的
答案投影 $\bar\pi:L_U\to L_O$。独立均匀地选
\[
 A:K\to\F_2^{r_{\rm s}},\qquad
 T_O:L_O\to W,\qquad M_O:L_O\to K
\]
并令 $S_O=AM_O$。两个输入分别是
\begin{equation}
 (U,A,T_O\bar\pi,S_O\bar\pi)
 \quad\text{与}\quad
 (O,A,T_O,S_O).
 \label{eq:clean-advice-inputs}
\end{equation}
$S_O$ 之外的矩阵 $M_O$ 不属于任一输入。每个证明者必须返回实际
答案，等价地其各自答案空间中首坐标为 1 的向量。

\begin{lemma}[干净坐标界]
\label{lem:clean}
若 $\OPT(\mathcal E)\le 4/5$，则建议实验中的每对局部随机化策略
满足
\begin{equation}
 \Pr\bigl[\bar\pi(a_U)=a_O\bigr]
 \le \left(1-\frac{\beta\,2^{-\dim W-r_{\rm s}}}{3600}\right)^k
 \le
 \exp\!\left(
 -\frac{2^{-\dim W-r_{\rm s}}}{3600}\,k^{1/3}
 \right).
 \label{eq:clean-agreement-bound}
\end{equation}
特别地，在建议空间固定时，该上界沿立方数随 $k\to\infty$ 对奇偶
实例一致地趋于零。
\end{lemma}

\begin{proof}
固定建议实验中任意一对局部随机化策略。先对 $A$ 条件化，把两个可见
建议映射合并为
\[
 D_O=(T_O,S_O):L_O\longrightarrow E_A,
 \qquad E_A=W\oplus\im A.
\]
在问题与投影选择的条件下，$D_O$ 的截距 $c_O$ 与所有块斜率列都是
$E_A$ 中独立的均匀元素。事实上，$T_O$ 与 $M_O$ 的列独立均匀，而
把 $A$ 作用到均匀列上给出 $\im A$ 中的均匀元素。该系数律只依赖
块是方程还是单例，不依赖任何方程 ID、变量 ID 或所选位置。特别地，
$c_O$ 与所有块数据独立。

称坐标 $i$ \emph{干净}，若它被投影为单例，且其唯一的建议斜率
$d_i\in E_A$ 为零。在 $A$ 的条件下，干净示性独立，公共概率为
\begin{equation}
 p_A=\frac{\beta}{\abs{E_A}}
     =\beta\,2^{-\dim W-\rank A}.
 \label{eq:clean-probability}
\end{equation}
更精确地，对每条方程出现 $e$ 与位置 $j\in[3]$，
\begin{equation}
 \Pr[e_i=e,\ i\text{ 是位置 }j\text{ 处的干净单例}\mid A]
 =\omega_e\,\frac{\beta}{3\abs{E_A}}.
 \label{eq:clean-joint-law}
\end{equation}
于是在「干净」的条件下，一个坐标恰具有普通的「方程对变量」分布。

设 $I\subseteq[k]$ 为干净集。为分析，令 $R$ 包含 $A$、$I$、原截距
$c_O$，以及每个 $i\notin I$ 的如下数据：其方程出现、块是被保留还是
被投影、后者情形下所选的位置、以及其建议斜率列。$R$ 不包含任何
见证、解码器输出、一致事件或应答随机性。

对每个正概率的取值 $R=r$，把 $r$ 中记录的干净集 $I$ 视为固定。
当 $I$ 非空时，我们将在定义 $G_{\mathcal E}^{\otimes\abs{I}}$ 的
两个策略时把 $r$ 用作常数数据。每个策略将由其自己的乘积博弈问题
重构一份原始建议输入，应用相应的原始应答规则，并把答案保留在干净
槽位上。

在 $R=r$ 的条件下，对
\[
 (e_i,v_{e_i,j_i}),\qquad i\in I,
\]
是 $G_{\mathcal E}$ 的独立问题对。事实上，在 $A$ 的条件下，坐标
数据独立，而 Equation~\eqref{eq:clean-joint-law} 表明以干净掩码为
条件在每个干净坐标上保持普通律。外部坐标数据与独立截距不对这些
对施加进一步条件。

现在验证两个重构，包括跨块共享的仿射截距。右端为 $b_e$ 的方程的
实际比特是 $x,y,b_e t_0+x+y$。把单例斜率 $d_i$ 拉回，在位置一、
二、三分别产生两个斜率
\begin{equation}
 (d_i,0),\quad (0,d_i),\quad (d_i,d_i)
 \label{eq:clean-padded-slopes}
\end{equation}
中的相应一个。$D_U=D_O\bar\pi$ 的截距为
\begin{equation}
 c_U=c_O+
 \sum_{\substack{i\text{ 投影为单例}\\j_i=3}}
 b_{e_i}d_i.
 \label{eq:clean-intercept}
\end{equation}
对干净坐标，Equation~\eqref{eq:clean-padded-slopes} 中每个斜率都是
零，其对 Equation~\eqref{eq:clean-intercept} 的贡献为零。该式中
所有其余项都由 $r$ 决定。

方程侧收到其对 $i\in I$ 的出现 $e_i$。连同 $r$，它们指定整个
$U$。其每个干净块中的两个建议斜率为零，而 $r$ 决定所有其他斜率与
完整截距。这些系数在不使用任何干净所选位置或变量问题的前提下
重构两张补零建议映射。变量侧收到其对每个 $i\in I$ 的变量问题。
连同 $r$，它们指定整个 $O$。其干净单例斜率为零，$r$ 提供所有其他
斜率与原截距。这些系数在不使用任何干净方程的前提下重构其两张建议
映射。不同块之间重复的 ID 不带来额外限制：问题是有放回独立采样的，
答案使用 Section~\ref{sec:game} 指定的局部副本。

这些重构定义了所述的乘积博弈策略：每侧把其原始应答规则作用于重构
后的输入，并把答案按原顺序保留在槽位 $I$ 中。其方程答案是满足
三元组，因为原答案合法。对干净问题与应答随机性的每组实现，建议
实验中的完整一致蕴含每个干净槽位被接受。若 $I=\varnothing$，条件
成功概率至多 1。否则，
\[
 \Pr[\bar\pi a_U=a_O\mid R=r]
 \le \val(G_{\mathcal E}^{\otimes\abs{I}})
 \le \left(1-\frac{1}{3600}\right)^{\abs{I}}.
\]
第一个不等式用刚构造的策略与条件乘积律。第二个用
Theorem~\ref{thm:repetition} 与 Equation~\eqref{eq:base-game-gap}。
原始应答规则只在其原始局部输入处求值；$r$ 索引辅助的乘积博弈策略。
因此最终界对每个 $r$ 都成立，包括 $I=\varnothing$（此时右端按 1
解释）。

最后，在 $A$ 的条件下，$I$ 的大小服从
$\operatorname{Binomial}(k,p_A)$。先在 $A$ 的条件下对 $R$ 取平均、
再对 $A$ 取平均，得
\begin{align*}
 \Pr\bigl[\bar\pi(a_U)=a_O\bigr]
 &\le \E_A\left(1-\frac{p_A}{3600}\right)^k\\
 &\le\left(1-\frac{\beta\,2^{-\dim W-r_{\rm s}}}{3600}\right)^k\\
 &\le \exp\!\left(
 -\frac{k\beta\,2^{-\dim W-r_{\rm s}}}{3600}
 \right),
\end{align*}
其中用了 $\rank A\le r_{\rm s}$。代入 $k\beta=k^{1/3}$ 证明
Equation~\eqref{eq:clean-agreement-bound}。
\end{proof}

\subsection{矩阵博弈间隙}

现在可以把每个建议策略的上界与从高接受率矩阵染色抽取出的策略
比较。

\begin{proposition}[可调完备性的固定间隙]
\label{prop:matrix-gap}
固定 $p_*>0$，并按 Equation~\eqref{eq:matrix-parameters} 的顺序选取
潜在与短码数据。对每个足够大的立方数 $k$，由位置互异的带权奇偶
实例 $\mathcal E$ 构造的矩阵博弈 $\mathcal G$ 满足
\[
 \begin{aligned}
 \OPT(\mathcal E)\ge1-\xi
 &\quad\Longrightarrow\quad
 \val(\mathcal G)\ge1-k\xi-p_*/2,\\
 \OPT(\mathcal E)\le4/5
 &\quad\Longrightarrow\quad
 \val(\mathcal G)\le99/100.
 \end{aligned}
\]
$k$ 的下阈值只依赖固定的潜在与短码数据，不依赖 $\mathcal E$
或 $\xi$。
\end{proposition}

\begin{proof}
完备性即 Lemma~\ref{lem:outer-completeness}。对可靠性，固定
$\mathcal G$ 的一个染色。其恢复表函数满足 Section~\ref{sec:game}
的折叠与共享恒等式。Proposition~\ref{prop:matrix-decoding} 给出
常数 $\gamma>0$，并且对每个足够大的立方数 $k$，只要那些函数的
矩阵测试接受率至少 $0.99$，就给出一致概率至少 $\gamma$ 的局部
建议策略。对所有足够大的立方数 $k$，
Equation~\eqref{eq:clean-agreement-bound} 的上界小于 $\gamma$。
于是只要 $\OPT(\mathcal E)\le4/5$，Lemma~\ref{lem:clean} 就与那些
策略的存在矛盾。所有阈值对奇偶实例一致。
\end{proof}

顺序很重要：$p_*$ 可以在 gadget 参数固定之前取得任意小，而 $\xi$
可以在 $k$ 固定之后取得任意小。因此同一个 NO 间隙与任何指定的小
完备性误差相容。

\section{完成归约}
\label{sec:finish}

矩阵博弈具有固定的可靠性间隙与任意小的完备性误差。现在得到
Theorem~\ref{thm:ugc} 所需的形式：无权、简单二部博弈、任意指定的
误差。把出现权重取整并对每个出现细分，将保持固定间隙。一次字母表
无关的重复即完成证明。

\subsection{去掉权重与重复边}

我们对带权置换约束使用如下初等构造。在这一中间阶段，约束出现可以
是自环，且多个出现可以连接相同的顶点。

\begin{lemma}[取整与细分]
\label{lem:unweighted-subdivision}
设 $H$ 为有限置换约束博弈，$h\ge1$ 个正有理出现权重之和为 1。对
每个正有理 $\tau$，令 $Q=\lceil h/\tau\rceil$。存在一个确定性构造
得到无权简单二部博弈 $B$（$4Q$ 条边、同一字母表），使
\begin{equation}
 \left|\val(B)-\left(1-\frac{1-\val(H)}4\right)\right|
 \le\frac{\tau}{4}.
 \label{eq:rounded-subdivision-value}
\end{equation}
若 $H$ 的约束是某个有限二元向量空间上的平移，则 $B$ 的约束也是。
对固定的 $\tau$，当 $H$ 以顶点、出现、置换表与二进制有理权重列出
时，该构造的运行时间在其编码长度的多项式以内。
\end{lemma}

\begin{proof}
把出现权重写成 $w_1,\ldots,w_h$。先取 $n_e=\lfloor Qw_e\rfloor$，
再把小数部分最大的 $Q-\sum_e\lfloor Qw_e\rfloor$ 个表项各加一
（平局按固定顺序处理）。则
\[
 n_e\ge0,\qquad \sum_e n_e=Q,\qquad
 |w_e-n_e/Q|\le1/Q.
\]
权重差之和为零。因此对每个染色，其被满足出现集的权重变化至多
\[
 \frac12\sum_{e=1}^h|w_e-n_e/Q|
 \le\frac{h}{2Q}\le\tau.
\]
把每个出现 $e$ 换成 $n_e$ 个副本（$n_e=0$ 时省略）。所得无权多重图
$H'$ 有 $Q$ 个出现，且 $|\val(H')-\val(H)|\le\tau$。

把 $H'$ 的每个出现 $e$ 定向为从 $u$ 到 $v$，把其约束写成
$a_v=\rho_e(a_u)$。对\emph{每个出现}用三个新鲜内部顶点把它替换为
路径
\[
 u\;--\;p_e\;--\;q_e\;--\;r_e\;--\;v.
\]
把所有原始顶点与所有 $q_e$ 放在左侧，所有 $p_e,r_e$ 放在右侧。
沿从 $u$ 到 $v$ 的方向，前三条边上放恒等置换，最后一条上放
$\rho_e$。当以左到右方向表达一条边时，必要时反转其置换。给全部
$4Q$ 条边相等的权重。

固定原始顶点处的标记。被满足的出现延拓为四条被满足的边。被违背的
出现不能延拓为四条，但沿前三条边传播标记恰满足三条。由于内部顶点
两两不同，这些延拓可以独立选取。对原始标记取最大值给出精确恒等式
\begin{equation}
 \val(B)=1-\frac{1-\val(H')}{4}.
 \label{eq:bipartization-value}
\end{equation}
这证明 Equation~\eqref{eq:rounded-subdivision-value}。

自环变成带三个不同新顶点的四元环。平行出现的内部不相交。于是每对
左右顶点至多承载一条边，故 $B$ 是简单图。逆二元平移仍是平移，恒等
即零平移；所述约束形式得以保持。

最后，$Q\le h/\tau+1$ 对固定 $\tau$ 是多项式有界的。整数除法与
有理余数的精确比较在多项式时间与位复杂度内算出诸 $n_e$。列出其
$Q$ 个副本与 $3Q$ 个新顶点因而也是多项式的。
\end{proof}

\subsection{最终归约}

\begin{proof}[Theorem~\ref{thm:ugc} 的证明]
固定 $\eps,\delta\in(0,1/2)$。选正有理 $\eps_0\le\eps$ 与
$\delta_0\le\delta$。在应用取整细分引理时将取 $\tau\le1/200$。对
任何满足 NO 界 $\val(H)\le99/100$ 的矩阵博弈 $H$，所得博弈 $B$
满足
\begin{equation}
 \val(B)\le1-\frac{1/100-\tau}{4}\le1-\frac1{800}.
 \label{eq:post-subdivision-gap}
\end{equation}
该间隙不依赖字母表。因此我们可以用字母表无关的重复，在潜在构造
固定字母表之前选定其指数。选整数 $t\ge1$ 使
\begin{equation}
 \left(1-\frac1{10240000}\right)^t\le\delta_0.
 \label{eq:final-repetition-gap}
\end{equation}
例如 $t\ge\lceil10240000\log(1/\delta_0)\rceil$ 即可。由
Theorem~\ref{thm:repetition}，该指数把每个值至多 $1-1/800$ 的投影
博弈归约到值至多 $\delta_0$。

选正有理参数
\begin{equation}
 p_*\le\frac{\eps_0}{t},\qquad
 \tau\le\min\left\{\frac1{200},\frac{\eps_0}{t}\right\}.
 \label{eq:final-errors}
\end{equation}
由 Lemma~\ref{lem:latent} 选潜在秩阈值 $r_*$，令 $\eta=2^{-r_*}/16$，
并用 $\alpha$ 的某个正有理下界选相应的短码参数 $\alpha,r_{\rm s}$。
然后按 Equation~\eqref{eq:matrix-parameters} 选足够大的 $\ell$，
固定相应的潜在数据 $(\mathcal V,K,C,\mu)$ 与分裂
$\mathcal V=K\oplus W$，其中 $K=\F_2^\ell$。

有理预算一旦固定，这些选取就是有效的。对有理 $p_*$，
Lemma~\ref{lem:latent} 的有效性断言经由一个会终止的确定性构造给出
全部有限的潜在数据，包括可计算的有理噪声律。
Theorem~\ref{thm:shortcode} 之后陈述的有效界给出 $\alpha$ 的有理
选取与维数阈值，包括 $\ell$ 固定之后的 $m_0(\eta,\ell)$。

接着选足够大的立方数 $k\ge8$ 满足 Proposition~\ref{prop:matrix-gap}。
这样的立方数可以通过核查其证明中的估计找到。记 $m=1+2k$ 与
$\beta=k^{-2/3}$；由于 $k$ 是立方数，$\beta$ 是有理的。要求
$m\ge m_0(\eta,\ell)$ 且
\[
 \frac{\alpha^2}{32}\,2^{(\ell-r_{\rm s})(m-r_{\rm s})}
 >(2r_{\rm s}+1)(\ell+m).
\]
后者是使 Lemma~\ref{lem:good-row-advice} 中并界小于 1 的充分有理
条件。对 Lemma~\ref{lem:projected-witness} 中的见证转移，要求
\[
 \frac14\left[
 \left(1+\frac{|\mathcal V|-1}{3}\beta^2\right)^k-1\right]
 <\frac{\gamma_1^2}{(1+\alpha/4)^2},
\]
其中 $\gamma_1$ 来自 Equation~\eqref{eq:projected-witness-mass}。
左端是全变差界的平方。为最终矛盾，要求
\[
 \left(1-\frac{\beta\,2^{-\dim W-r_{\rm s}}}{3600}\right)^k
 <\gamma,
\]
$\gamma$ 来自 Proposition~\ref{prop:matrix-decoding}。所有这些条件
都是有理的，且对足够大的立方数 $k$ 成立。枚举立方数直至成立即可
找到合适的 $k$；搜索代价只依赖两个指定的误差，并在使用输入公式
之前发生。

最后，固定 $k$ 后，选正有理
\begin{equation}
 \xi<\min\left\{\frac1{100},\frac{\eps_0}{2kt}\right\}.
 \label{eq:final-parity-error}
\end{equation}
每个选取只依赖两个指定的误差。特别地，$r_*$ 在 $\ell$ 之前固定，
建议维数在 $k$ 之前固定，而 $k$ 在奇偶误差 $\xi$ 之前固定。

给定 3SAT 公式，应用 Theorem~\ref{thm:parity} 与
Lemma~\ref{lem:distinct-parity}，构造 Section~\ref{sec:game} 的
矩阵博弈 $H$。Proposition~\ref{prop:matrix-gap} 给出
\[
 \begin{array}{ll}
 \text{YES：}&\val(H)\ge1-d,
       \quad d=k\xi+p_*/2\le\eps_0/t,\\[2pt]
 \text{NO：}&\val(H)\le99/100.
 \end{array}
\]
以 $\tau$ 应用 Lemma~\ref{lem:unweighted-subdivision}，得到无权
简单二部博弈 $B$。它满足
\[
 \begin{array}{ll}
 \text{YES：}&\val(B)\ge1-(d+\tau)/4,\\[2pt]
 \text{NO：}&\val(B)\le1-1/800,
 \end{array}
\]
其中 NO 界即 Equation~\eqref{eq:post-subdivision-gap}。最终输出为
$G=B^{\otimes t}$。由 Equations~\eqref{eq:projection-repetition}
与 \eqref{eq:final-repetition-gap}，其 NO 值至多 $\delta_0\le\delta$。
在 YES 输入上，逐坐标使用 $B$ 的染色。并界给出
\[
 \val(G)\ge1-\frac{t(d+\tau)}4
 \ge1-\eps_0/2\ge1-\eps.
\]

输出具有所断言的图与约束形式。其两个顶点集是 $B$ 两侧的 $t$ 重
乘积。由于 $B$ 是简单图，一对乘积顶点在每个坐标处至多确定一条边，
故至多一条乘积边。独立均匀的边采样在这些乘积边上均匀。由
Section~\ref{sec:game}，$H$ 的约束是 $K$ 上的平移；取整与细分保持
它们。偏移为 $c_1,\ldots,c_t$ 的平移之积是 $K^t=\F_2^{\ell t}$ 上
的平移 $(c_1,\ldots,c_t)$。这就是定理中的字母表 $\F_2^s$，其中
$s=\ell t\ge1$ 由指定误差固定。

\paragraph{多项式编码。}
设预处理后的奇偶实例有 $N_{\mathcal E}$ 个正方程出现。方程元组有
$N_{\mathcal E}^k$ 个。每个元组的表数、噪声结果数与扰动向量数在
上述参数选取之后固定，故列出的矩阵测试出现数 $h$ 是固定常数乘
$N_{\mathcal E}^k$。省略零概率结果，重复出现分开保留。每个权重
是 $k$ 个奇偶权重、一个固定有理噪声概率与固定均匀采样因子之积，
具有多项式二进制长度。

精确表键与折叠代表由固定维数中的线性代数计算（保留原变量与出现
ID）。至多 $2h$ 个端点出现。即使对每对键都比较以识别相同顶点，也
只需多项式时间；给所得的互异顶点连续编号。最大余数取整用
$Q=\lceil h/\tau\rceil=O(h)$ 个副本；细分创建 $3Q$ 个新顶点与
$4Q$ 条边。由于 $t$ 固定，最终乘积有 $(4Q)^t$ 条边与多项式多个
顶点。固定字母表 $\F_2^{\ell t}$ 上的完整置换表规模为常数，端点
指标有 $O(t\log(h+1))$ 比特。这些观察给出全程确定性多项式时间的
显式构造，完成证明。
\end{proof}

\section{近似推论}\label{sec:consequences}

Theorem~\ref{thm:ugc} 把已有的条件难度结果转化为 NP 难度结果。底层
的条件归约与近似阈值出自所引作者。我们给出下述精确实例模型所需的
附加论证。同时记录姊妹手稿的两个结果。对最大化问题，比值 \(r\)
指返回至少 \(r\) 倍最优值的解；对最小化问题，因子 \(c\) 指返回至多
\(c\) 倍最优值的解。下文所有因子与容差都独立于输入规模固定。在每
个使用 Theorem~\ref{thm:ugc} 的归约中，它们在字母表被选定之前决定
唯一游戏误差。若 \(\mathrm P\ne\mathrm{NP}\)，所得难度陈述排除
所示严格阈值处的确定性多项式时间近似。

\subsection{约束满足与顶点覆盖}

对 \textsc{Max-Cut}（最大化二分后跨越的边数），Khot、Kindler、
Mossel 与 O'Donnell 的归约给出：在 \((\alpha_{\mathrm{GW}},1)\) 中
每个固定比值处的近似都是 NP 难的，其中 Goemans--Williamson 常数为
\(\alpha_{\mathrm{GW}}\simeq0.87856\)~\cite[Theorem~1]{KKMO07}。
这匹配 Goemans 与 Williamson 的半定近似阈值~\cite{GW95}。该归约
使用 Mossel、O'Donnell 与 Oleszkiewicz 的 Majority Is Stablest
定理~\cite{MOO10}。对 \textsc{Vertex Cover}（最小化与每条边相交的
顶点数），Khot 与 Regev 的归约给出 \((1,2)\) 中每个固定因子处的
NP 难度~\cite{KhotRegev2008}。

约束满足问题（CSP）在此由有限的非空域 \(D\) 与 \(D\) 上有限多
个有界元的布尔谓词族 \(\Lambda\) 指定。带权 \(\Lambda\)-实例有
有限变量集与有限多个「来自 \(\Lambda\) 的谓词作用于变量元组」，
其非负有理权重之和为 1。赋值把变量映到 \(D\)；其值是被满足的
总权重，\(\operatorname{OPT}\) 是该值在所有赋值上的最大值。
Raghavendra 定理以一个基础半定松弛刻画了每个这种固定语言的最优
近似（至任意小的正误差）
\cite[Theorem~1.1 and Corollary~1.3]{Raghavendra08}。
顶点覆盖结果使用其自己的归约。把 Raghavendra 松弛的值记为
\(\operatorname{SDP}_{R}\)，如~\cite[Section~4.5]{RaghavendraThesis09}
所定义。

Raghavendra 定理的固定实例形式与下文排序定理使用唯一游戏的强化
二部表述。在其 YES 情形，一个染色满足每个左顶点关联边的至少
\(1-\zeta\) 比例；在其 NO 情形，每个染色至多满足 \(\nu\) 比例的
边。这分别是左顶点与边的均匀比例。
Lemma~\ref{lem:simple-strong-left} 在显式非空简单图（其公共正左度
可以达到任何固定下界并被任何固定正整数整除）上给出该承诺。它遵循
Khot 与 Regev 的归一化与序列构造
~\cite[Lemmas~3.3, 3.4, and~3.6]{KhotRegev2008}，外加使图为简单的
一个附加限制。其误差参数与度在字母表选取之前固定，因此
Theorem~\ref{thm:ugc} 提供所需前提。

\begin{corollary}[Raghavendra 推论]\label{cor:raghavendra-consequence}
设 \(J\) 为固定的 \(\Lambda\)-实例，
\(A=\operatorname{OPT}(J)<B=\operatorname{SDP}_{R}(J)\)。对每个
固定 \(0<\eta<(B-A)/2\)，区分满足
\[
 \operatorname{OPT}(I)\ge B-\eta
 \qquad\text{的 \(\Lambda\)-实例 \(I\) 与满足}\qquad
 \operatorname{OPT}(I)\le A+\eta
\]
的实例是 NP 难的。
\end{corollary}
\begin{proof}
按该论文的公共元约定
\cite[Definitions~7.2.1--7.2.2]{RaghavendraThesis09}，把 \(D\)
等同于固定的 \([q]\)，并把每个谓词用被忽略的尾随参数补齐到公共
正元数：新增位置填入 \(J\) 的某个固定变量（因 \(A<B\) 而存在）。
这保持每个被应用收益及其依赖集不变，因而保持
\(\operatorname{OPT}\) 与 \cite[Section~4.5]{RaghavendraThesis09}
的松弛。在该松弛中，每个局部分布定义在其被应用收益真正依赖的
变量上的赋值上。在发出一个出现时，只丢弃新增的尾随位置，保留
\(\Lambda\) 中的原谓词及其权重。

以容差 \(\eta/2\) 应用 Raghavendra 的固定实例定理
\cite[Theorem~7.9]{RaghavendraThesis09}；另见 \cite[Theorem~1.1
and Lemma~2.3]{Raghavendra08}。固定见证与容差决定其强化的唯一
游戏误差，于是上述表述与 Theorem~\ref{thm:ugc} 提供其前提，字母表
为固定的 \([M]\)。其验证者实例在 YES 情形最优值大于 \(B-\eta/2\)，
在 NO 情形至多 \(A+\eta/2\)。

我们把验证者的概率权重换成所述的有理权重。对这次固定调用，
\cite[Sections~7.3.1 and~7.5]{RaghavendraThesis09} 的构造在字母表
\([M]\) 固定之后给出有限集 \(\mathcal S\) 上的局部概率律
\(\theta\)。记录 \(s\in\mathcal S\) 指定 \(J\) 的一个被应用谓词、
其在 \(D^M\) 中的 \(a(s)\) 个有序神谕查询串、以及其应答填充的
位置。该律可以依赖 \(J,\eta,M\)，但没有变化的输入参数进入它。
有理概率向量在这个有限单纯形中稠密，故取
\(\widehat\theta\in\mathbb Q_{\ge0}^{\mathcal S}\) 使
\[
 \sum_{s\in\mathcal S}\widehat\theta(s)=1,
 \qquad
 \delta:=\frac12\sum_{s\in\mathcal S}
       |\widehat\theta(s)-\theta(s)|<\frac{\eta}{2}.
\]
这张有理表是归约的固定有限数据；其描述长度可以依赖
\(J,\eta,M\)。

设 \(L_\Phi,V_\Phi\) 为上述非空简单博弈的左右两侧，\(h\ge1\) 为其
公共左度。列出所引论文 Section~7.5 中使用的变量
\(V_\Phi\times D^M\)。其外层验证者选均匀左顶点 \(u\)，并对每个
神谕查询独立地选 \(u\) 的均匀邻居、按该边的置换重标记查询串。对
每个 \(s\in\mathcal S\)、\(u\in L_\Phi\) 与有序邻居元组
\(\in N(u)^{a(s)}\)，在所得查询变量上发出 \(s\) 规定的
\(\Lambda\)-谓词。测试中未被使用的任何语法变量，其所有位置都填
第一个列出的输出变量；被应用谓词与这类变量无关，故其值不变。给
出现权重
\[
 \frac{\widehat\theta(s)}{|L_\Phi|h^{a(s)}}.
\]
这些权重是和为 1 的非负有理数。保持相同出现分开，得到
\(|L_\Phi|\sum_{s\in\mathcal S}h^{a(s)}\) 个出现与
\(|V_\Phi||D|^M\) 个变量。局部集与每个 \(a(s)\) 固定，而图是
多项式规模的；因此这是多项式出现列表。显式的邻居列表与边置换使我
们能以多项式时间写出它。每个权重有多项式二进制长度：可以用固定
分母乘 \(|L_\Phi|h^{a(s)}\) 表示。

对赋值 \(\sigma\)，设 \(g_\sigma(s)\in[0,1]\) 为其关于这些外层
选择的条件接受概率。原实例与有理化实例用同一变量集，且
\[
 \bigl|\operatorname{val}_{\widehat\theta}(\sigma)
       -\operatorname{val}_{\theta}(\sigma)\bigr|
 =\left|\sum_{s\in\mathcal S}
       \bigl(\widehat\theta(s)-\theta(s)\bigr)g_\sigma(s)\right|
 \le\delta<\frac{\eta}{2}.
\]
取最大值即对其最优值给出同样的界。因此引入的间隙给出 YES 情形
最优值至少 \(B-\eta\)、NO 情形至多 \(A+\eta\)，如所需。
\end{proof}

\subsection{排序约束}

排序 CSP 使用输入变量的全局序，而不是从固定有限域取值的赋值。
固定整数 \(k\ge2\) 与非空真子集 \(\Pi\subset S_k\)（\(S_k\) 为
\(\{1,\ldots,k\}\) 的置换群）。在 Guruswami、H\aa stad、Manokaran、
Raghavendra 与 Charikar 的有限概率加权模型
~\cite[Definition~8.2]{GuruswamiEtAl2011OrderingCSP} 中，我们使用
显式列出的有限变量集 \(V\) 与有限的约束出现列表。其非负有理权重
用二进制分子分母编码且和为 1。每个出现是互异变量的有序 \(k\)
元组 \((v_1,\ldots,v_k)\)，线性序 \(\prec\) 在
\[
 v_{\pi(1)}\prec\cdots\prec v_{\pi(k)}
 \quad\text{对某个 }\pi\in\Pi
\]
时满足它。记 \(\OPT_\Pi(I)\) 为被满足权重的最大值。均匀随机序满足
期望权重 \(\rho_\Pi=|\Pi|/k!\)。由于 \(\OPT_\Pi(I)\le1\)，这是
\(\rho_\Pi\)-近似。对固定的 \(k\)，条件期望使该保证去随机化。当
序的任何前缀已暴露后，一条约束与之相容的相对顺序至多 \(k!\) 种。
其条件成功概率因此可以精确计算；选下一个变量以保持期望值只需
多项式多次、多项式位长的精确有理运算。\(V\) 的规模属于输入，
故这不是前一推论中固定有限域上的 CSP。

\begin{corollary}[排序 CSP 推论]\label{cor:ordering-csp}
固定 \(k\ge2\) 与 \(\varnothing\ne\Pi\subsetneq S_k\)，令
\(\rho_\Pi=|\Pi|/k!\)。对每个固定 \(0<\varepsilon<(1-\rho_\Pi)/2\)，
区分满足
\[
 \OPT_\Pi(I)\ge1-\varepsilon
 \qquad\text{与}\qquad
 \OPT_\Pi(I)\le\rho_\Pi+\varepsilon
\]
的有限带权 \(\Pi\)-排序 CSP 实例 \(I\) 是 NP 难的。因此，每个
固定比值 \(r\in(\rho_\Pi,1]\) 处的近似都是 NP 难的。连同随机排序
保证，这给出最优阈值 \(\rho_\Pi\)。
\end{corollary}
\begin{proof}
Proposition~\ref{prop:finite-ordering-transfer} 从简单强左唯一游戏
实例构造所述有限实例。它使用 Guruswami 等人的局部间隙与解析估计
~\cite{GuruswamiEtAl2011OrderingCSP}，并在字母表选取之前固定所需的
博弈误差与度条件。Lemma~\ref{lem:simple-strong-left} 从
Theorem~\ref{thm:ugc} 提供该博弈；复合归约即对原谓词 \(\Pi\) 证得
所示间隙。

对固定的 \(r\in(\rho_\Pi,1]\)，把正误差取得更小使
\(\varepsilon<(r-\rho_\Pi)/(1+r)\)。则
\(r(1-\varepsilon)>\rho_\Pi+\varepsilon\)，于是 \(r\)-近似将区分
两种情形。
\end{proof}

对有限带权有向图模型中的 \textsc{Maximum Acyclic Subgraph}，取
\(k=2\)、\(\Pi=\{12\}\)（一行记号）。约束 \((u,v)\) 要求边指向前
方；前向边构成无圈子图，而每个无圈子图都有拓扑序。故其阈值为
\(1/2\)。对 \textsc{Betweenness}，约束 \((a,b,c)\) 要求 \(b\)
位于 \(a\) 与 \(c\) 之间，故 \(\Pi=\{123,321\}\subset S_3\)，阈值
为 \(2/3!=1/3\)
\cite[Theorem~1.1 and Section~1.2]{GuruswamiEtAl2011OrderingCSP}。
这些有限带权推论具有固定正误差的近完美完备性。

\subsection{割、删除与聚类推论}

对 Chawla、Krauthgamer、Kumar、Rabani 与 Sivakumar 的割问题
~\cite[Section~1]{CKKRS2006}，设 \(G=(V,E)\) 为有限无向图，正有理
边代价 \(c_e\)，\(\mathcal D=\{\{s_i,t_i\}:1\le i\le k\}\) 含
\(k\ge1\) 个互异的需求对（\(s_i\ne t_i\)）。对 \(M\subseteq E\)，
记 \(c(M)=\sum_{e\in M}c_e\)，\(q_{\mathcal D}(M)\) 计
\(G-M\) 中断开的需求对数。带权 \textsc{Multicut} 在
\(q_{\mathcal D}(M)=k\) 的条件下最小化 \(c(M)\)。其非均匀
\textsc{Sparsest Cut} 搜索问题在 \(q_{\mathcal D}(M)>0\) 的割集
\(M\) 上最小化 \(c(M)/q_{\mathcal D}(M)\)。只有被选的对携带单位
需求；它们不必是全部顶点对。

对带权 Min-\(2\mathrm{CNF}^{\equiv}\) Deletion，输入是带正有理
权重的子句 \((\ell\equiv\ell')\)（\(\ell,\ell'\) 为布尔文字），
赋值支付假子句的总权重。允许的子句因此表达变量的相等或不等。

在 \textsc{Correlation Clustering} 的任意带权图形式中，每条出现的
边有 \(+\) 或 \(-\) 标签与非负有理权重。分成任意多个簇的划分支付
「端点落在不同簇的 \(+\) 边」与「端点落在同一簇的 \(-\) 边」。
缺失的边不产生惩罚。这是 \cite[Section~2.1]{DemaineEtAl2006Clustering}
的最小分歧目标。

\begin{corollary}[割、删除与聚类难度]
\label{cor:cut-clustering}
对每个固定实常数 \(C>1\)，上述四个最小化问题中每一个在因子 \(C\)
内的近似都是 NP 难的。对 \textsc{Sparsest Cut}，这是搜索任务在
多项式时间 Cook 归约下的难度：算法必须输出割集，而不只是最优值的
估计。因此除非 \(\mathrm P=\mathrm{NP}\)，这些问题都不容许任何
固定常数因子的确定性多项式时间近似。
\end{corollary}
\begin{proof}
\cite{CKKRS2006} 的 Corollary~1.3 在其普通的边值唯一游戏假设下给
出前三个结论。对请求的因子 \(C\)，在任一更大的固定因子 \(L\) 处
使用其难度陈述：\(C\)-近似也是 \(L\)-近似。为提供该假设，把
Theorem~\ref{thm:ugc} 的二部约束定向为左到右，把其固定字母表重新
编号为 \([d]\)（\(d=2^s\ge2\)），并给其边权和为 1 的相等权重。
必要时用孤立问题补齐较小的顶点侧，这保持值与多项式规模。
\textsc{Sparsest Cut} 的搜索版本限定见所引来源的 Section~1.3。

对 \textsc{Correlation Clustering}，\cite[Section~2.1]{CKKRS2006}
中困难的 \textsc{Multicut} 实例的需求对是 \(d\)-立方内互异的
对径对，而 Multicut 图的边是立方体的边或连接不同立方体。此处
\(d\ge2\)（该来源 Lemma~1.5 不改变这一点），故需求对是非边。
Demaine、Emanuel、Fiat 与 Immorlica 的变换
~\cite[Theorem~4.7 and Corollary~4.8]{DemaineEtAl2006Clustering}
因此不会产生异号的平行边。它保持最优值，并把每个聚类转换为代价
不超过它的 multicut，故同样的固定因子难度得以转移。
\end{proof}

独立手稿 \cite[Theorem~1.1]{OpenAIUniformSparsest2026} 证明：对
「非负有理容量、每对互异顶点之间恰为单位需求」的二分目标，每个
固定因子 \(C>1\) 处都是 NP 难的。这一均匀需求定理不依赖
Theorem~\ref{thm:ugc}，也不由上述非均匀归约提供。

\paragraph{核聚类。}
独立的姊妹应用关心如下最大化问题。对\emph{恒等目标核聚类}，固定
整数 \(k\ge3\)。输入是显式表示的有理对称矩阵
\(A=(a_{pq})_{p,q=1}^N\succeq0\) 且 \(A\mathbf1=0\)（\(\mathbf1\)
为全一向量）。赋值 \(\sigma:\{1,\ldots,N\}\to\{1,\ldots,k\}\) 的
值为
\[
 \val_A(\sigma)=\sum_{p,q=1}^N
 a_{pq}\mathbf1_{\{\sigma(p)=\sigma(q)\}},
 \qquad
 \OPT_k(A)=\max_\sigma\val_A(\sigma).
\]
求和包含有序对与对角线；允许空簇，且不按簇大小归一化。损失因子
\(\alpha\ge1\) 指返回至少 \(\OPT_k(A)/\alpha\) 的值。高斯螺旋桨
姊妹篇用 Theorem~\ref{thm:ugc} 证明：对每个固定 \(k\ge3\)，达到
每个固定损失因子
\[
 1\le\alpha<\alpha_k,
 \qquad \alpha_k=\frac{8\pi}{9}\left(1-\frac1k\right)
\]
是 NP 难的~\cite[Theorem~6.1]{OpenAIGaussianPropeller2026}。
这是低于螺旋桨舍入因子 \(\alpha_k\) 的严格难度。

\appendix
\section{推论的有限表述}
\label{app:finite-consequences}

Section~\ref{sec:consequences} 的部分归约使用这样的二部唯一游戏
实例：近满足的染色满足大多数左顶点关联的全部约束。我们先得到具有
该性质、且公共左度有指定下界的简单图。随后用它给出排序推论的有限
有理表述。局部排序构造遵循 Guruswami、H\aa stad、Manokaran、
Raghavendra 与 Charikar~\cite{GuruswamiEtAl2011OrderingCSP}，使用
下文所引的低影响结果。复合过程指明图、独立的邻居选取，以及到互异
变量上的约束的转换。

\subsection{唯一游戏的简单强左形式}

把二部博弈的约束从左侧定向到右侧。若某个染色的满足一个左顶点关联
的每条边，则称该左顶点被\emph{强满足}。

\begin{lemma}[简单强左博弈]\label{lem:simple-strong-left}
固定 \(0<\zeta,\nu<1\) 与整数 \(d,H\ge1\)。存在整数
\(h,R\ge1\)（\(h\ge H\) 且 \(d\mid h\)）与一个确定性多项式时间
归约，从 \(\mathrm{3SAT}\) 到 \([R]\) 上非空、显式、无权的简单
二部唯一游戏实例 \(\Phi_\varphi=(A,B,E)\)，具有以下性质。每个左
顶点的度为 \(h\)。若 \(\varphi\) 可满足，则一个染色强满足 \(A\)
的至少 \(1-\zeta\) 比例。若 \(\varphi\) 不可满足，则
\(\val(\Phi_\varphi)<\nu\)。度 \(h\) 与 Theorem~\ref{thm:ugc}
所用的两个误差可以只由 \(\zeta,\nu,d,H\) 选定，先于该定理给出
字母表 \(R\)。
\end{lemma}

\begin{proof}
我们改编 Khot 与 Regev 的归一化与序列构造
\cite[Lemmas~3.3, 3.4, and~3.6]{KhotRegev2008}。下文对右顶点使用
的多份副本保证最终图为简单图。

选 \(d\) 的倍数 \(h\) 使
\[
 h\ge\max\{H,2\},\qquad \frac1h<\frac{\nu}{4}.
\]
对 \((T)_h=T(T-1)\cdots(T-h+1)\)，选固定整数 \(T\ge h\) 使
\[
 p_{\mathrm{graph}}:=\frac{(T)_h}{T^h}>\frac34.
\]
这样的选取存在，因为 \(1-(T)_h/T^h\le\binom h2/T\)。
选正有理数 \(\beta,\gamma,u\) 与正整数 \(L_2\)，使
\[
 \beta<\frac{\nu}{4},\qquad
 h^2\gamma<\frac{\nu}{4},\qquad
 (h+1)u<\frac{\zeta}{2},\qquad
 \frac1{L_2}<\min\left\{\frac{\beta\gamma}{2},
                            \frac{\zeta}{2h}\right\}.
\]
最后选正有理数 \(\gamma_0,\zeta_0<1/2\) 使
\[
 2\gamma_0<\frac{\beta\gamma}{2},
 \qquad \sqrt{2\zeta_0}<u.
\]
以上选取都在以完备性误差 \(\zeta_0\)、可靠性误差 \(\gamma_0\)
调用 Theorem~\ref{thm:ugc} 之前完成。

把其输出写成 \(\Phi_0=(X_0,Y,E_0)\)，每个置换定向为左到右，并把
固定字母表重新编号为 \([R]\)。令 \(M_0=|E_0|>0\)，给每条边权重
\(1/M_0\)，令 \(w_x\) 为 \(x\) 处的关联总权重。简单性给出每对正
权重处恰有一个置换。创建
\[
 n_x=\lfloor 2|X_0|w_x\rfloor
\]
个 \(x\) 的副本，各带局部边权 \(w_{xy}/w_x\)。只有 \(w_x>0\) 的
顶点有副本。设 \(X_1\) 为副本集，则
\[
 |X_0|\le |X_1|=\sum_x n_x\le2|X_0|,
\]
且每个副本处的关联权重为 1。

此处及下文，局部被满足权重使用该左顶点处的标记与共享右顶点处的
标记。给定副本博弈的任意染色，保留右标记，并给每个 \(x\) 的所有
副本一个使其公共局部目标最大化的标记。相应的 \(\Phi_0\) 的染色
\(L\) 满足
\[
 \sum_{x'\in X_1} w_{\mathrm{sat}}(x')
 \le
 \sum_{x:n_x>0}\frac{n_x}{w_x}w_{\mathrm{sat},L}(x)
 \le 2|X_0|\,w_{\mathrm{sat},L}(\Phi_0).
\]
因此副本博弈的归一化值在 NO 输入上至多 \(2\gamma_0\)。在 YES
输入上，提升 \(\Phi_0\) 的一个未满足权重至多 \(\zeta_0\) 的染色。
对未满足权重作同样计算，得平均局部损失至多 \(2\zeta_0<u^2\)。
Markov 不等式留下 \(X_1\) 中超过 \(1-u\) 比例的副本，其局部被满足
权重至少 \(1-u\)。

为去掉权重，令 \(\alpha=L_2|Y|\)。在每个 \(x\in X_1\) 处，把其
归一化局部权重记为 \(\widetilde w_{xy}\)。选一个局部权重为正的
邻居 \(y_0\)，对每个 \(y\ne y_0\) 把 \(\alpha\widetilde w_{xy}\)
向下取整，并把剩余的出现都指派给 \(y_0\)。同一对上的所有出现保留
该对的置换。这给出左度 \(\alpha\) 的无权多重图
\(\Phi_{\mathrm{occ}}\)。对邻居的任何子集，其旧权重与新出现比例
之差的绝对值至多 \(|Y|/\alpha=1/L_2\)：若该子集不含 \(y_0\)，对
子集内求取整损失之和；否则对其补集求和。这对每个染色的被满足子集
都适用。故
\begin{equation}
 \val(\Phi_{\mathrm{occ}})\le 2\gamma_0+\frac1{L_2}<\beta\gamma
 \quad\text{在 NO 输入上。}
\label{eq:strong-left-base-soundness}
\end{equation}
在 YES 输入上，上述染色在超过 \(1-u\) 比例的 \(X_1\) 处的被满足
出现比例至少 \(1-a\)，其中 \(a=u+1/L_2\)。

把 \(x\) 处的 \(\alpha\) 个出现用 \(j\) 索引，右端点为 \(y_j\)、
置换为 \(\pi_j\)。把每个 \(y\in Y\) 换成 \((y,t)\)（\(t\in[T]\)），
并把出现 \(j\) 换成从 \(x\) 到 \((y_j,t)\) 的全部 \(T\) 个出现
（置换 \(\pi_j\)）。称所得左 \(\alpha T\)-正则多重图为
\(\Phi_{\mathrm{tag}}\)。对 \(\Phi_{\mathrm{tag}}\) 的任意染色，
保留其左标记，并对每个原 \(y\) 独立地给 \(y\) 一个均匀所选副本
\((y,t)\) 的标记。\(\Phi_{\mathrm{occ}}\) 中的期望边值（计其出现）
恰是 \(\Phi_{\mathrm{tag}}\) 中的给定边值。故
\begin{equation}
 \val(\Phi_{\mathrm{tag}})\le\val(\Phi_{\mathrm{occ}})<\beta\gamma
 \quad\text{在 NO 输入上。}
\label{eq:strong-left-tag-value}
\end{equation}
YES 染色的与标签无关的提升保持每个局部被满足比例。

考虑完整的序列图：其左顶点为
\[
 \bigl(x,(j_1,t_1),\ldots,(j_h,t_h)\bigr),
 \qquad j_i\in[\alpha],\quad t_i\in[T],
\]
位置 \(i\) 经 \(\pi_{j_i}\) 连到 \((y_{j_i},t_i)\)。这是 Khot 与
Regev 的序列构造；其左顶点标识符包含整个有序出现序列。我们以允许
左标记依赖整个序列的形式记录可靠性估计。每个 \(x\) 恰贡献
\((\alpha T)^h\) 个完整序列顶点，度全为 \(h\)。因此均匀边值是：
对均匀 \(x\)、其处一列有序的 \(h\) 个独立均匀出现、以及该元组中
的均匀位置，取被满足比例的平均。

固定 \(\Phi_{\mathrm{tag}}\) 中的任意右标记。\(x\) 处的每个出现
决定唯一满足它的左标记。设 \(p_x(\ell)\) 为决定标记 \(\ell\) 的
比例。诸 \(p_x(\ell)\) 之和为 1，且在每个 \(x\) 处选最好的左标记
给出
\[
 \mathbb E_x\max_\ell p_x(\ell)\le\val(\Phi_{\mathrm{tag}})<\beta\gamma.
\]
因此少于 \(\beta\) 比例的 \(x\) 有
\(\max_\ell p_x(\ell)\ge\gamma\)。在每个其余的 \(x\) 处，\(h\) 个
独立出现所决定的两个标记一致的概率至多
\[
 \binom h2\sum_\ell p_x(\ell)^2
 \le \binom h2\max_\ell p_x(\ell)<h^2\gamma.
\]
若没有任何两个被决定的标记一致，则序列顶点的任一标记至多满足其
\(h\) 条边中的一条。在其余序列上其被满足比例至多为 1。先对序列
取平均、再对 \(x\) 取平均，并用 \(1/h+h^2\gamma<1\)，把完整序列
图的值界定在
\begin{equation}
 s:=\beta+\frac1h+(1-\beta)h^2\gamma<\frac{3\nu}{4}.
\label{eq:strong-left-full-value}
\end{equation}
该估计对右标记的每种取法、以及左标记的每种依赖序列的取法都
成立。

只保留 \(t_1,\ldots,t_h\) 两两不同的序列顶点。对每个 \(x\) 与每
个基序列 \((j_1,\ldots,j_h)\)，恰有 \((T)_h\) 个标签序列存活。
每个 \(x\) 因此贡献 \(\alpha^h(T)_h\) 个存活顶点，占其完整集合的
比例 \(p_{\mathrm{graph}}\)。每个存活顶点的 \(h\) 个右邻居互异
（因标签互异）。不同的出现序列仍是不同的左顶点，故该存活图是
简单图且公共左度为 \(h\)。

给基于 \(x\) 的每个序列顶点 \(x\) 的旧标记，并保持与标签无关的
右标记。对这个提升的 YES 染色，位置 \(i\) 处的满足只依赖
\(j_i\)。把标签限制为互异，使基出现序列保持独立均匀。若
\(\Phi_{\mathrm{occ}}\) 中 \(x\) 处的被满足比例为 \(f_x\)，则恰
有 \(f_x^h\) 比例的其存活序列顶点被强满足。由于每个 \(x\) 的贡献
相同，强满足比例至少为
\[
 (1-u)(1-a)^h
 \ge 1-u-ha
 =1-(h+1)u-\frac h{L_2}>1-\zeta.
\]
此处 \(0<a<1\) 由参数选取得到。

对 NO：把存活图的任意染色延拓到完整图——给被丢弃的左顶点任意
标记并保留所有右标记。所有左度都是 \(h\)，故其完整边值至少是其
存活边值的 \(p_{\mathrm{graph}}\) 倍。由
\eqref{eq:strong-left-full-value}，后者至多
\[
 \frac{s}{p_{\mathrm{graph}}}<\frac{3\nu/4}{3/4}=\nu.
\]
该延拓容许存活的左标记任意依赖其标签。

所有操作都是有限且确定性的。设 \(N\) 界定 \(\Phi_0\) 的显式规模，
则 \(|X_1|=O(N)\)、\(\alpha=L_2|Y|=O(N)\)。存活图有
\(|X_1|\alpha^h(T)_h=O(N^{h+1})\) 个左顶点、每个左顶点 \(h\) 条
边。其标识符与取整所用的有理算术具有多项式位长，其置换表使用从
\(\Phi_0\) 继承的固定字母表。这证明了所述多项式界并完成归约。
\end{proof}

\subsection{有限排序转移}

Guruswami 等人的排序定理使用多尺度局部实例与低影响舍入估计。其
一般的外层复合被描述为类似其 Section~7 中的双问题构造。下面的
命题针对 Section~\ref{sec:consequences} 的有理出现列表模型给出该
复合。它把为排序间隙选取的常数与随后选取的唯一游戏字母表分开。

\begin{proposition}[有限排序转移]
\label{prop:finite-ordering-transfer}
固定 \(k\ge2\) 与 \(\varnothing\ne\Pi\subsetneq S_k\)，令
\(\rho_\Pi=|\Pi|/k!\)，并固定 \(0<\delta<(1-\rho_\Pi)/2\)。存在
只依赖 \(k,\Pi,\delta\) 的整数 \(q,H\ge1\) 与正有理数
\(\zeta,\nu<1\)，对每个固定整数 \(R\ge1\) 有如下性质。

设 \(\Phi=(A,B,E)\) 为 \([R]\) 上的非空显式无权简单二部唯一游戏
实例（从 \(A\) 定向到 \(B\)），其公共左度 \(h\) 满足 \(h\ge H\)
且 \(q\mid h\)。存在确定性变换（对固定 \(R\) 在 \(\Phi\) 的显式
规模上多项式），产生显式的 \(\Pi\)-排序约束列表：约束内变量互异，
带非负有理权重（每个用二进制分子与正二进制分母编码），权重之和为
1。若某染色强满足至少 \((1-\zeta)|A|\) 个左顶点，则
\(\OPT_\Pi(I)\ge1-\delta\)。若 \(\val(\Phi)\le\nu\)，则
\(\OPT_\Pi(I)\le\rho_\Pi+\delta\)。常数 \(q,H,\zeta,\nu\) 与
\(R,h\) 都无关。
\end{proposition}

输出变量将是 \((b,z)\in B\times[m]^R\)（\([m]\) 为下文选定的固定
局部符号集）。在每个 \(a\in A\) 处，把全局序限制到
\(N(a)\times[m]^R\)，并把该局部有序表分成 \(q\) 个等长连续块。
对邻居取块隶属的平均（其边置换作用于 \(z\)），得到 \(z\) 的一个
单纯形值函数，其分量在均匀 \(z\) 下均值为 \(1/q\)。对固定的
\(a\)，在局部查询串的条件下，独立的邻居选取使块秩独立、概率由该
函数给出。局部估计在小影响下界定其期望粗收益；碰撞估计容许相关
的串，并界定共享块的概率（粗化可能在共享块上改变诱导序）。平均秩
区间中的大影响反而给出博弈标记，而许多这样的顶点与低边值矛盾。
随后以互异邻居为条件得到合法约束，并把粗完备映射细化成严格序。

对 \(P=\mathbf1_\Pi\)，选 \(\pi_0\in\Pi\) 并设
\[
 P_0(\sigma)=P(\pi_0\circ\sigma),\qquad \sigma\in S_k.
\]
则 \(P_0(\mathrm{id})=1\)，其平均为 \(\rho_\Pi\)。
对整数秩向量 \(r=(r_1,\ldots,r_k)\)，记 \(\operatorname{Ref}(r)\)
为与其弱序相容的置换：它们按秩递增列出位置，相同秩的位置可以任意
排序。定义
\begin{equation}
 \overline{P_0}(r)=
 \frac1{|\operatorname{Ref}(r)|}
 \sum_{\sigma\in\operatorname{Ref}(r)}P_0(\sigma).
\label{eq:ordering-tie-payoff}
\end{equation}
这是 \cite[Section~8.2]{GuruswamiEtAl2011OrderingCSP} 使用的均匀
平局延拓。它落在 \([0,1]\) 中，且秩互异时与 \(P_0\) 一致。

该归一化在元组位置经一次置换后保持谓词。给定原始元组
\(w=(w_1,\ldots,w_k)\)，定义发出的元组 \(v\) 为
\begin{equation}
 v_{\pi_0(j)}=w_j\qquad(j\in[k]).
\label{eq:ordering-slot-permutation}
\end{equation}
若严格序在 \(w\) 上诱导 \(\sigma\)，则它在 \(v\) 上诱导
\(\pi_0\circ\sigma\)。更一般地，对位置重标记给出双射
\[
 \operatorname{Ref}\bigl(f(v_1),\ldots,f(v_k)\bigr)
 =
 \{\pi_0\circ\sigma:
       \sigma\in\operatorname{Ref}(f(w_1),\ldots,f(w_k))\}
\]
（对每个秩映射 \(f\)）。这些集合上的均匀平均因此给出相等的
\(P\)-与 \(P_0\)-收益，包括平局。\eqref{eq:ordering-slot-permutation}
中的位置置换保持互异性与权重。

对有限局部实例 \(J=(U,\mathcal C,\omega)\)，把出现
\(c\in\mathcal C\) 的元组记为 \(t_c=(u_{c,1},\ldots,u_{c,k})\)，
其概率权重记为 \(\omega_c\)。其 \(q\)-排序值为
\[
 \operatorname{opt}_q(J)=
 \max_{g:U\to[q]}
 \sum_{c\in\mathcal C}\omega_c\,
 \overline{P_0}\bigl(g(u_{c,1}),\ldots,g(u_{c,k})\bigr).
\]

\begin{lemma}[有理局部排序数据]\label{lem:ordering-local-data}
对每个整数 \(q\ge1\) 与有理 \(\eta>0\)，存在有限局部实例
\(J=(U,\mathcal C,\omega)\)（\(|U|=m\ge2\)）与 \([m]^k\) 上的
概率律 \(\mu_c\)，具有以下性质。每个元组 \(c\) 有 \(k\) 个互异
变量；权重 \(\omega_c\) 与 \(\mu_c\) 的全部原子都是有理的。
\(\mu_c\) 的每个一位边际在 \([m]\) 上均匀，且存在固定的
\(0<\alpha\le1/2\) 使 \(\mu_c(x)>\alpha\)（对所有 \(c,x\)）。
此外
\begin{equation}
 \operatorname{opt}_q(J)\le\rho_\Pi+\eta,\qquad
 v_{\mathrm{loc}}:=
 \sum_c\omega_c\mathbb E_{X\sim\mu_c}\overline{P_0}(X)
 \ge1-\eta.
\label{eq:ordering-local-gap}
\end{equation}
这些局部律为 \cite[Section~8.3]{GuruswamiEtAl2011OrderingCSP} 的
松弛提供公共的可行 Gram 数据：存在单位向量 \(I\) 与向量
\(b_{u,j}\)（\(u\in U,j\in[m]\)）使
\[
 \sum_j b_{u,j}=I,\quad
 \langle b_{u,i},b_{u,j}\rangle=0\ (i\ne j),\quad
 \|b_{u,j}\|^2=\frac1m,
\]
且对 \(c\) 的每两个位置 \(r,s\)，
\begin{equation}
 \langle b_{u_{c,r},i},b_{u_{c,s},j}\rangle
 =\Pr_{X\sim\mu_c}[X_r=i,\ X_s=j].
\label{eq:ordering-local-moments}
\end{equation}
所有这些数据都由 \(k,\Pi,q,\eta\) 固定。
\end{lemma}

\begin{proof}
\cite[Theorem~11.1, Definition~11.2, and the proof of
Theorem~2.5]{GuruswamiEtAl2011OrderingCSP} 中显式的递增元组构造
（应用于 \(P_0\)）给出有限实例 \(J^*\)：
\(\operatorname{opt}_q(J^*)\le\rho_\Pi+\eta\)，且存在满足每条约束
的线性序。其分布只用有限集上的均匀选取，故其权重是有理的。其元组
位置落在相继的不相交区间中，因而互异。

使用 \cite[Lemma~8.9]{GuruswamiEtAl2011OrderingCSP} 证明中的有限
副本平滑，特化到这条完整序上。取 \(J^*\) 的 \(L\) 个不相交副本
（各权重 \(1/L\)），把它们的完整序按某个序 \(\sigma:U\to[m]\)（
\(m=L|U^*|\)）连续排列。所得实例 \(J\) 仍满足
\(\operatorname{opt}_q(J)\le\rho_\Pi+\eta\)，因为任何 \(q\)-赋值
的值是其各副本上值的平均。选 \(L\) 与正有理 \(s_0<1\) 使
\(1/L+s_0<\eta\)。定义随机赋值 \(X:U\to[m]\) 如下：以概率
\(1-s_0\) 取 \(\sigma\) 的秩的一个均匀循环移位；以概率 \(s_0\)
从 \([m]\) 独立均匀地给变量赋值。

每个变量在混合的任一部分都均匀。一次循环切割至多切开 \(L\) 个
连续副本中的一个，故每个移位的收益至少 \(1-1/L\)。混合的收益
因此至少 \((1-s_0)(1-1/L)>1-\eta\)。令 \(\mu_c\) 为其在 \(c\)
的变量上的限制。若 \(n_c(x)\) 是 \(m\) 个循环移位中限制为 \(x\)
的个数，则
\[
 \mu_c(x)=(1-s_0)\frac{n_c(x)}m+\frac{s_0}{m^k}.
\]
这些原子是有理的且超过 \(\alpha:=s_0/(2m^k)\)。分母 \(m^k\)
合法，因为 \(c\) 中的变量互异。

最后，把有限混合视为概率空间，在其 \(L^2\) 空间中取
\(I=\mathbf1\) 与 \(b_{u,j}=\mathbf1_{\{X(u)=j\}}\)。均匀边际、
同一变量不同标签的不相交事件、以及到每个 \(c\) 的限制，给出全部
所示的 Gram 恒等式。这就为这个特定的局部实例构造了所需的可行
数据。
\end{proof}

对均匀乘积 \([m]^R\) 上的标量函数，用
\[
 \operatorname{Inf}_\ell(F)
 =\mathbb E\!\left[\operatorname{Var}_{z_\ell}
       (F\mid z_{-\ell})\right].
\]
噪声算子 \(T_\theta\)（\(0\le\theta\le1\)）以概率 \(\theta\)
独立保留每个坐标、否则均匀重采样，再对函数取平均。在正交规范乘积
展开中
\begin{equation}
 T_\theta F=\sum_\sigma\theta^{|\sigma|}\widehat F(\sigma)\chi_\sigma,
 \qquad
 \operatorname{Inf}_\ell(F)
 =\sum_{\sigma:\sigma_\ell\ne0}\widehat F(\sigma)^2.
\label{eq:ordering-noise-influence}
\end{equation}
特别地，进一步平滑会降低每个影响。

\begin{lemma}[均匀碰撞估计]\label{lem:ordering-collision}
对每个 \(0<\kappa<1/2\)，存在 \(\mu_0(\kappa)>0\)。对每个
\(m\ge2\)，存在非负函数 \(r_c^{m,\kappa}(t)\to0\)
（\(t\downarrow0\)）具有下述性质（对所有 \(R\ge1\) 一致）。设
\(F,G:[m]^R\to[0,1]\) 有相同的均值 \(0<\mu\le\mu_0(\kappa)\) 且
\[
 \operatorname{Inf}_\ell(T_{1-\kappa}F)\le t,\qquad
 \operatorname{Inf}_\ell(T_{1-\kappa}G)\le t
 \quad(\ell\in[R]).
\]
对两个边际都在 \([m]^R\) 上均匀的任何耦合 \((Z,Z')\)，
\begin{equation}
 \mathbb E\bigl[
   (T_{1-2\kappa}F)(Z)(T_{1-2\kappa}G)(Z')\bigr]
 \le \mu^{1+\kappa/2}+r_c^{m,\kappa}(t).
\label{eq:ordering-collision}
\end{equation}
函数 \(r_c^{m,\kappa}\) 与 \(R,\mu,F,G\) 及耦合无关。
\end{lemma}

\begin{proof}
这是 \cite[Lemmas~3.4 and~3.5]{GuruswamiEtAl2011OrderingCSP} 此处
所需的维数一致形式。对 \(H=T_{1-\kappa}F\)，傅里叶公式给出
\[
 \|T_{1-2\kappa}F\|_2^2
 \le \operatorname{Stab}_{1-\kappa}(H),
\]
其中 \(\operatorname{Stab}_\theta(H)
=\sum_\sigma\theta^{|\sigma|}\widehat H(\sigma)^2\)。
事实上，
\((1-\kappa)^3-(1-2\kappa)^2
=\kappa(1-\kappa-\kappa^2)>0\)。
\(H\) 的每个完整坐标影响界蕴含 Mossel、O'Donnell 与 Oleszkiewicz
定理~\cite{MOO10} 中相应的低次影响界。以乘积维数 \(R\)、单坐标
原子界 \(1/m\)、相关 \(1-\kappa\) 应用该定理。对足够小的 \(t\)
它给出
\[
 \operatorname{Stab}_{1-\kappa}(H)
 \le \Gamma_{1-\kappa}(\mu)
 +C\frac{\log(2m)}{\kappa}
       \frac{\log\log(1/t)}{\log(1/t)}.
\]
截断与该误差都与 \(R\) 和 \(\mu\) 无关。此处
\(\Gamma_\theta(\mu)\) 是高斯测度为 \(\mu\) 的集合的高斯半空间
噪声稳定性。

对固定的 \(\kappa\)，同一来源的 Theorem~B.5 给出：当
\(\mu\downarrow0\) 时
\(\Gamma_{1-\kappa}(\mu)=\mu^{2/(2-\kappa)+o(1)}\)。由于
\[
 \frac{2}{2-\kappa}-\left(1+\frac{\kappa}{2}\right)
 =\frac{\kappa^2}{2(2-\kappa)}>0,
\]
存在与 \(m,R\) 无关的截断 \(\mu_0(\kappa)>0\)，在其之下
\(\Gamma_{1-\kappa}(\mu)\le\mu^{1+\kappa/2}\)。这证明对 \(F\) 与
\(G\) 的相应平方 \(L^2\) 界，误差 \(r_c^{m,\kappa}(t)\to0\) 与
所述变量无关。在所示界的范围内选截断 \(t_0<e^{-e}\)，对
\(0<t<t_0\) 用其非负误差，对 \(t\ge t_0\) 置
\(r_c^{m,\kappa}(t)=1\)。后一范围平凡成立，因为乘积至多 1。在任何
具有指定边际的耦合下用 Cauchy--Schwarz 即证
\eqref{eq:ordering-collision}。
\end{proof}

记 \(\mathcal S_q=\{p\in[0,1]^q:\sum_jp_j=1\}\)，并在 \([q]^k\)
上定义 \(\overline{P_0}\) 的多重线性延拓
\begin{equation}
 P_0^{\mathrm{ml}}(p_1,\ldots,p_k)
 =\sum_{r\in[q]^k}\overline{P_0}(r)
       \prod_{i=1}^k p_i^{r_i},
 \qquad p_i\in\mathcal S_q.
\label{eq:ordering-multilinear}
\end{equation}

\begin{lemma}[单纯形值的局部估计]\label{lem:ordering-rounding}
设 \(J,\mu_c,\alpha\) 为 Lemma~\ref{lem:ordering-local-data} 的
数据，固定 \(0<\lambda<1\)。对整数 \(R\ge1\)，以概率 \(\omega_c\)
选 \(c\)。在这一选定的 \(c\) 的条件下，从 \(\mu_c\) 独立采样
\(R\) 个坐标元组，得串 \(Z_1,\ldots,Z_k\in[m]^R\)。再以概率
\(\lambda\) 独立地把每个位置与坐标均匀重采样，得
\(\widetilde Z_i\)。设单个函数
\(F=(F^1,\ldots,F^q):[m]^R\to\mathcal S_q\) 对每个 \(c\) 与每个
位置通用，且 \(0<t<1/(kq)\)、
\[
 \operatorname{Inf}_\ell(T_{1-\lambda}F^j)\le t
 \quad(\ell\in[R],\ j\in[q]).
\]
则
\begin{equation}
 \mathbb E_{c,Z,\widetilde Z}
 P_0^{\mathrm{ml}}\bigl(F(\widetilde Z_1),\ldots,
                         F(\widetilde Z_k)\bigr)
 \le \operatorname{opt}_q(J)+r_s(kqt),
\label{eq:ordering-rounding}
\end{equation}
其中 \(r_s(u)\to0\)（\(u\downarrow0\)），与 \(R\) 无关。函数
\(r_s\) 只依赖固定的局部数据与 \(k,q,\lambda\)。
\end{lemma}

\begin{proof}
这是 \cite[Lemma~10.5]{GuruswamiEtAl2011OrderingCSP} 的单纯形值
估计，其充分的标量影响界及其依赖被显式化。令 \(H=T_{1-\lambda}F\)。
位置间噪声的独立性与多重线性把
\eqref{eq:ordering-rounding} 的左边变成
\begin{equation}
 \mathbb E_c\mathbb E_{Z\sim\mu_c^{\otimes R}}
 P_0^{\mathrm{ml}}\bigl(H(Z_1),\ldots,H(Z_k)\bigr).
\label{eq:ordering-noise-factorization}
\end{equation}

对每个坐标，积分系综由 \(\mu_c\) 下的示性
\(\mathbf1_{\{X_i=j\}}\) 组成；其原子概率大于 \(\alpha\)。公共
Gram 数据 \eqref{eq:ordering-local-moments} 为 \(J\) 的所有变量
提供同一高斯系综，在每条约束上有相同的一、二阶矩（如
\cite[Lemma~10.4]{GuruswamiEtAl2011OrderingCSP}）。显式地，对 Gram
向量张成中的标准高斯向量 \(g\)，变量
\[
 G_{u,j}=\langle I,b_{u,j}\rangle+
   \langle b_{u,j}-\langle I,b_{u,j}\rangle I,g\rangle
\]
具有那些矩。在 \(R\) 个坐标中使用该系综的独立副本。每个 \(H^j\)
经由其多重线性乘积展开在这些系综上求值；等价地，把 \([m]\) 上的
正交规范基写成符号示性的线性组合，并对积分与高斯系综使用同一个
多项式。

对固定的 \(c\)，设 \(\mathcal T_\theta^c\) 以概率 \(\theta\) 保留
每个完整的坐标元组、否则从 \(\mu_c\) 重采样。均匀的一位边际给出
（逐分量）
\[
 \mathcal T_{1-\lambda}^c[F^j(Z_i)]
 =(T_{1-\lambda}F^j)(Z_i)=H^j(Z_i).
\]
于是这些正是不变性定理在元组系综乘积上所需的平滑分量。同样的
边际给出下面的精确影响恒等式；上标指计算影响所用的乘积测度：
\begin{equation}
 \operatorname{Inf}^{\mu_c^{\otimes R}}_\ell
       \bigl(H^j(Z_i)\bigr)
 =
 \operatorname{Inf}^{[m]^R}_\ell(H^j).
\label{eq:ordering-vector-influence}
\end{equation}
事实上，以除第 \(\ell\) 个之外的每个坐标元组为条件，位置 \(i\)
处的第 \(\ell\) 个符号均匀，而该位置的其他符号具有均匀乘积律。
对向量，影响是其分量影响之和
\cite[Definition~2.6]{IsakssonMossel2009}。因此向量
\((H^j(Z_i))_{i\in[k],j\in[q]}\) 的每个坐标影响至多 \(kqt\)；每
个分量的方差至多 1。

对剩余的尾部假设，设 \((F^j(Z_i))_S\)（\(S\subseteq[R]\)）为元组
乘积 \(\mu_c^{\otimes R}\) 上的 Hoeffding 分量，上标 \(>d\) 保留
\(|S|>d\) 的项。上面的元组噪声恒等式把分量 \(S\) 乘以
\((1-\lambda)^{|S|}\)。由于未平滑的 \(F^j\) 取值于 \([0,1]\)，
对每个实 \(d>0\)，
\[
\begin{aligned}
 \operatorname{Var}\bigl((H^j(Z_i))^{>d}\bigr)
 &=\sum_{|S|>d}(1-\lambda)^{2|S|}
       \bigl\|(F^j(Z_i))_S\bigr\|_2^2\\
 &\le (1-\lambda)^{2d}\operatorname{Var}(F^j(Z_i))
 \le (1-\lambda)^{2d}.
\end{aligned}
\]
此处的范数与方差都用 \(\mu_c^{\otimes R}\)。这给出
\cite[Theorem~3.6]{IsakssonMossel2009} 中的尾部界。它们是
\cite[Theorem~10.1]{GuruswamiEtAl2011OrderingCSP} 中向量值不变性
估计的充分假设。

设 \(\mathsf p:\mathbb R^q\to\mathcal S_q\) 为欧氏投影。它不扩张。
函数
\[
 \Psi(x_1,\ldots,x_k)=
 P_0^{\mathrm{ml}}(\mathsf p(x_1),\ldots,\mathsf p(x_k))
\]
的 Lipschitz 常数至多 \(\sqrt{kq}\)：在单纯形上，对位置裂项把收益
差界定在 \(\sum_i\|x_i-y_i\|_1\)，投影与 Cauchy--Schwarz 给出该
断言。对每个 \(c\) 分别把所引不变性估计用于该函数，向量维数
\(kq\)、影响阈值 \(kqt\)。其显示的误差至多
\[
 r_s(kqt),\qquad
 r_s(u)=C_{kq}\sqrt{kq}\,
             u^{\lambda/(18\log(1/\alpha))}.
\]
特别地，误差与 \(R\) 无关。

在积分系综上，\(H(Z_i)\in\mathcal S_q\)，故 \(\Psi\) 等于
\eqref{eq:ordering-noise-factorization} 中的收益。在公共高斯系综
上，把每个 \(u\in U\) 的 \(H(G_u)\) 投影到 \(\mathcal S_q\)，并
按这些概率独立地为每个 \(u\) 抽一个 \(q\)-标记。这是 \(J\) 的一
个随机赋值。每个局部元组中的变量互异，故其在那个元组处的条件期望
收益是多重线性收益 \(\Psi\)。每个实现的赋值的价值至多
\(\operatorname{opt}_q(J)\)。把不变性比较对 \(c\) 取平均即证
\eqref{eq:ordering-rounding}。
\end{proof}

\begin{proof}[Proposition~\ref{prop:finite-ordering-transfer} 的
证明]
我们先推导对每个 \(R\) 都成立的估计，再独立于 \(R\) 选取其数值
参数。暂固定 \(0<\kappa<1/2\)、满足 \(1/q\le\mu_0(\kappa)\) 的
\(q\)，以及来自 Lemma~\ref{lem:ordering-local-data} 的局部数据。
使用 Lemma~\ref{lem:ordering-rounding} 的局部实验，取实际重采样
概率 \(\lambda=2\kappa\)。再固定辅助常数 \(0<t<1/(kq)\) 与
\(\chi>0\)。

对输入图 \(\Phi\)，把边 \(ab\) 上的置换记为 \(\pi_{a\to b}\)。
其对串的作用由
\[
 (\pi_{a\to b}z)_{\pi_{a\to b}(\ell)}=z_\ell
 \qquad(\ell\in[R])
\]
固定。输出变量集将是 \(\mathcal V=B\times[m]^R\)。考虑如下独立
邻居实验：从 \(A\) 均匀选 \(a\)，从 \(N(a)\) 独立均匀地选
\(b_1,\ldots,b_k\)，并独立地执行局部实验，得
\(\widetilde Z_1,\ldots,\widetilde Z_k\)。其原始输出位置为
\begin{equation}
 w_i=(b_i,\pi_{a\to b_i}\widetilde Z_i)\in\mathcal V.
\label{eq:ordering-iid-queries}
\end{equation}
对任何秩映射 \(f:\mathcal V\to\mathbb Z\)，定义
\[
 V(f)=\mathbb E\,
   \overline{P_0}\bigl(f(w_1),\ldots,f(w_k)\bigr),
 \qquad V_a(f)=\mathbb E[
   \overline{P_0}(f(w_1),\ldots,f(w_k))\mid a].
\]
有界的收益 \eqref{eq:ordering-tie-payoff} 在实验重复某个输出 ID
时也适用。我们将对一个粗秩映射得到下界，对每个单射序得到上界，
随后把两个界转移到互异变量的约束上。

设染色 \(L\) 强满足至少 \(1-\zeta\) 比例的左顶点。定义粗秩映射
\(f(b,z)=z_{L(b)}\)。在每个强满足的 \(a\) 与每个邻居 \(b\) 处，
\[
 f(b,\pi_{a\to b}z)=z_{\pi_{a\to b}^{-1}(L(b))}=z_{L(a)}.
\]
于是对邻居的每种取法，全部 \(k\) 个位置读取同一个标记坐标。以
概率 \((1-2\kappa)^k\) 这 \(k\) 个符号都未被重采样；在该事件下
其元组的律为 \(\mu_c\)。收益的非负性给出
\begin{equation}
 V(f)\ge (1-\zeta)(1-2\kappa)^k v_{\mathrm{loc}}.
\label{eq:ordering-iid-completeness}
\end{equation}

对可靠性，固定单射序 \(O:\mathcal V\to[|\mathcal V|]\)。在固定的
\(a\) 处，其局部集 \(N(a)\times[m]^R\) 恰有 \(M=hm^R\) 个互异
ID。按 \(O\) 排序该集，并把有序表分成 \(q\) 个连续块
\(C_{a,1},\ldots,C_{a,q}\)（各 \(M/q\) 大小）。由 \(q\mid h\)，
这些大小是整数。每块是局部集与一个全局秩区间的交。定义
\begin{equation}
 F_a^j(z)=\frac1h\sum_{b\in N(a)}
  \mathbf1_{\{(b,\pi_{a\to b}z)\in C_{a,j}\}},
 \qquad F_a=(F_a^1,\ldots,F_a^q).
\label{eq:ordering-block-functions}
\end{equation}
对每个 \(z\) 有 \(F_a(z)\in\mathcal S_q\)。每条边置换是串上的
双射，故
\begin{equation}
 \mathbb E_z F_a^j(z)=\frac{|C_{a,j}|}{hm^R}=\frac1q.
\label{eq:ordering-block-mass}
\end{equation}

对整数秩的任何区间 \(I\)，令
\[
 F_a^I(z)=\frac1h\sum_{b\in N(a)}
   \mathbf1_{\{O(b,\pi_{a\to b}z)\in I\}}.
\]
若对每个 \(I\) 与每个 \(\ell\in[R]\) 都有
\(\operatorname{Inf}_\ell(T_{1-\kappa}F_a^I)\le t\)，称 \(a\)
\emph{在阈值 \(t\) 处伪随机}。在这样的 \(a\) 处，该界对每个块
函数 \(F_a^j\) 成立。

设 \(j_a(w)\) 为局部 ID \(w\) 的块指标。在局部查询串的条件下，
邻居选取的独立性给出精确恒等式
\begin{equation}
 \mathbb E_{b_1,\ldots,b_k}
 \overline{P_0}\bigl(j_a(w_1),\ldots,j_a(w_k)\bigr)
 =
 P_0^{\mathrm{ml}}\bigl(F_a(\widetilde Z_1),\ldots,
                         F_a(\widetilde Z_k)\bigr).
\label{eq:ordering-neighbor-factorization}
\end{equation}
该恒等式包括邻居相等的结果。进一步平滑降低块影响，故
\(\operatorname{Inf}_\ell(T_{1-2\kappa}F_a^j)\le t\)。把
Lemma~\ref{lem:ordering-rounding} 用于单个函数 \(F_a\)，把
\eqref{eq:ordering-neighbor-factorization} 中的期望粗收益界定在
\(\operatorname{opt}_q(J)+r_s(kqt)\) 以内。

当 \(k\) 个粗秩互异时，其诱导的置换就是细秩 \(O(w_i)\) 诱导的
那个；否则两个收益相差至多 1。这也涵盖重复的输出 ID——它必然
产生粗平局。对两个位置 \(i\ne i'\)，噪声与邻居选择的独立性给出
其粗平局概率为
\[
 \sum_{j=1}^q
 \mathbb E_c\mathbb E_{Z\sim\mu_c^{\otimes R}}
 \bigl[
  (T_{1-2\kappa}F_a^j)(Z_i)
  (T_{1-2\kappa}F_a^j)(Z_{i'})
 \bigr].
\]
在 \(c\) 的条件下，两个串都有均匀乘积边际。其相关性由
Lemma~\ref{lem:ordering-collision} 允许。把该引理用于
\eqref{eq:ordering-block-mass} 的固定均值 \(1/q\)，记
\(r_c=r_c^{m,\kappa}\)。上一个展示式至多
\(q^{-\kappa/2}+q r_c(t)\)。对顶点对的并界给出：在每个伪随机
\(a\) 处，
\begin{equation}
 V_a(O)\le
 S_{\mathrm{loc}}:=\operatorname{opt}_q(J)
 +\binom{k}{2}\bigl(q^{-\kappa/2}+q r_c(t)\bigr)
 +r_s(kqt).
\label{eq:ordering-pseudorandom-bound}
\end{equation}

接下来把超出该界的部分转换为 \(\Phi\) 的染色。对每个 \(b\in B\)，
把 \(O\) 限制到 \(\{b\}\times[m]^R\) 并压缩其秩，得 \([m]^R\) 的
严格序 \(o_b\)。令
\[
 \mathcal L_b=\left\{\ell\in[R]:
   \operatorname{Inf}_\ell
    \bigl(T_{1-\kappa}\mathbf1_{\{o_b(z)\in I\}}\bigr)>t/2
   \text{ 对某个秩区间 }I\right\}.
\]
少数有影响坐标引理
\cite[Definition~4.2 and Lemma~4.3]{GuruswamiEtAl2011OrderingCSP}
给出
\begin{equation}
 |\mathcal L_b|\le\frac{400}{\kappa(t/2)^3}
 =\frac{3200}{\kappa t^3}.
\label{eq:ordering-right-candidates}
\end{equation}
该界对 \(R\) 一致。

若 \(a\) 非伪随机，选一个见证区间 \(I_a\) 与坐标 \(\ell_a\) 使
\(\operatorname{Inf}_{\ell_a}(T_{1-\kappa}F_a^{I_a})>t\)。噪声与
坐标置换及平均交换，而影响是凸的范数平方。故
\begin{equation}
 \operatorname{Inf}_{\ell_a}(T_{1-\kappa}F_a^{I_a})
 \le\frac1h\sum_{b\in N(a)}
 \operatorname{Inf}_{\pi_{a\to b}(\ell_a)}
 \left(T_{1-\kappa}
       \mathbf1_{\{O(b,z)\in I_a\}}\right).
\label{eq:ordering-influence-average}
\end{equation}
每项至多 1。平均大于 \(t\) 迫使超过 \(t/2\) 比例的项超过
\(t/2\)。全局秩区间限制到单个右码上是其相对序的区间，故这些邻居
满足 \(\pi_{a\to b}(\ell_a)\in\mathcal L_b\)。

给每个这样的 \(a\) 标记 \(\ell_a\)，其他左顶点任意标记。在每个
\(\mathcal L_b\ne\varnothing\) 的右顶点处独立地选 \(\mathcal L_b\)
中的均匀标记；集合为空时用任意标记。由
\eqref{eq:ordering-right-candidates}，每个非伪随机 \(a\) 处各边
的期望被满足比例大于
\[
 \frac t2\,\frac{\kappa t^3}{3200}
 =\frac{\kappa t^4}{6400}.
\]
一个随机的右染色被用于所有左顶点。期望的线性性与公共左度把所得
平均等同于普通的均匀边比例。

若 \(V(O)>S_{\mathrm{loc}}+\chi\)，则超过 \(\chi\) 比例的左顶点
非伪随机：设该比例为 \(g\)，由 \(0\le V_a(O)\le1\) 与
\eqref{eq:ordering-pseudorandom-bound} 得
\(V(O)\le(1-g)S_{\mathrm{loc}}+g\le S_{\mathrm{loc}}+g\)。于是
前面的解码器给出边值大于 \(\chi\kappa t^4/6400\) 的染色。因此
\begin{equation}
 \val(\Phi)\le\nu<\frac{\chi\kappa t^4}{6400}
 \quad\Longrightarrow\quad
 V(O)\le S_{\mathrm{loc}}+\chi\quad\text{对每个单射 }O.
\label{eq:ordering-iid-soundness}
\end{equation}

剩下把独立实验转化为严格排序约束。设 \(D\) 为事件
\(b_1,\ldots,b_k\) 互异。对 \(h\ge k\)，
\begin{equation}
 p_D:=\Pr(D)=\frac{(h)_k}{h^k},\qquad
 \delta_h:=1-p_D\le\frac{\binom{k}{2}}h.
\label{eq:ordering-distinct-probability}
\end{equation}
该概率在每个 \(a\) 处相同，且 \(D\) 与局部实验独立。在 \(D\) 的
条件下，\(a\) 仍均匀，且其有序互异邻居元组在其 \((h)_k\) 种可能
中均匀。此时所有输出 ID 互异（因其首坐标 \(b_i\) 互异）。把这个
有限的带权严格 \(P_0\) 约束分布记为 \(I_0\)，其最大严格收益记为
\(\OPT_{P_0}(I_0)\)。

设 \(V_D(f)\) 为该条件化之后的平局延拓期望收益，
\(V_{\neg D}(f)\) 为以其补为条件者。对任何秩映射 \(f\)（粗或
单射），
\begin{equation}
 V(f)=p_DV_D(f)+(1-p_D)V_{\neg D}(f),
 \qquad |V_D(f)-V(f)|\le\delta_h,
\label{eq:ordering-conditioning}
\end{equation}
因为所有收益都落在 \([0,1]\)。先把它用于
\eqref{eq:ordering-iid-completeness} 的粗映射。在该映射的每个
有限纤维内选互异 ID 的一个均匀随机置换（各纤维独立），并按秩顺序
连接纤维。在每个条件化元组上，这些置换对其 ID 的限制在每个并列
秩类内均匀。于是这一个随机全局序的期望严格 \(P_0\)-收益等于
\(V_D(f)\)。某个全局序至少达到该值。对可靠性：单射 \(O\) 在每个
条件化元组上有互异的秩，故其平局延拓收益即其严格 \(P_0\)-收益。
Equations \eqref{eq:ordering-iid-completeness}、
\eqref{eq:ordering-iid-soundness} 与 \eqref{eq:ordering-conditioning}
因此给出严格界
\begin{equation}
\begin{aligned}
 \mathrm{YES}：\quad&
 \OPT_{P_0}(I_0)\ge(1-\zeta)(1-2\kappa)^k v_{\mathrm{loc}}-\delta_h,\\
 \mathrm{NO}：\quad&
 \OPT_{P_0}(I_0)\le S_{\mathrm{loc}}+\chi+\delta_h
 \quad\text{若 }\nu<\chi\kappa t^4/6400.
\end{aligned}
\label{eq:ordering-conditioned-gap}
\end{equation}
这些界针对条件化后的 \(P_0\) 元组。应用
\eqref{eq:ordering-slot-permutation} 使每个出现以完全相同的值
使用原谓词 \(\Pi\)。称发出的实例为 \(I\)。

现在为指定的 \(\delta\) 作出选取。选正有理 \(\eta\) 使
\(10\eta<\delta\)，选有理 \(0<\kappa<1/2\) 使 \(2\kappa k<\eta\)。
令 \(K=\binom{k}{2}\)。选足够大的整数 \(q\) 使
\[
 \frac1q\le\mu_0(\kappa),\qquad Kq^{-\kappa/2}<\eta.
\]
在误差 \(\eta\) 处取 Lemma~\ref{lem:ordering-local-data} 的数据
\(J,m,\alpha\)。两个解析误差函数现在有固定参数且对 \(R\)
一致。选足够小的正有理 \(t<1/(kq)\) 使
\[
 Kq r_c(t)<\eta,\qquad r_s(kqt)<\eta.
\]
置 \(\chi=\eta\)。选正有理 \(\zeta,\nu<1\) 使
\[
 \zeta<\eta,\qquad
 \nu<\frac{\eta\kappa t^4}{6400},
\]
并选整数 \(H\ge\max\{k,q\}\) 使 \(K/H<\eta\)。
这些选取只依赖 \(k,\Pi,\delta\)。特别地，它们先于
Lemma~\ref{lem:simple-strong-left} 中 \(h\) 的选取，也先于
Theorem~\ref{thm:ugc} 提供的字母表。

对每个 \(h\ge H\) 有 \(\delta_h<\eta\)。局部间隙
\eqref{eq:ordering-local-gap}、Bernoulli 不等式与
\eqref{eq:ordering-conditioned-gap} 给出
\[
 \mathrm{YES}：\quad
 \OPT_\Pi(I)\ge1-\zeta-2\kappa k-\eta-\delta_h
 >1-4\eta>1-\delta,
\]
而 \eqref{eq:ordering-pseudorandom-bound} 给出
\[
 \mathrm{NO}：\quad
 \OPT_\Pi(I)\le S_{\mathrm{loc}}+\chi+\delta_h
 <\rho_\Pi+6\eta<\rho_\Pi+\delta.
\]
这同时对每个 \(R\) 证明了数值断言。

最后指明编码后的变换。一旦 \(R\) 固定，\(c\) 的选取、乘积
\(\mu_c^{\otimes R}\) 与独立的有理 \(2\kappa\)-噪声就构成局部
结果 \(\xi\) 上的一个固定的有限有理律 \(\theta_R\)。这些就是
上述的实际局部结果；解析比较所用的高斯变量不属于输出。把 \(B\)
重新编号为 \([|B|]\)，显式列出 \(\mathcal V=B\times[m]^R\)。
对每个 \(a\in A\)、每个有序互异邻居元组、以及每个
\(\xi\in\operatorname{supp}\theta_R\)，列出
\eqref{eq:ordering-iid-queries} 的元组（经
\eqref{eq:ordering-slot-permutation} 置换），权重为
\begin{equation}
 \frac{\theta_R(\xi)}{|A|(h)_k}.
\label{eq:ordering-rational-weight}
\end{equation}
权重和为 1 且非负有理。每个列出的元组有互异的 ID；同一元组可以
出现在多个条目中。变量有 \(|B|m^R\) 个，条目至多
\(|A|(h)_k|\operatorname{supp}\theta_R|\) 个。对固定的
\(k,m,R,\theta_R\)，即使 \(h\) 变化，这些计数也在显式图规模上
多项式。\(\theta_R\) 的固定公共分母表明每个权重有
\(O(\log|A|+k\log h)\) 比特加一个固定常数。ID 有
\(O(\log|B|)\) 比特加一个固定常数。因此这些条目的精确枚举对
固定 \(R\) 是确定性的多项式，如所断言。
\end{proof}

同一出现列表编码也为 Section~\ref{sec:consequences} 中的确定性
随机序保证提供了依据。逐次暴露下一个变量来构建序。若某约束恰有
\(s\) 个变量已被暴露，其条件成功概率是 \(\Pi\) 中「前 \(s\) 个
位置恰为这些变量在已暴露序中的角色」的置换个数除以 \((k-s)!\)。
可以通过检查固定的集合 \(S_k\) 计算它。对下一个变量的每种取法，
把这些概率对列出的有理权重求和，并选条件期望不低于各取法平均值的
下一个变量。这保持初始期望 \(|\Pi|/k!\)。所有比较都是精确的，且
在输入位长上多项式：权重分母之积的位长至多其总编码长度，而阶乘
分母对固定的 \(k\) 是常数规模。变量列表是显式的，故前缀步数与
候选取法的数目也是多项式的。

\bibliographystyle{plain}
\bibliography{references}
\end{document}
